\documentclass[UTF8, a4paper, 12pt, fontset=fandol, scheme=chinese]{ctexart}
% ============================================================
%  SY02-22《Springer 理论与轨道积分讲义》译稿导言区
%  底本：Zhiwei Yun, Lectures on Springer theories and orbital integrals,
%        arXiv:1602.01451（PCMI 2015 暑期学校讲义，62 页）
%  体例：复现原文 pcmi.cls 的编号机制——\numberwithin{equation}{subsection}，
%        定理类环境与公式共用一个计数器（印成「定理 1.5.1」「(1.5.2)」），
%        并且 \subsubsection 也消耗这个计数器（印成「1.5.3. 约定」）。
% ============================================================

% ---------- 页面与字体 ----------
\usepackage{geometry}
\geometry{a4paper, left=2.6cm, right=2.6cm, top=2.8cm, bottom=2.6cm}
\linespread{1.22}
\setlength{\parskip}{0.25em}
\setlength{\emergencystretch}{1.5em}
\sloppy

% ---------- 数学 ----------
\usepackage{amsmath}
\usepackage{amsthm}
\usepackage{amssymb}
\usepackage{amsfonts}
\usepackage{mathtools}
\usepackage{mathrsfs}
\usepackage{stmaryrd}
\usepackage{bm}
\allowdisplaybreaks[2]

% ---------- 图表 ----------
\usepackage{graphicx}
\usepackage{xcolor}
\usepackage{tikz}
\usetikzlibrary{cd, arrows.meta, positioning, calc}
\usepackage[most]{tcolorbox}
\usepackage{enumitem}
\usepackage{xspace}   % 原文宏 \sB 等用到 \xspace
\setlist{itemsep=2pt, topsep=3pt, parsep=0pt}
\setlist[enumerate,1]{label=(\arabic*)}

% ---------- 交换图：\xymatrix 用的 xy 必须在 tikz 之后 ----------
\usepackage[all]{xy}

% ---------- 超链接 ----------
\usepackage[hyperfootnotes=false, colorlinks=true, linkcolor=blue!45!black, citecolor=green!35!black,
            urlcolor=teal!60!black, bookmarksnumbered=true]{hyperref}
\usepackage{bookmark}

% ---------- 中文微调 ----------
\xeCJKsetup{CJKecglue={\hskip 0.24em plus 0.05em minus 0.05em}}
\ctexset{
  section    = {name={,}, aftername={.\enspace}, format=\Large\bfseries\centering},
  subsection = {aftername={.\enspace}, format=\large\bfseries},
  subsubsection = {aftername={.\enspace}, format=\normalsize\bfseries}
}

% ---------- 定理环境：全部挂在 equation 计数器上（原文即如此） ----------
% 原文用 \swapnumbers 把编号印在「Theorem」之前；中译本按汉语习惯印成
% 「定理 1.5.1（Springer）」，编号值不变。定理头不用句点收尾。
\makeatletter
\newtheoremstyle{cnplain}{3pt plus 1pt minus 1pt}{3pt plus 1pt}{}{}{\bfseries}{\relax}{0.6em}{%
  \thmname{#1}\thmnumber{\@ifnotempty{#1}{\ }#2}\@ifempty{#3}{}{（#3）}}
\newtheoremstyle{cndefn}{3pt plus 1pt minus 1pt}{3pt plus 1pt}{}{}{\bfseries}{\relax}{0.6em}{%
  \thmname{#1}\thmnumber{\@ifnotempty{#1}{\ }#2}\@ifempty{#3}{}{（#3）}}
\newtheoremstyle{cnrmk}{3pt plus 1pt minus 1pt}{3pt plus 1pt}{}{}{\bfseries}{\relax}{0.6em}{%
  \thmname{#1}\thmnumber{\@ifnotempty{#1}{\ }#2}\@ifempty{#3}{}{（#3）}}
\makeatother

\numberwithin{equation}{subsection}
\theoremstyle{cnplain}
\newtheorem{theorem}[equation]{定理}
\newtheorem{lemma}[equation]{引理}
\newtheorem{prop}[equation]{命题}
\newtheorem{cor}[equation]{推论}
\newtheorem{conj}[equation]{猜想}
\newtheorem{fact}[equation]{事实}
\newtheorem{ansatz}[equation]{拟设}
\theoremstyle{cndefn}
\newtheorem{defn}[equation]{定义}
\newtheorem{cons}[equation]{构造}
\theoremstyle{cnrmk}
\newtheorem{remark}[equation]{注}
\newtheorem{exam}[equation]{例}
\newtheorem{ques}[equation]{问题}

\renewcommand{\proofname}{证明}
\renewcommand{\qedsymbol}{$\square$}
\renewcommand{\refname}{参考文献}
\renewcommand{\contentsname}{目\quad 录}

% ---------- \subsubsection 消耗 equation 计数器（原文导言区的核心技巧） ----------
% 原文：\let\oldsubsubsection\subsubsection 后重定义，先令 subsubsection 计数器
% 取 equation 的当前值，再调原命令（自增 1，于是编号等于新的 equation 值），
% 最后再 \refstepcounter{equation} 跳过这个号。空标题 \subsubsection{} 同样占号。
\makeatletter
\renewcommand{\thesubsubsection}{\thesubsection.\arabic{subsubsection}}
\let\oldsubsubsection\subsubsection
\renewcommand{\subsubsection}[1]{%
  \setcounter{subsubsection}{\value{equation}}%
  \oldsubsubsection{#1}%
  \refstepcounter{equation}%
}
\makeatother

% ---------- 译注标记（独立脚注序列） ----------
\usepackage{amscd}
\usepackage{manyfoot}
\DeclareNewFootnote[plain]{trans}
\renewcommand{\thefootnotetrans}{译\arabic{footnotetrans}}
\newcommand{\zy}[1]{\footnotetrans{#1}}
\newtcolorbox{yizhu}[1][]{
  enhanced, breakable, colback=orange!4!white, colframe=orange!55!black,
  boxrule=0.5pt, left=6pt, right=6pt, top=4pt, bottom=4pt,
  fonttitle=\bfseries\small, title=#1,
  before skip=6pt, after skip=6pt
}

\providecommand{\cprime}{'}
\providecommand{\vprime}{''}
\providecommand{\sprime}{'}

% ---------- 原文宏 ----------
% ============================================================
%  SY02-22《Springer 理论与轨道积分讲义》译稿：原文自定义宏
%  由 工程文件/scripts/port_orig_macros.py 从原文导言区（340--797 行）抽取，
%  只保留数学记号宏；原文的 \newtheorem、\usepackage、\swapnumbers、
%  subsubsection 重定义等已由 preamble.tex 统一处理。
% ============================================================

%%%%%%%% mathbb %%%%%%%%

\def\AA{\mathbb{A}}
\def\BB{\mathbb{B}}
\def\CC{\mathbb{C}}
\def\DD{\mathbb{D}}
\def\EE{\mathbb{E}}
\def\FF{\mathbb{F}}
\def\GG{\mathbb{G}}
\def\HH{\mathbb{H}}
\def\II{\mathbb{I}}
\def\JJ{\mathbb{J}}
\def\KK{\mathbb{K}}
\def\LL{\mathbb{L}}
\def\MM{\mathbb{M}}
\def\NN{\mathbb{N}}
\def\OO{\mathbb{O}}
\def\PP{\mathbb{P}}
\def\QQ{\mathbb{Q}}
\def\RR{\mathbb{R}}
\def\SS{\mathbb{S}}
\def\TT{\mathbb{T}}
\def\UU{\mathbb{U}}
\def\VV{\mathbb{V}}
\def\WW{\mathbb{W}}
\def\XX{\mathbb{X}}
\def\YY{\mathbb{Y}}
\def\ZZ{\mathbb{Z}}

%%%%%%%% mathcal %%%%%%%%

\def\calA{\mathcal{A}}
\def\calB{\mathcal{B}}
\def\calC{\mathcal{C}}
\def\calD{\mathcal{D}}
\def\calE{\mathcal{E}}
\def\calF{\mathcal{F}}
\def\calG{\mathcal{G}}
\def\calH{\mathcal{H}}
\def\calI{\mathcal{I}}
\def\calJ{\mathcal{J}}
\def\calK{\mathcal{K}}
\def\calL{\mathcal{L}}
\def\calM{\mathcal{M}}
\def\calN{\mathcal{N}}
\def\calO{\mathcal{O}}
\def\calP{\mathcal{P}}
\def\calQ{\mathcal{Q}}
\def\calR{\mathcal{R}}
\def\calS{\mathcal{S}}
\def\calT{\mathcal{T}}
\def\calU{\mathcal{U}}
\def\calV{\mathcal{V}}
\def\calW{\mathcal{W}}
\def\calX{\mathcal{X}}
\def\calY{\mathcal{Y}}
\def\calZ{\mathcal{Z}}


%%%%%%%% mathscr %%%%%%%%

%IM Changed \mathscr to \mathcal since the former looks 
%       completely wrong with the other fonts in the PCMI style.

\long\def\commentout#1{}
\commentout{
\newcommand{\sA}{\ensuremath{\mathscr{A}}\xspace}
\newcommand{\sB}{\ensuremath{\mathscr{B}}\xspace}
\newcommand{\sC}{\ensuremath{\mathscr{C}}\xspace}
\newcommand{\sD}{\ensuremath{\mathscr{D}}\xspace}
\newcommand{\sE}{\ensuremath{\mathscr{E}}\xspace}
\newcommand{\sF}{\ensuremath{\mathscr{F}}\xspace}
\newcommand{\sG}{\ensuremath{\mathscr{G}}\xspace}
\newcommand{\sH}{\ensuremath{\mathscr{H}}\xspace}
\newcommand{\sI}{\ensuremath{\mathscr{I}}\xspace}
\newcommand{\sJ}{\ensuremath{\mathscr{J}}\xspace}
\newcommand{\sK}{\ensuremath{\mathscr{K}}\xspace}
\newcommand{\sL}{\ensuremath{\mathscr{L}}\xspace}
\newcommand{\sM}{\ensuremath{\mathscr{M}}\xspace}
\newcommand{\sN}{\ensuremath{\mathscr{N}}\xspace}
\newcommand{\sO}{\ensuremath{\mathscr{O}}\xspace}
\newcommand{\sP}{\ensuremath{\mathscr{P}}\xspace}
\newcommand{\sQ}{\ensuremath{\mathscr{Q}}\xspace}
\newcommand{\sR}{\ensuremath{\mathscr{R}}\xspace}
\newcommand{\sS}{\ensuremath{\mathscr{S}}\xspace}
\newcommand{\sT}{\ensuremath{\mathscr{T}}\xspace}
\newcommand{\sU}{\ensuremath{\mathscr{U}}\xspace}
\newcommand{\sV}{\ensuremath{\mathscr{V}}\xspace}
\newcommand{\sW}{\ensuremath{\mathscr{W}}\xspace}
\newcommand{\sX}{\ensuremath{\mathscr{X}}\xspace}
\newcommand{\sY}{\ensuremath{\mathscr{Y}}\xspace}
\newcommand{\sZ}{\ensuremath{\mathscr{Z}}\xspace}
}%commentout

\newcommand{\sA}{\ensuremath{\mathcal{A}}\xspace}
\newcommand{\sB}{\ensuremath{\mathcal{B}}\xspace}
\newcommand{\sC}{\ensuremath{\mathcal{C}}\xspace}
\newcommand{\sD}{\ensuremath{\mathcal{D}}\xspace}
\newcommand{\sE}{\ensuremath{\mathcal{E}}\xspace}
\newcommand{\sF}{\ensuremath{\mathcal{F}}\xspace}
\newcommand{\sG}{\ensuremath{\mathcal{G}}\xspace}
\newcommand{\sH}{\ensuremath{\mathcal{H}}\xspace}
\newcommand{\sI}{\ensuremath{\mathcal{I}}\xspace}
\newcommand{\sJ}{\ensuremath{\mathcal{J}}\xspace}
\newcommand{\sK}{\ensuremath{\mathcal{K}}\xspace}
\newcommand{\sL}{\ensuremath{\mathcal{L}}\xspace}
\newcommand{\sM}{\ensuremath{\mathcal{M}}\xspace}
\newcommand{\sN}{\ensuremath{\mathcal{N}}\xspace}
\newcommand{\sO}{\ensuremath{\mathcal{O}}\xspace}
\newcommand{\sP}{\ensuremath{\mathcal{P}}\xspace}
\newcommand{\sQ}{\ensuremath{\mathcal{Q}}\xspace}
\newcommand{\sR}{\ensuremath{\mathcal{R}}\xspace}
\newcommand{\sS}{\ensuremath{\mathcal{S}}\xspace}
\newcommand{\sT}{\ensuremath{\mathcal{T}}\xspace}
\newcommand{\sU}{\ensuremath{\mathcal{U}}\xspace}
\newcommand{\sV}{\ensuremath{\mathcal{V}}\xspace}
\newcommand{\sW}{\ensuremath{\mathcal{W}}\xspace}
\newcommand{\sX}{\ensuremath{\mathcal{X}}\xspace}
\newcommand{\sY}{\ensuremath{\mathcal{Y}}\xspace}
\newcommand{\sZ}{\ensuremath{\mathcal{Z}}\xspace}


%%%%%%%%% bold face %%%%%%%%%%

\def\bA{\mathbf{A}}
\def\bB{\mathbf{B}}
\def\bC{\mathbf{C}}
\def\bD{\mathbf{D}}
\def\bE{\mathbf{E}}
\def\bF{\mathbf{F}}
\def\bG{\mathbf{G}}
\def\bH{\mathbf{H}}
\def\bI{\mathbf{I}}
\def\bJ{\mathbf{J}}
\def\bK{\mathbf{K}}
\def\bL{\mathbf{L}}
\def\bM{\mathbf{M}}
\def\bN{\mathbf{N}}
\def\bO{\mathbf{O}}
\def\bP{\mathbf{P}}
\def\bQ{\mathbf{Q}}
\def\bR{\mathbf{R}}
\def\bS{\mathbf{S}}
\def\bT{\mathbf{T}}
\def\bU{\mathbf{U}}
\def\bV{\mathbf{V}}
\def\bW{\mathbf{W}}
\def\bX{\mathbf{X}}
\def\bY{\mathbf{Y}}
\def\bZ{\mathbf{Z}}

%%%%%%%% frak %%%%%%%%


\newcommand\frA{\mathfrak{A}}
\newcommand\frB{\mathfrak{B}}
\newcommand\frC{\mathfrak{C}}
\newcommand\frD{\mathfrak{D}}
\newcommand\frE{\mathfrak{E}}
\newcommand\frF{\mathfrak{F}}
\newcommand\frG{\mathfrak{G}}
\newcommand\frH{\mathfrak{H}}
\newcommand\frI{\mathfrak{I}}
\newcommand\frJ{\mathfrak{J}}
\newcommand\frK{\mathfrak{K}}
\newcommand\frL{\mathfrak{L}}
\newcommand\frM{\mathfrak{M}}
\newcommand\frN{\mathfrak{N}}
\newcommand\frO{\mathfrak{O}}
\newcommand\frP{\mathfrak{P}}
\newcommand\frQ{\mathfrak{Q}}
\newcommand\frR{\mathfrak{R}}
\newcommand\frS{\mathfrak{S}}
\newcommand\frT{\mathfrak{T}}
\newcommand\frU{\mathfrak{U}}
\newcommand\frV{\mathfrak{V}}
\newcommand\frW{\mathfrak{W}}
\newcommand\frX{\mathfrak{X}}
\newcommand\frY{\mathfrak{Y}}
\newcommand\frZ{\mathfrak{Z}}

\newcommand\fra{\mathfrak{a}}
\newcommand\frb{\mathfrak{b}}
\newcommand\frc{\mathfrak{c}}
\newcommand\frd{\mathfrak{d}}
\newcommand\fre{\mathfrak{e}}
\newcommand\frf{\mathfrak{f}}
\newcommand\frg{\mathfrak{g}}
\newcommand\frh{\mathfrak{h}}
\newcommand\fri{\mathfrak{i}}
\newcommand\frj{\mathfrak{j}}
\newcommand\frk{\mathfrak{k}}
\newcommand\frl{\mathfrak{l}}
\newcommand\frakm{\mathfrak{m}}
\newcommand\frn{\mathfrak{n}}
\newcommand\frp{\mathfrak{p}}
\newcommand\frq{\mathfrak{q}}
\newcommand\frr{\mathfrak{r}}
\newcommand\frs{\mathfrak{s}}
\newcommand\frt{\mathfrak{t}}
\newcommand\fru{\mathfrak{u}}
\newcommand\frv{\mathfrak{v}}
\newcommand\frw{\mathfrak{w}}
\newcommand\frx{\mathfrak{x}}
\newcommand\fry{\mathfrak{y}}
\newcommand\frz{\mathfrak{z}}






%%%%%%%% tilde %%%%%%%%%

\newcommand\tilA{\widetilde{A}}
\newcommand\tilB{\widetilde{B}}
\newcommand\tilC{\widetilde{C}}
\newcommand\tilD{\widetilde{D}}
\newcommand\tilE{\widetilde{E}}
\newcommand\tilF{\widetilde{F}}
\newcommand\tilG{\widetilde{G}}
\newcommand\tilH{\widetilde{H}}
\newcommand\tilI{\widetilde{I}}
\newcommand\tilJ{\widetilde{J}}
\newcommand\tilK{\widetilde{K}}
\newcommand\tilL{\widetilde{L}}
\newcommand\tilM{\widetilde{M}}
\newcommand\tilN{\widetilde{N}}
\newcommand\tilO{\widetilde{O}}
\newcommand\tilP{\widetilde{P}}
\newcommand\tilQ{\widetilde{Q}}
\newcommand\tilR{\widetilde{R}}
\newcommand\tilS{\widetilde{S}}
\newcommand\tilT{\widetilde{T}}
\newcommand\tilU{\widetilde{U}}
\newcommand\tilV{\widetilde{V}}
\newcommand\tilW{\widetilde{W}}
\newcommand\tilX{\widetilde{X}}
\newcommand\tilY{\widetilde{Y}}
\newcommand\tilZ{\widetilde{Z}}

\newcommand\tila{\widetilde{a}}
\newcommand\tilb{\widetilde{b}}
\newcommand\tilc{\widetilde{c}}
\newcommand\tild{\widetilde{d}}
\newcommand\tile{\widetilde{e}}
\newcommand\tilf{\widetilde{f}}
\newcommand\tilg{\widetilde{g}}
\newcommand\tilh{\widetilde{h}}
\newcommand\tili{\widetilde{i}}
\newcommand\tilj{\widetilde{j}}
\newcommand\tilk{\widetilde{k}}
%\newcommand\till{\widetilde{l}}
\newcommand\tilm{\widetilde{m}}
\newcommand\tiln{\widetilde{n}}
\newcommand\tilo{\widetilde{o}}
\newcommand\tilp{\widetilde{p}}
\newcommand\tilq{\widetilde{q}}
\newcommand\tilr{\widetilde{r}}
\newcommand\tils{\widetilde{s}}
\newcommand\tilt{\widetilde{t}}
\newcommand\tilu{\widetilde{u}}
\newcommand\tilv{\widetilde{v}}
\newcommand\tilw{\widetilde{w}}
\newcommand\tilx{\widetilde{x}}
\newcommand\tily{\widetilde{y}}
\newcommand\tilz{\widetilde{z}}

%%%%%%%% hat %%%%%%%%%

\def\hatG{\widehat{G}}


%%%%%%%% check %%%%%%%%


\def\dG{G^{\vee}}
\def\dB{B^{\vee}}
\def\dT{T^{\vee}}

%%%%%%%% standard operators %%%%%%%%


\newcommand\aff{\textup{aff}}
\newcommand\Aff{\textup{Aff}}
\newcommand\alg{\textup{alg}}
\newcommand\Alg{\textup{Alg}}
\newcommand\AS{\textup{AS}}
\newcommand\Av{\textup{Av}}
\newcommand\av{\textup{av}}
\newcommand\BM{\textup{BM}}
\newcommand{\Bun}{\textup{Bun}}
\newcommand{\can}{\textup{can}}
%\renewcommand{\char}{\textup{char}}  dangerous, conflicts with xy package
\newcommand{\codim}{\textup{codim}}
\newcommand{\Coh}{\textup{Coh}}
\newcommand{\coker}{\textup{coker}}
\newcommand{\Coor}{\textup{Coor}}
\newcommand{\Corr}{\textup{Corr}}
\newcommand{\diag}{\textup{diag}}
\DeclareMathOperator\Dyn{Dyn}
\newcommand\ev{\textup{ev}}
\newcommand{\even}{\textup{even}}
\newcommand{\Fil}{\textup{Fil}}
\newcommand{\Fl}{\textup{Fl}}
\newcommand\Forg{\textup{Forg}}
\newcommand\Fr{\textup{Fr}}
\newcommand\Frac{\textup{Frac}}
\newcommand\Frob{\textup{Frob}}
\newcommand\Gal{\textup{Gal}}
\newcommand\geom{\textup{geom}}
\newcommand\glob{\textup{glob}}
\newcommand{\Gr}{\textup{Gr}}
\newcommand{\gr}{\textup{gr}}
\newcommand{\Hecke}{\textup{Hecke}}
\newcommand{\Hilb}{\textup{Hilb}}
\newcommand\IC{\textup{IC}}
\newcommand\id{\textup{id}}
\renewcommand{\Im}{\textup{Im}}
\newcommand{\Ind}{\textup{Ind}}
\newcommand{\ind}{\textup{ind}}
\newcommand\inv{\textup{inv}}
\newcommand\Irr{\textup{Irr}}
\newcommand\Jac{\textup{Jac}}
\newcommand\Lie{\textup{Lie}\ }
\newcommand\Loc{\textup{Loc}}
\newcommand\loc{\textup{loc}}
\newcommand\Mod{\textup{Mod}}
\newcommand{\mult}{\textup{mult}}
\newcommand\nil{\textup{nil}}
\newcommand{\Nm}{\textup{Nm}}
\newcommand{\odd}{\textup{odd}}
\newcommand{\op}{\textup{op}}
\newcommand\Out{\textup{Out}}
\newcommand{\Perf}{\textup{Perf}}
\newcommand\Perv{\textup{Perv}}
\newcommand{\Pic}{\textup{Pic}}
\newcommand{\cPic}{\overline{\textup{Pic}}}
\newcommand\pr{\textup{pr}}
\newcommand\pre{\textup{pre}}
\newcommand\pro{\textup{pro}}
\newcommand{\Quot}{\textup{Quot}}
\newcommand\rank{\textup{rank}\ }
\newcommand{\red}{\textup{red}}
\newcommand{\reg}{\textup{reg}}
\newcommand\Rep{\textup{Rep}}
\newcommand{\Res}{\textup{Res}}
\newcommand\res{\textup{res}}
\newcommand\rk{\textup{rk}}
\newcommand\rs{\textup{rs}}
\newcommand\Sht{\textup{Sht}}
\newcommand\Span{\textup{Span}}
\newcommand\Spec{\textup{Spec}\ }
\newcommand\St{\textup{St}}
\newcommand\st{\textup{st}}
\newcommand\Stab{\textup{Stab}}
\newcommand\Supp{\textup{Supp}}
\newcommand\Sym{\textup{Sym}}
\newcommand{\tors}{\textup{tors}}
\newcommand\Tot{\textup{Tot}}
\newcommand{\Tr}{\textup{Tr}}
\newcommand{\tr}{\textup{tr}}
\newcommand\triv{\textup{triv}}
\newcommand\unip{\textup{unip}}
\newcommand{\univ}{\textup{univ}}
\newcommand{\val}{\textup{val}}
\newcommand{\Vect}{\textup{Vect}}
\newcommand{\vol}{\textup{vol}}


\newcommand\Aut{\textup{Aut}}
\newcommand\Hom{\textup{Hom}}
\newcommand\End{\textup{End}}
\newcommand\Isom{\textup{Isom}}
\newcommand\Map{\textup{Map}}
\newcommand\Mor{\textup{Mor}}
\newcommand\uAut{\underline{\Aut}}
\newcommand\uHom{\underline{\Hom}}
\newcommand\uEnd{\underline{\End}}
\newcommand\uMor{\underline{\Mor}}
\newcommand{\RHom}{\bR\Hom}
\newcommand\RuHom{\bR\uHom}
\newcommand{\Ext}{\textup{Ext}}
\newcommand{\uExt}{\underline{\Ext}}


%%%%%%%% Lie groups and algebras %%%%%%%
\newcommand\GL{\textup{GL}}
\newcommand\gl{\mathfrak{gl}}
\newcommand\PGL{\textup{PGL}}
\newcommand\Ug{\textup{U}}
\newcommand\ug{\mathfrak{u}}
\newcommand\SL{\textup{SL}}
\renewcommand\sl{\mathfrak{sl}}
\newcommand\SU{\textup{SU}}
\newcommand\su{\mathfrak{su}}
\newcommand\GU{\textup{GU}}
\newcommand\SO{\textup{SO}}
\newcommand\GO{\textup{GO}}
\newcommand\Og{\textup{O}}
\newcommand\so{\mathfrak{so}}
\newcommand\Sp{\textup{Sp}}
\renewcommand\sp{\mathfrak{sp}}
\newcommand\GSp{\textup{GSp}}

\newcommand{\Gm}{\GG_m}
\def\Ga{\GG_a}

\newcommand{\ad}{\textup{ad}}
\newcommand{\Ad}{\textup{Ad}}
\renewcommand\sc{\textup{sc}}
\newcommand{\der}{\textup{der}}

\newcommand\xch{\mathbb{X}^*}
\newcommand\xcoch{\mathbb{X}_*}


%%%%%%% common symbols %%%%%%%%

\newcommand{\incl}{\hookrightarrow}
\newcommand{\isom}{\stackrel{\sim}{\to}}
\newcommand{\bij}{\leftrightarrow}
\newcommand{\surj}{\twoheadrightarrow}
\newcommand{\ep}{\epsilon}
\renewcommand{\l}{\lambda}
\renewcommand{\L}{\Lambda}
\newcommand{\om}{\omega}
\newcommand{\Om}{\Omega}
\newcommand{\one}{\mathbf{1}}


\newcommand{\Ql}{\QQ_{\ell}}
\newcommand{\Qlbar}{\overline{\QQ}_\ell}
\newcommand{\const}[1]{\overline{\QQ}_{\ell,#1}}
\newcommand{\twtimes}[1]{\stackrel{#1}{\times}}
\newcommand{\jiao}[1]{\langle{#1}\rangle}
\newcommand{\wt}{\widetilde}
\newcommand{\wh}{\widehat}
\newcommand{\bs}{\backslash}
\newcommand\mat[4]{\left(\begin{array}{cc} #1 & #2 \\ #3 & #4 \end{array}\right)}  % 2-by-2 matrix
\newcommand\un{\underline}
\newcommand\ov{\overline}
\newcommand\bu{\bullet}
\newcommand\chark{\textup{char}(k)}
\newcommand{\tl}[1]{[\![#1]\!]}
\newcommand{\lr}[1]{(\!(#1)\!)}


\newcommand\fdeg{\textup{fdeg}}
\newcommand\Wa{W_{\textup{aff}}}

\newcommand{\homo}[2]{\mathbf{H}_{#1}({#2})}   % sheaf
\newcommand{\homog}[2]{\textup{H}_{#1}({#2})}  % plain group
\newcommand{\coho}[2]{\mathbf{H}^{#1}({#2})}    % sheaf
\newcommand{\cohog}[2]{\textup{H}^{#1}({#2})}     % plain group
\newcommand{\cohoc}[2]{\textup{H}_{c}^{#1}({#2})}     % compact support
\newcommand{\hBM}[2]{\textup{H}^{\textup{BM}}_{#1}({#2})}  % Borel-Moore

\newcommand\upH{\textup{H}}

\newcommand{\oll}[1]{\overleftarrow{#1}}
\newcommand{\orr}[1]{\overrightarrow{#1}}
\newcommand{\olr}[1]{\overleftrightarrow{#1}}

%%%%% bar %%%%%%%%
\newcommand{\kbar}{\overline{k}}
\newcommand{\Qbar}{\overline{\QQ}}


%%%%%% Galois groups %%%%%%%%
\newcommand{\Gk}{\Gal(\kbar/k)}
\newcommand{\GQ}{\Gal(\overline{\QQ}/\QQ)}
\newcommand{\GQp}{\Gal(\overline{\QQ}_{p}/\QQ_{p})}

%%%%%%% Spaces %%%%%%%%%%

\newcommand{\tN}{\wt{\sN}}
\newcommand{\tX}{\wt{\sX}}
\newcommand{\tY}{\wt{\sY}}
\newcommand{\tB}{\wt{\sB}}
\newcommand{\tg}{\wt{\frg}}
\newcommand{\Ah}{\calA^{\heartsuit}}
\DeclareMathOperator\subreg{subreg}

\newcommand{\Prym}{\textup{Prym}}
\newcommand{\cPrym}{\overline{\textup{Prym}}}
\newcommand{\ur}{\textup{ur}}
\newcommand{\Mpar}{\calM^{\textup{par}}}
\newcommand{\fpar}{f^{\textup{par}}}


% ---------- 原文类文件提供、但正文另需的宏 ----------
\providecommand{\scr}[1]{\mathscr{#1}}
\providecommand{\bR}{\mathbf{R}}
\providecommand{\bL}{\mathbf{L}}

% 参考文献用：原文 amsrefs 条目里的 \MR 与 \bibarxiv 只保留正文，不排版
\providecommand{\MR}[1]{}
\providecommand{\bibarxiv}[1]{arXiv:#1}
\providecommand{\biburl}[1]{\url{#1}}


\begin{document}
\pagestyle{plain}

% ============ 扉页 ============
\begin{titlepage}
\centering
\vspace*{2.2cm}

{\Huge\bfseries Springer 理论与\\[6pt] 轨道积分讲义}

\vspace{1.0cm}
{\LARGE\itshape Lectures on Springer theories\\[4pt] and orbital integrals}

\vspace{0.9cm}
{\large 中文译注本}

\vspace{1.4cm}
{\large 恽之玮（Zhiwei Yun）}

\vspace{0.4cm}
{\large PCMI 2015 暑期学校讲义}

\vspace{1.6cm}
{\normalsize 原文：arXiv:1602.01451，62 页，2016 年}

\vspace{0.3cm}
{\normalsize 译者：数译（DSH agent）\quad 版本：v1.0\quad 日期：2026 年 11 月}

\vspace{1.2cm}
\rule{0.72\textwidth}{0.4pt}

\vfill
{\small 本译本仅供学术研读；原文版权归作者所有。}
\end{titlepage}

% ============ 摘要 ============
\renewcommand{\abstractname}{摘\quad 要}
\begin{abstract}
本讲义是作者 2015 年帕克城数学研究所（Park City Mathematics Institute，PCMI）研究生暑期学校短期课程的扩充讲义。
主要内容为：Springer 纤维、仿射 Springer 纤维与 Hitchin 纤维的几何；这些对象产生的（仿射）Weyl 群表示；
仿射 Springer 纤维与轨道积分之间的联系。
\end{abstract}

% ============ 译者说明 ============
\section*{译者说明}
\addcontentsline{toc}{section}{译者说明}

\textbf{原文性质。}这是恽之玮（Zhiwei Yun）在 2015 年 PCMI 暑期学校（Park City，专题「几何朗兰兹纲领与自守形式」）
的系列讲义，共四讲加一篇引言，收入 IAS/Park City Mathematics Series。四讲分别处理 Springer 纤维、仿射 Springer 纤维、
轨道积分与 Hitchin 纤维化，主线是三类「纤维」上的 Weyl 群作用，以及它们同轨道积分、基本引理的联系。

\textbf{编号约定。}原文用 PCMI 的类文件，全部定理类环境、公式、图、表共用一个计数器，按「节.小节.序号」三段编号
（例如「定理 1.5.1」「(1.5.2)」）；原文还让编号的小标题（subsubsection）消耗同一个计数器，因此「1.5.3. 约定」
这样的标题本身也占一个号。译文逐处复现这套机制，编号与原文一一对应。原文的引言用 0 编节，故有「0.1」「0.2」。

\textbf{译注标记。}文中「译\emph{n}」是译者所加，用独立脚注序列，不占用作者的脚注编号；作者脚注保留原编号。
译注只写三类内容：易错预警、含糊之处的点明、背景直觉。每条译注给出依据或标明「译者按（上下文推断，无外部依据）」。

\textbf{术语约定。}核心术语在全书首次出现时保留英文原词，其后只用中文。Springer fiber 译「Springer 纤维」，
nilpotent cone 译「幂零锥」，perverse sheaf 译「反常层」，orbital integral 译「轨道积分」，stable orbital integral
译「稳定轨道积分」，endoscopy 译「内窥」，Hitchin fibration 译「Hitchin 纤维化」，affine Grassmannian 译
「仿射 Grassmann 簇」，parahoric subgroup 保留作「parahoric 子群」。完整对照见 \verb|资料卡/04-术语文献勘误.md|。

\newpage
\tableofcontents
\newpage

% ============ 正文 ============
\setcounter{section}{-1}

\section{引言}

\subsection{本讲义的主题} 这是 2015 年 PCMI 研究生暑期学校上一门短期课程的讲义。目标有二。其一，向听众介绍一些有表示论应用的几何对象；Springer 纤维及其推广正是这样的好例子。其二，介绍轨道积分，着重它与仿射 Springer 纤维的联系，从而为 B-C.~Ng\^o 关于基本引理的短期课程提供背景材料。

本讲义的几何部分（除 \S\ref{s:orb} 以外）研究三类「纤维」：Springer 纤维、仿射 Springer 纤维与 Hitchin 纤维，复杂度依次递增。我们会讨论它们的几何性质，例如连通性与不可约分支。我们会在这些簇上构造若干群作用，并借助这些作用研究几个不平凡的例子。最重要的是，我们会研究这些纤维的上同调上的某些 Weyl 群作用，这些作用并不来自簇本身的自同构。三类纤维在表示论上的意义，以及它们之间的类比，可以概括为下表。

\begin{table}[htbp]
\centering
\begin{tabular}{|c|c|c|c|}
\hline
 & Springer 纤维 & 仿射 Springer 纤维 & Hitchin 纤维\\
\hline
域 & $k$ & 局部域 $F=k\lr{t}$ & 整体域 $k(X)$\\
\hline
对称性 & $W$ & $\tilW$ & $\tilW$\\
\hline
「扩张」对称性 & 分次 AHA & 分次 DAHA & 分次 DAHA\\
\hline
表示论 & $G(k)$ 的特征标 & $G(F)$ 的轨道积分 & $k(X)$ 上 $G$ 的迹公式\\
\hline
\end{tabular}
\end{table}

这里 AHA 指仿射 Hecke 代数，DAHA 指双仿射 Hecke 代数；$X$ 是 $k$ 上的一条代数曲线；$W$ 与 $\tilW$ 分别是 Weyl 群与扩张仿射 Weyl 群。

本讲义不追求对所有论断给出完整证明，而是着重指出有趣的现象与例子。不过，对几个关键结果我们给出大体完整的证明，例如
\begin{itemize}
\item 定理 \ref{th:Spr corr}（Springer 对应）；
\item 定理 \ref{th:ft quot}（仿射 Springer 纤维的有限性）；
\item 定理 \ref{th:orb coho}（稳定轨道积分的上同调解释）。
\end{itemize}

\subsection{对读者的预备要求} 本讲义面向对几何表示论感兴趣的研究生新生。我们假定读者了解一些基本的代数几何（概形理论、模问题、有限域上的点计数等），尽管偶尔会用到代数叠的语言。我们也假定读者了解一些 Lie 理论（代数闭域上与局部域上的约化群），但只要熟悉 $\GL_{n}$ 与 $\SL_{n}$，就足以读懂本讲义的大部分内容。

下面说明本讲义使用的上同调理论。我们处理的是任意域 $k$ 上的代数簇，而非 $\CC$ 上的，因此使用系数在 $\ell$ 进层（通常是常层 $\Ql$）中的平展上同调。我们把 $k$ 上概形 $X$ 的 $\Ql$ 系数平展上同调简记为 $\cohog{*}{X}$。不熟悉平展上同调的读者，不妨先取 $k=\CC$，把 $\cohog{*}{X}$ 理解为 $X(\CC)$ 以 $\QQ$ 为系数的奇异上同调。反常层只在 \S\ref{ss:Spr corr} 中用到。

\subsection*{致谢} 感谢 2015 年 PCMI 暑期项目的共同组织者、各位授课人与工作人员。也感谢听讲者提出的反馈。特别感谢陈竞仁（Jingren Chi）仔细通读了本讲义初稿，并提出有益的建议。

\section{第一讲：Springer 纤维}

Springer 纤维是几何表示论中经典而基本的对象。Springer \cite{Springer} 首先发现，它们的上同调群实现 Weyl 群的表示，这一现象称为 Springer 对应。作为奇异的代数簇，Springer 纤维本身即是有趣的几何对象。经由特征层，它们还与有限 Lie 型群的表示论相联系。

\subsection{设定}\label{ss:notation} 本节中 $k$ 为代数闭域。设 $G$ 是 $k$ 上连通约化群，其伴随群单（于是伴随群 $G^{\ad}$ 由七族 Dynkin 图中的一族确定）。设 $\chark$ 相对 $G$ 足够大。记 $r$ 为 $G$ 的秩。

有时方便固定 $\frg$ 的一个 Cartan 子代数 $\frt$，等价地固定一个极大环面 $T\subset G$。一旦如此，就可以谈论 $T$ 在 $\frg$ 上作用的根。

设 $\sB$ 为 $G$ 的旗簇：这是参数化 $G$ 的 Borel 子群的 $G$ 齐性射影簇。取定 Borel 子群 $B\subset G$，可把 $\sB$ 与 $G/B$ 等同。

记 $\frg=\Lie G$。对 $X\in\frg$，记 $C_{G}(X)$ 为 $X$ 在 $G$ 中的中心化子，即 $G$ 在 $\frg$ 上伴随作用下 $X$ 的稳定化子。

设 $\sN\subset\frg$ 为幂零元构成的子簇。这是一个锥：它在 $\Gm$ 按缩放作用在 $\frg$ 上时稳定。已知在伴随作用下，$\sN$ 上只有有限多个 $G$ 轨道。

\subsection{Springer 纤维}

\subsubsection{Springer 消解} 余切丛 $\tN:=T^{*}\sB$ 参数化这样的数对 $(e, B)$：$e\in\sN$，$B$ 是 $G$ 的 Borel 子群，且 $e\in \frn$，其中 $\frn$ 是 $\Lie B$ 的幂零根。{\em Springer 消解}是遗忘映射
\begin{equation*}
\pi:\tN\to\sN
\end{equation*}
它把 $(e, B)$ 送到 $e$。此映射是射影的。

对 $e\in\sN$，纤维 $\sB_{e}:=\pi^{-1}(e)$ 称为 $e$ 的 {\em Springer 纤维}。按定义，$\sB_{e}$ 是 $\sB$ 中由那些满足「$e$ 落在 $\Lie B$ 的幂零根中」的 Borel 子群 $B\subset G$ 构成的闭子概形。

\subsubsection{Grothendieck 修改}\label{sss:Gr alt}
考虑由数对 $(X,B)$ 构成的簇，其中 $X\in\frg$，$B\in\sB$，且 $X\in\Lie B$。遗忘映射 $(X,B)\mapsto X$
\begin{equation*}
\pi_{\frg}: \tg\to \frg
\end{equation*}
称为 {\em Grothendieck 修改} \footnote{术语「修改」（alteration）指源在 $k$ 上光滑的真、一般有限的映射。}，又称 {\em Grothendieck 同时消解}\zy{两名为同一映射 $\pi_{\frg}:\tg\to\frg$ 的两种通行叫法，不是两个对象：以真、一般有限的覆盖视之，称「修改」；以幂零锥上的消解视之，称「同时消解」。来源：原文 \S1.2 该段文脉（\texttt{ZhiweiYunPCMI.tex} 第 935--951 行，含作者脚注）。}。

设 $\frc=\frg\sslash G\cong \frt\sslash W$ 为 $\frg$ 在 $G$ 伴随作用下的范畴商。于是有交换图
\begin{equation}\label{tgg}
\xymatrix{\tg\ar[d]^{\pi_{\frg}}\ar[r]^{\wt\chi} & \frt\ar[d]\\
\frg \ar[r]^{\chi} & \frc}
\end{equation}
其中 $\chi:\frg\to \frc$ 是自然商映射，$\wt\chi:\tg\to \frt$ 把 $(X,B)\in\tg$ 送到 $X$ 在 $\frb/\frn$ 中的像（这里 $\frb=\Lie B$，幂零根为 $\frn$）；选取含 $\frt$ 的 Borel 子群后，$\frb/\frn$ 可与 $\frt$ 典范等同。把左列限制到 $\frg$ 的正则元上时，图 \eqref{tgg} 是笛卡儿方块 \footnote{若 $X\in\frg$ 在 $G$ 中的中心化子维数为 $r$（$G$ 的秩），则称 $X$ 为{\em 正则}元。}。特别地，若把图 \eqref{tgg} 限制到正则半单轨迹 $\frc^{\rs}\subset \frc$ \footnote{若对任意根 $\alpha$ 都有 $\alpha(X)\neq0$，则称 $X\in\frt$ 为{\em 正则半单}元。用 $\frt^{\rs}\subset\frt$ 记 $\frt$ 中正则半单元构成的开子概形。由于 $\frt^{\rs}$ 在 $W$ 下稳定，它是某个开子集 $\frc^{\rs}\subset\frc$ 的原像。开子集 $\frg^{\rs}=\chi^{-1}(\frc^{\rs})$ 按定义就是 $\frg$ 的{\em 正则半单}轨迹。}，则 $\pi_{\frg}|_{\tg^{\rs}}:\tg^{\rs}\to \frg^{\rs}$ 是 $W$-挠子；也就是说，$W$ 作用在 $\pi_{\frg}^{-1}(\frg^{\rs})$ 上，保持到 $\frg^{\rs}$ 的投影，并使 $\pi_{\frg}$ 的纤维成为 $W$ 的主齐性空间。在正则轨迹 $\frg^{\reg}\subset \frg$ 上方，$\pi_{\frg}$ 出现分支，但仍有限。

对 $X\in \frg$，把纤维 $\pi_{\frg}^{-1}(X)$ 简记为 $\tB_{X}$。把 $\pi_{\frg}$ 限制到 $\sN$ 上，幂零元 $e\in\sN$ 对应的纤维 $\tB_{e}$ 是 $\sB$ 中由满足 $e\in\Lie B$ 的 $B$ 构成的闭子簇。显然 $\sB_{e}$ 是 $\tB_{e}$ 的子概形，且两者有相同的约化结构。然而作为概形，$\tB_{e}$ 与 $\sB_{e}$ 一般不同。参见下文 \S\ref{sss:regular Spr} 与习题 \ref{exI:nonred Spr}。

\subsection{Springer 纤维的例子}

\subsubsection{} 当 $e=0$ 时，$\sB_{e}=\sB$。

\subsubsection{正则幂零元}\label{sss:regular Spr} 唯一的稠密开 $G$ 轨道由{\em 正则幂零元}组成，即满足 $\dim C_{G}(e)=r$ 的那些 $e$。当 $e$ 是正则幂零元时，$\sB_{e}$ 是单点：恰有一个 Borel 子群的 Lie 代数包含 $e$。这个 Borel 子群是什么？

由 Jacobson--Morosov 定理，可以把 $e$ 扩充为 $\frg$ 中的一个 $\sl_{2}$ 三元组 $(e,h,f)$。$h$ 在 $\frg$ 上的伴随作用有整数值的权，它把 $\frg$ 分解为权空间 $\frg(n)$，$n\in\ZZ$。令 $\frb=\oplus_{n\geq 0}\frg(n)$。这是 $\frg$ 的一个 Borel 子代数，对应的 Borel 子群就是 $\sB_{e}$ 中唯一的点。

当 $e$ 正则时，Grothendieck 修改的纤维 $\tB_{e}$ 是非约化概形，其底约化概形是单点。$\tB_{e}$ 的坐标环同构于余不变代数 $\Sym(\frt^{*})/(\Sym(\frt^{*})^{W}_{+})$；有趣的是，它也同构于复旗簇 $\calB_{\CC}$ 的 $k$ 系数上同调环。

\subsubsection{} 当 $G=\SL(V)$，$\dim V=n$ 时，$\sB$ 是完备旗 $0=V_{0}\subset V_{1}\subset V_{2}\subset\cdots\subset V_{n-1}\subset V_{n}=V$ 的模空间。Springer 纤维 $\sB_{e}$ 由满足 $eV_{i}\subset V_{i-1}$ 的旗组成。

\subsubsection{}\label{sss:SL3 subreg} 考虑 $G=\SL_{3}$，并在 $V$ 的标准基 $\{v_{1}, v_{2}, v_{3}\}$ 下取 $e=\begin{pmatrix} 0 & 0 & 1\\  0 & 0& 0 \\  0 & 0& 0 \end{pmatrix}$。此时 $\sB_{e}$ 是两个 $\PP^{1}$ 之并：第一个 $\PP^{1}$ 由旗 $0\subset V_{1}\subset \jiao{v_{1}, v_{2}}\subset V$ 组成，其中 $V_{1}$ 在固定平面 $\ker(e)=\jiao{v_{1},v_{2}}$ 内变动；第二个 $\PP^{1}$ 由旗 $0\subset \jiao{v_{1}}\subset V_{2}\subset V$ 组成，其中 $V_{2}$ 是包含 $\Im(e)=\jiao{v_{1}}$ 的变动超平面。

\subsubsection{} 考虑 $G=\SL_{4}$ 及幂零元 $e: v_{3}\mapsto v_{1}\mapsto 0, v_{4}\mapsto v_{2}\mapsto 0$，基取为 $V$ 的标准基 $\{v_{1},v_{2},v_{3},v_{4}\}$。若旗 $0\subset V_{1}\subset V_{2}\subset V_{3}\subset V$ 属于 $\sB_{e}$，则 $V_{1}\subset \ker(e)=\jiao{v_{1},v_{2}}$，且 $V_{3}\supset\Im(e)=\jiao{v_{1},v_{2}}$。记 $H:=\jiao{v_{1},v_{2}}$。分两种情形：
\begin{enumerate}
\item $V_{2}=H$。此时 $V_{1}$ 可取 $H$ 中任一直线，$V_{3}$ 可取任一包含 $V_{2}$ 的超平面。我们得到 $\sB_{e}$ 的一个闭子簇，同构于 $\PP(H)\times \PP(H)\cong \PP^{1}\times \PP^{1}$。把这个闭子簇记为 $C_{1}$。
\item $V_{2}\neq H$。这种情形定义了 $\sB_{e}$ 的一个开子概形 $U$。设直线 $V_{1}\subset H$ 由 $av_{1}+bv_{2}$ 张成，其中 $[a:b]\in\PP(H)$，则 $V_{2}$ 在 $V/H=\jiao{v_{3},v_{4}}$ 中的像由 $av_{3}+bv_{4}$ 张成。固定 $V_{1}$ 后，$V_{2}$ 的选取由 $\Hom(\jiao{av_{3}+bv_{4}}, H/V_{1})\cong\Hom(V_{1},H/V_{1})$ 给出。一旦 $V_{2}$ 固定，$V_{3}=V_{2}+H$ 也随之固定。因此 $U$ 同构于 $\PP(H)\cong \PP^{1}$ 上直线丛 $\calO(2)$ 的全空间。
\end{enumerate}
由上述讨论可见 $\sB_{e}$ 的维数为 $2$，$C_{1}$ 是 $\sB_{e}$ 的不可约分支，$U$ 的闭包也是，记后者为 $C_{2}$。于是 $C_{1}\cup C_{2}=\sB_{e}$，而 $C_{1}\cap C_{2}$ 是 $C_{1}\cong\PP^{1}\times\PP^{1}$ 中的对角线。

\subsubsection{$A$ 型 Springer 纤维的分支}\label{sss:Spr type A} 当 $G=\SL_{n}=\SL(V)$ 时，Spaltenstein \cite{Spa76} 与 Steinberg \cite[]{St} 用 $n$ 阶标准 Young 表给出了 $\sB_{e}$ 的不可约分支的刻画。这与后文讨论的 Springer 对应有关，见 \S\ref{sss:A Spr corr}。下面我们遵循 \cite[Ch II, \S5]{Spaltenstein} 的叙述。

固定幂零元 $e\in\sN$，其 Jordan 型是 $n$ 的分拆 $\l$。这就是说，若分拆写成 $n=\l_{1}+\l_{2}+\cdots$，则 $e$ 的 Jordan 块大小依次为 $\l_{1},\l_{2},\cdots$。我们将构造一个（非代数的）映射 $\sB_{e}\to ST(\l)$，其中 $ST(\l)$ 是分拆 $\l$ 的标准 Young 表构成的离散集。对每个满足 $eV_{i}\subset V_{i-1}$ 的完备旗 $0=V_{0}\subset V_{1}\subset\cdots\subset V_{n-1}\subset V_{n}=V$，$e$ 诱导出 $V/V_{n-i}$ 的幂零自同态。设 $\mu_{i}$ 是 $e$ 在 $V/V_{n-i}$ 上的 Jordan 型，则 $\mu_{i}$ 是 $i$ 的分拆。分拆的递增序列 $\mu_{1},\mu_{2},\cdots, \mu_{n}=\l$ 满足：$\mu_{i}$ 由 $\mu_{i-1}$ 把某一部分加 $1$ 得到（包括新增一个大小为 $1$ 的部分）。由此得到 $\l$ 的 Young 图 $Y(\l)$ 的子图递增序列 $Y_{1},\cdots, Y_{n}=Y(\l)$。把 $Y_{i}-Y_{i-1}$ 中唯一的格子标上 $i$，就得到一张标准 Young 表。

Spaltenstein \cite[Ch II, Prop 5.5]{Spaltenstein} 证明：每张标准 Young 表在 $\sB_{e}$ 中的原像的闭包都是不可约分支。而且 $\sB_{e}$ 的全部不可约分支都这样得到，维数都等于 $d_{e}=\frac{1}{2}\sum_{i}\l^{*}_{i}(\l^{*}_{i}-1)$，其中 $\l^{*}$ 是 $\l$ 的共轭分拆。特别地，最高维上同调 $\cohog{2d_{e}}{\sB_{e}}$ 的维数等于 $\#ST(\l)$，后者也是对称群 $S_{n}$ 的一个不可约表示的维数。这一陈述是 Springer 对应的数值影子：后者断言 $\cohog{2d_{e}}{\sB_{e}}$ 自然地是对称群 $S_{n}$ 的不可约表示。

Spaltenstein \cite[Ch II, Prop 5.9]{Spaltenstein} 还证明：存在 $\sB$ 到仿射空间的分层，使 $\sB_{e}$ 是若干层之并。这说明上同调上的限制映射 $\cohog{*}{\sB}\to \cohog{*}{\sB_{e}}$ 是满射。

\subsubsection{}\label{sss:Sp4} 考虑 $G=\Sp(V)$，其中 $V$ 是 $2n$ 维辛向量空间。此时 $\sB$ 是满足 $V^{\bot}_{i}=V_{2n-i}$（$i=1,\cdots, n$）的完备旗
\begin{equation*}
0=V_{0}\subset V_{1}\subset V_{2}\subset\cdots\subset V_{n}\subset\cdots\subset V_{2n-1}\subset V_{2n}=V
\end{equation*}
的模空间。Springer 纤维 $\sB_{e}$ 由那些对所有 $i$ 满足 $eV_{i}\subset V_{i-1}$ 的旗组成。

考虑 $\dim V=4$ 的情形。选取 $V$ 的基 $\{v_{1},v_{2},v_{3},v_{4}\}$，使 $V$ 上的辛形式 $\om$ 满足：当 $i=1,2$ 时 $\om(v_{i},v_{5-i})=1$，而当 $i+j\neq 5$ 时 $\om(v_{i},v_{j})=0$。设 $e$ 是 $\frg=\sp(V)$ 中由 $e:v_{4}\mapsto v_{1}\mapsto 0, v_{3}\mapsto 0,v_{2}\mapsto 0$ 给出的幂零元。此时 $\sB_{e}$ 中的旗 $0\subset V_{1}\subset V_{2}\subset V_{1}^{\bot}\subset V$ 必须满足 $\jiao{v_{1}}\subset V_{2}\subset\jiao{v_{1},v_{2},v_{3}}$，而这也是它属于 $\sB_{e}$ 的唯一条件（习题 \ref{exI:211Sp4}）。这样的 $V_{2}$ 对应于直线 $V_{2}/\jiao{v_{1}}\subset \jiao{v_{2},v_{3}}$，即 $\PP^{1}=\PP(\jiao{v_{2},v_{3}})$ 中一点。在此 $\PP^{1}$ 之上有万有秩二丛 $\calV_{2}$，它在 $V_{2}/\jiao{v_{1}}$ 处的纤维是二维空间 $V_{2}$。再选取 $V_{1}$，相当于在 $\calV_{2}$ 的射影化中取一点。正合列 $0\to\jiao{v_{1}}\to V_{2}\to V_{2}/\jiao{v_{1}}\to0$ 给出 $\PP^{1}$ 上向量丛的正合列 $0\to\calO\to\calV_{2}\to\calO(-1)\to0$。因此 $\calV_{2}$ 同构于 $\calO(-1)\oplus\calO$，而 $\sB_{e}\cong \PP(\calO(-1)\oplus\calO)$ 是一张 Hirzebruch 曲面。

\subsubsection{次正则 Springer 纤维}\label{sss:subreg} \S\ref{sss:SL3 subreg} 中的例子是{\em 次正则 Springer 纤维}的最简单情形。$\sN$ 中恰有一个余维 $2$ 的幂零轨道 $\calO_{\subreg}$，称为{\em 次正则幂零轨道}。对 $e\in\calO_{\subreg}$，已知 $\sB_{e}$ 是若干 $\PP^{1}$ 之并，其构型如下。可以对 $\sB_{e}$ 作对偶图：顶点为 $\sB_{e}$ 的不可约分支，若两个分支相交（实际上交于单点）则连一条边。

为简单起见设 $G$ 为伴随型。设 $G'$ 是另一个伴随单群，其型如下定义。若 $G$ 单链，取 $G'=G$。若 $G$ 是 $B_{n}$、$C_{n}$、$F_{4}$、$G_{2}$ 型，则分别取 $G'$ 为 $A_{2n-1}$、$D_{n+1}$、$E_{6}$、$D_{4}$ 型。可以证明 $\sB_{e}$ 总是若干 $\PP^{1}$ 之并，其对偶图是 $G'$ 的 Dynkin 图。非单链情形的规则是：每个长单根对应 $2$ 个或 $3$ 个 $\PP^{1}$，每个短单根对应唯一一个 $\PP^{1}$。这样一组 $\PP^{1}$ 称为 Dynkin 曲线，参见 \cite[\S3.10, Definition and Prop 2]{StConj} 与 \cite[\S6.3]{Slodowy}。

例如，当 $G$ 为 $A_{n}$ 型时，$\sB_{e}$ 是由 $n$ 个 $\PP^{1}$ 组成的链：$\sB_{e}=C_{1}\cup C_{2}\cup\cdots\cup C_{n}$，其中 $C_{i}\cap C_{i+1}$ 为一点，其余两两不交。

Brieskorn \cite{Brieskorn} 遵循 Grothendieck 的建议，把幂零锥沿次正则轨道的奇点与 Klein 奇点联系起来，并且在 $\frg$ 内部实现了该奇点的半通用形变。设 $G$ 单链。可以构造过 $e$ 的二维横截片 $S_{e}\subset\sN$，使得 $S_{e}$ 除 $e$ 外都由正则元组成。此时 $S_{e}$ 是正规曲面，在 $e$ 处有 Klein 奇点，其型与 $G$ 的 Dynkin 图相同。\footnote{Klein 奇点指与 $\AA^{2}$ 被 $\SL_{2}$ 的有限子群作商后在 $(0,0)$ 处的奇点解析同构的曲面奇点。} 其原像 $\tilS_{e}:=\pi^{-1}(S_{e})\subset\tN$ 恰是 $S_{e}$ 的极小消解，因此 $\sB_{e}$ 是例外除子之并。把 $\sB_{e}$ 的分支与 $G$ 的单根对应起来后，例外除子的相交矩阵恰为 $G$ 的 Cartan 矩阵的负矩阵。Slodowy \cite{Slodowy} 把上述图景推广到非单链群；关于次正则轨道与 Klein 奇点之间联系的精美叙述，参见他的书 \cite{Slodowy}。

\subsection{Springer 纤维的几何性质}

\subsubsection{连通性} Springer 纤维 $\sB_{e}$ 是连通的。参见 \cite[p.131, Prop 1]{StConj}、\cite[Ch II, Cor 1.7]{Spaltenstein} 以及习题 \ref{exI:conn}。

\subsubsection{等维性} Spaltenstein \cite{Spa77}、\cite[Ch II, Prop 1.12]{Spaltenstein} 证明：$\sB_{e}$ 的所有不可约分支维数相同。

\subsubsection{维数公式} 设 $d_{e}$ 为 $\sB_{e}$ 的维数。Steinberg \cite[Thm 4.6]{St} 与 Springer \cite{Springer} 证明
\begin{equation}\label{Spr dim}
d_{e}=\frac{1}{2}(\dim\sN-\dim G\cdot e)=\frac{1}{2}(\dim C_{G}(e)-r).
\end{equation}

\subsubsection{中心化子的作用} 群 $G\times \Gm$ 作用在 $\sN$ 上（其中 $\Gm$ 通过缩放作用）。设 $\tilG_{e}=\Stab_{G\times \Gm}(e)$ 为稳定化子，则 $\tilG_{e}$ 作用在 $\sB_{e}$ 上。注意 $\tilG_{e}$ 总是满射到 $\Gm$，其核为 $C_{G}(e)$。

$\tilG_{e}$ 在 $\cohog{*}{\sB_{e}}$ 上的作用经由有限群 $A_{e}=\pi_{0}(\tilG_{e})=\pi_{0}(G_{e})$ 分解。注意 $A_{e}$ 不仅依赖于 $G$ 的同源类，还依赖于 $G$ 的同构类。例如当 $e$ 正则时，$A_{e}=\pi_{0}(ZG)$，其中 $ZG$ 是 $G$ 的中心。$A_{e}$ 在 $\cohog{*}{\sB_{e}}$ 上的作用进一步经由 $A_{e}$ 在伴随群 $G^{\ad}$ 中的像分解。

\subsubsection{纯性} Springer \cite{SpringerPure} 证明：$\sB_{e}$ 的上同调总是纯的（当 $k=\CC$ 时按 Hodge 理论的意义，当 $k$ 为有限域时按 Frobenius 权的意义）。

\subsubsection{} 设 $e\in\sN$。考虑由包含 $\sB_{e}\incl\sB$ 诱导的限制映射 $i^{*}_{e}: \cohog{*}{\sB}\to \cohog{*}{\sB_{e}}$。则 $i^{*}_{e}$ 的像恰为 $\cohog{*}{\sB_{e}}$ 在 $A_{e}$ 下的不变量。这是 Hotta 与 Springer 的定理 \cite[Theorem 1.1]{HottaSpringer}。特别地，当 $G$ 为 $A$ 型时，这样的限制映射总是满射。

\subsubsection{奇偶消没} De Concini、Lusztig 与 Procesi \cite{DLP} 证明：对一切奇数 $i$ 与任意 $e\in\sN$，$\cohog{i}{\sB_{e}}$ 消没。当 $k=\CC$ 时他们还证明了更强的结论：奇数 $i$ 时 $\cohog{i}{\sB_{e},\ZZ}$ 消没，偶数 $i$ 时无挠。

\subsection{Springer 对应}\label{ss:Spr corr} 设 $W$ 为 $G$ 的 Weyl 群。1976 年，Springer \cite{Springer} 作出了根本性的观察：$\cohog{*}{\sB_{e}}$ 上存在自然的 $W$ 作用，尽管 $W$ 并不作为簇的自同构作用在 $\sB_{e}$ 上。

\begin{theorem}[Springer {\cite[Thm 6.10]{Springer}}]\label{th:Spr corr}
\begin{enumerate}
\item 对每个幂零元 $e$，$\cohog{*}{\sB_{e}}$ 上有自然的 $W$ 分次作用，且与 $A_{e}$ 的作用交换。
\item 对每个幂零元 $e$ 与 $A_{e}$ 的每个不可约表示 $\rho$，重数空间 $M(e,\rho):=\Hom_{A_{e}}(\rho, \cohog{2d_{e}}{\sB_{e}})$ 或为零，或在 (1) 中的作用之下是 $W$ 的不可约表示。
\item $W$ 的每个不可约表示 $\chi$ 都在 $G$ 共轭意义下唯一地以 $M(e,\rho)$ 的形式出现。于是对应 $\chi\mapsto (e,\rho)$ 给出单射
\begin{equation}\label{Spr corr}
\Irr(W)\incl \{(e,\rho)\}/G.
\end{equation}
\end{enumerate}
\end{theorem}

\subsubsection{约定}\label{sss:two actions} 事实上 $\cohog{*}{\sB_{e}}$ 上有两个自然的 $W$ 作用，二者相差与 $W$ 的符号表示作张量。本讲义采用由以下性质规范的那个作用：$W$ 的平凡表示对应正则幂零元 $e$ 与平凡 $\rho$；$W$ 的符号表示对应 $e=0$。但要注意，Springer 的原始论文 \cite{Springer} 用的是另一个作用。

\subsubsection{$e=0$ 的情形} 取 $e=0$，Springer 定理给出 $W$ 在 $\cohog{*}{\calB}$ 上的分次作用。这是什么作用？首先，可以改用 $G/T$ 而非 $\sB=G/B$ 来几何地看出这个作用。事实上，由于 $N_{G}(T)/T=W$，$N_{G}(T)$ 在 $G/T$ 上的右作用诱导出 $W$ 在 $G/T$ 上的作用，进而诱导出 $W$ 在 $\cohog{*}{G/T}$ 上的作用。由于投影 $G/T\to G/B$ 是仿射空间丛，有 $\cohog{*}{\sB}\cong\cohog{*}{G/T}$。可以证明，在此同构之下，$W$ 在 $\cohog{*}{G/T}$ 上的作用恰对应 Springer 在 $\cohog{*}{\calB}$ 上的作用。

设 $S=\Sym(\xch(T)\otimes\Ql)$ 为分次对称代数，其中 $\xch(T)$ 的次数为 $2$。$W$ 在 $\xch(T)$ 上的反射表示随即诱导出 $W$ 在 $S$ 上的分次作用。回忆旗簇上同调环的 Borel 表示
\begin{equation}\label{Borel HB}
\cohog{*}{\calB,\Ql}\cong S/(S^{W}_{+})
\end{equation}
其中 $S_{+}\subset S$ 是由正次数元素张成的理想，$(S^{W}_{+})$ 表示 $S$ 中由 $S_{+}$ 的 $W$ 不变量生成的理想。其实 \eqref{Borel HB} 是 $W$ 模的同构（见 \cite[Prop 7.2]{Springer}）。由 Chevalley 的定理，$S/(S^{W}_{+})$ 同构于 $W$ 的正则表示，因此作为 $W$ 模，$\cohog{*}{\sB}$ 也同构于 $W$ 的正则表示。

\begin{remark} \eqref{Spr corr} 中的目标集可以与幂零轨道上不可约 $G$ 等变局部系统的同构类典范地等同。事实上，对幂零轨道 $\calO\subset\sN$ 上不可约 $G$ 等变局部系统 $\calL$，它在 $e\in\calO$ 处的茎给出中心化子 $C_{G}(e)$ 的一个不可约表示 $\rho$，该表示经由 $A_{e}$ 分解。注意当 $G$ 在固定同源类中变动时，$G$ 等变的概念随之改变。若给 $\Irr(W)$ 补上若干较小 Weyl 群的 $\Irr(W')$，可以把上述单射 \eqref{Spr corr} 扩充为双射。这就是 Lusztig \cite{LuIC} 发现的{\em 广义 Springer 对应}。
\end{remark}

Springer 对定理 \ref{th:Spr corr} 的原始证明使用 $\frg(\FF_{q})$ 上的三角和；当 $k$ 的特征为零时，他的证明还用到向有限域的归约。下述 Borho 与 MacPherson 的定理 \cite{BM} 可以对所有大特征或特征零的基域 $k$ 直接证明 Springer 对应。为陈述它，需要可构造 $\Ql$ 层（复形）与反常层的语言，参见标准文献 \cite{BBD} 与本卷中 de Cataldo 的讲义 \cite{dC}。

\begin{theorem}\label{th:BM} 复形 $\calS:=\bR\pi_{*}\Ql[\dim\sN]$ 是 $\sN$ 上的反常层，其自同态环典范同构于群代数 $\Ql[W]$。特别地，$W$ 作用在 $\bR\pi_{*}\Ql$ 的各茎上，即 $W$ 作用在 $\cohog{*}{\sB_{e}}$ 上，对一切 $e\in\sN$。
\end{theorem}

下面概述 $W$ 在 $\bR\pi_{*}\Ql[\dim\sN]$ 上作用的三种构造。

\subsubsection{经由中间扩充的构造}\label{sss:mid} 这一构造（更确切地说，是把 $\frg$ 换成 $G$ 的版本）属于 Lusztig \cite[\S3]{LuGreen}。Springer 纤维的维数公式 \eqref{Spr dim} 蕴含
\begin{itemize}
\item 映射 $\pi:\tN\to\sN$ 是半小的。\footnote{对不可约簇的真满射 $f:X\to Y$，若对任意 $d\geq 1$，集合 $\{y\in Y|\dim f^{-1}(y)\geq d\}$ 在 $Y$ 中的余维至少为 $2d$（相应地 $2d+1$），则称 $f$ 为{\em 半小}（相应地{\em 小}）映射。}
\end{itemize}
对 $\tB_{X}$ 的维数，公式 \eqref{Spr dim} 有适用于一切 $X\in \frg$ 的推广。用这个公式可以证明
\begin{itemize}
\item 映射 $\pi_{\frg}:\tg\to\frg$ 是小映射。
\end{itemize}
反常层理论中有一个众所周知的事实：$\pi_{\frg}$ 的小性（连同 $\tg$ 光滑这一事实）蕴含 $\calS_{\frg}:=\bR\pi_{\frg,*}\Ql[\dim\frg]$ 是反常层，且是它在 $\frg$ 任一稠密开子集上的限制的中间扩充。在正则半单轨迹 $\frg^{\rs}\subset \frg$ 上方，$\pi_{\frg}$ 是 $W$-挠子，因此 $\calS_{\frg}|_{\frg^{rs}}$ 是次数移位 $-\dim\frg$ 的局部系统，并带有 $W$ 的作用。由中间扩充的函子性，$\calS_{\frg}$ 带有 $W$ 的作用。取 $\calS_{\frg}$ 的茎，就得到 $W$ 在 $\cohog{*}{\tB_{X}}$ 上的作用，对一切 $X\in\frg$。

特别地，对幂零元 $e$，我们得到 $W$ 在 $\cohog{*}{\sB_{e}}=\cohog{*}{\tB_{e}}$ 上的作用，因为 $\tB_{e}$ 与 $\sB_{e}$ 有相同的约化结构。这正是 Springer 在其原始论文 \cite{Springer} 中定义的作用，{\em 它与我们的作用相差与 $W$ 的符号特征作张量}。

\subsubsection{经由 Fourier 变换的构造}\label{sss:Four} 由 $\pi$ 的半小性，复形 $\calS=\bR\pi_{*}\Ql[\dim\sN]$ 也是反常层。然而它不是从 $\calN$ 的某个开子集作中间扩充得来的。仿射空间上 $\Gm$ 等变层有 Fourier 变换的概念 \cite{Laumon}。可以证明 $\calS$ 同构于 $\calS_{\frg}$ 的 Fourier 变换，反之亦然。于是由 Fourier 变换的函子性，$\calS_{\frg}$ 上的 $W$ 作用诱导出 $\calS$ 上的 $W$ 作用。取 $\calS$ 在 $e$ 处的茎，就得到 $W$ 在 $\cohog{*}{\sB_{e}}$ 上的作用。按 \S\ref{sss:two actions} 的约定，这个作用是规范化的。

\subsubsection{经由对应的构造} 考虑 Steinberg 簇 $\St_{\frg}=\tg\times_{\frg}\tg$，它分类三元组 $(X,B_{1},B_{2})$，其中 $B_{1},B_{2}$ 是 $G$ 的 Borel 子群，且 $X\in\Lie B_{1}\cap\Lie B_{2}$。$\St_{\frg}$ 的不可约分支以 Weyl 群的元素为指标：对 $w\in W$，令 $\St_{w}$ 为 $w$ 在 $\tg^{\rs}$ 上作用的图像之闭包，则 $\St_{w}$ 是 $\St$ 的不可约分支，且这些穷尽 $\St$ 的全部不可约分支。上同调对应的形式体系使我们能从每个 $\St_{w}$ 得到复形 $\calS_{\frg}=\bR\pi_{\frg,*}\Ql[\dim\frg]$ 的一个自同态。要证明这些自同态合起来构成 $W$ 在 $\calS_{\frg}$ 上的作用，并不显然。论证的关键仍是映射 $\pi_{\frg}$ 的小性。定义了 $\calS_{\frg}$ 上的 $W$ 作用之后，可以用 \S\ref{sss:mid} 的办法，把 $W$ 在茎 $\calS_{\frg, e}$ 上的作用与符号表示作张量，也可以用 \S\ref{sss:Four} 的 Fourier 变换，来定义 $\cohog{*}{\sB_{e}}$ 上的 Springer 作用。关于这一构造的讨论参见 \cite[Remark 3.3.4]{GS}；亦可参见 \cite[\S3.4]{CG}，那里用 $\St_{w}$ 在幂零 Steinberg 簇 $\tN\times_{\sN}\tN$ 中的极限给出了类似但不同的构造。

注意上述三种构造都能证明 $\End(\calS)\cong\Ql[W]$，从而给出定理 \ref{th:BM} 的证明。

\subsubsection{经由单值的构造} 我们概述 Slodowy \cite[\S4]{Slo4} 的一个构造，它适用于 $k=\CC$。当时猜想这一构造给出的 $W$ 在 $\cohog{*}{\sB_{e}}$ 上的作用与定理 \ref{th:Spr corr} 中的相同。Rossmann 在 \cite[\S2]{Rossmann} 中给出了类似的构造，他把自己的作用等同于 Kazhdan 与 Lusztig 在 \cite{KL80} 中构造的作用，而后者已知与 Springer 的作用相同。因此这些构造给出的 $W$ 作用都与定理 \ref{th:Spr corr} 中的相同。

设 $e\in\calN$，设 $S_{e}\subset \frg$ 是 $e$ 的轨道的横截片。选取包含 $e$ 的 $\sl_{2}$ 三元组 $(e,h,f)$ 后，有这样的横截片的标准取法：$S_{e}=e+\frg_{f}$，其中 $\frg_{f}$ 是 $f$ 在 $\frg$ 中的中心化子。现考虑下图，其中除最右方块外各方块都是笛卡儿的
\begin{equation}\label{Slo}
\xymatrix{\sB_{e}\ar[d]\ar@{^{(}->}[r] & \tilS^{\nil}_{e}\ar[d]\ar@{^{(}->}[r] & \tilS_{e}\ar[d]\ar@{^{(}->}[r] & \tg\ar[d]^{\pi_{\frg}}\ar[r]^{\wt\chi} & \frt\ar[d]\\
\{e\}\ar@{^{(}->}[r] & S^{\nil}_{e}\ar@{^{(}->}[r] & S_{e}\ar@{^{(}->}[r] & \frg\ar[r]^{\chi} & \frc}
\end{equation}
这里最右方块是 \eqref{tgg}，$S^{\nil}_{e}=S_{e}\cap\sN$，而 $\tilS_{e}$ 与 $\tilS^{\nil}_{e}$ 分别是 $S_{e}$ 与 $S^{\nil}_{e}$ 在 $\pi_{\frg}$ 下的原像。取 $e$ 处的小球 $V_{e}\subset S_{e}$，再取 $0\in\frc$ 处更小的球 $V_{0}\subset \frc$。令 $U_{e}=V_{e}\cap \chi^{-1}(V_{0})\subset S_{e}$。则图 \eqref{Slo} 限制为图
\begin{equation}\label{Slo U}
\xymatrix{\sB_{e}\ar[d]\ar@{^{(}->}[r] & \tilU^{\nil}_{e}\ar[d]\ar@{^{(}->}[r] & \tilU_{e}\ar[d]\ar[r]^{\wt\chi_{e}} & \tilV_{0}\ar[d]\\
\{e\}\ar@{^{(}->}[r] & U^{\nil}_{e}\ar@{^{(}->}[r] & U_{e}\ar[r] & V_{0}}
\end{equation}
这里 $U^{\nil}_{e}=U_{e}\cap\sN$，而 $\tilV_{0}, \tilU_{e}$ 与 $\tilU^{\nil}_{e}$ 分别是 $V_{0}, U_{e}$ 与 $U^{\nil}_{e}$ 在 \eqref{Slo} 中竖直映射下的原像。关键的拓扑事实有两条
\begin{itemize}
\item 包含 $\sB_{e}\incl \tilU^{\nil}_{e}$ 有形变收缩，因此是同伦等价；
\item 映射 $\wt\chi_{e}: \tilU_{e}\to \tilV_{0}$ 是可平凡化的纤维丛（微分拓扑意义下）。
\end{itemize}
现在，由第二条事实，$\wt\chi_{e}$ 的一般纤维带有一个同伦意义下的 $W$ 作用：因为 \eqref{Slo U} 的最右方块在 $V_{0}\cap\frc^{\rs}$ 上方是笛卡儿的，且 $\tilV_{0}\to V_{0}$ 在 $V_{0}\cap\frc^{\rs}$ 上方是 $W$-挠子。由第一条事实，$\sB_{e}$ 与 $\tilU^{\nil}_{e}=\wt\chi^{-1}_{e}(0)$ 同伦等价；又因 $\wt\chi_{e}$ 是纤维丛，$\sB_{e}$ 也与 $\wt\chi_{e}$ 的一般纤维同伦等价。把这些事实合起来，就得到 $W$ 在 {\em $\sB_{e}$ 的同伦型}上的作用，这比 $W$ 作用在 $\sB_{e}$ 的上同调上更强。此构造的一个推论是：定理 \ref{th:Spr corr} 中 $\cohog{*}{\sB_{e}}$ 上的 $W$ 作用保持环结构。

\subsubsection{在定理 \ref{th:BM} 之下证明定理 \ref{th:Spr corr}}
把反常层 $\calS$ 按 $W$ 作用分解为同型分量
\begin{equation*}
\calS=\bigoplus_{\chi\in\Irr(W)}V_{\chi}\otimes\calS_{\chi}
\end{equation*}
其中 $V_{\chi}$ 是 $W$ 通过不可约表示 $\chi$ 作用的那个空间，而 $\calS_{\chi}=\Hom_{W}(V_{\chi},\calS)$ 是 $\sN$ 上的反常层。由 $\End(\calS)\cong\Ql[W]$ 可得每个 $\calS_{\chi}$ 都非零，且
\begin{equation}\label{HomSS}
\Hom(\calS_{\chi},\calS_{\chi'})=\begin{cases}\Ql & \chi=\chi' \\ 0 & \text{其他情形}\end{cases}
\end{equation}
分解定理 \cite[Th 6.2.5]{BBD} 蕴含每个 $\calS_{\chi}$ 都是半单反常层。因此 \eqref{HomSS} 说明 $\calS_{\chi}$ 是单的反常层。于是 $\calS_{\chi}$ 形如 $\IC(\ov\calO, \calL)$，其中 $\calO\subset\sN$ 是幂零轨道，$\calL$ 是 $\calO$ 上不可约 $G$ 等变局部系统。再者，由于 $\chi\neq\chi'$ 时 $\Hom(\calS_{\chi},\calS_{\chi'})=0$，这些单反常层 $\{\calS_{\chi}\}_{\chi\in\Irr(W)}$ 互不同构。把 \eqref{Spr corr} 的右边理解为幂零轨道上不可约 $G$ 等变局部系统的同构类，就证明了定理 \ref{th:Spr corr} 的第 (3) 部分。若 $\calS_{\chi}=\IC(\ov\calO,\calL)$ 且 $e\in\calO$，则 $\pi$ 的半小性使我们可以把茎 $\calL_{e}$ 与 $\cohog{2d_{e}}{\sB_{e}}$ 的一个 $A_{e}$ 同型子空间等同起来。这就证明了定理 \ref{th:Spr corr} 的第 (2) 部分。\qed

下面再举几个 Springer 对应的例子。

\subsubsection{$A$ 型}\label{sss:A Spr corr} 当 $G=\SL_{n}$ 时，所有 $A_{e}$ 都平凡。Springer 对应给出 $W=S_{n}$ 的不可约表示与 $\frg=\sl_{n}$ 的幂零轨道之间的双射，两者都由 $n$ 的分拆参数化。在 \S\ref{sss:Spr type A} 中已经看到：若 $e$ 的 Jordan 型为 $\l$，则最高维上同调 $\cohog{2d_{e}}{\sB_{e}}$ 有一组以 $\l$ 的标准 Young 表为指标的基，而后者也给出 $S_{n}$ 对应于分拆 $\l$ 的不可约表示的一组基。例如，对 $G=\SL_{3}$，Springer 对应为
\begin{itemize}
\item 平凡表示 $\bij$ 正则轨道，分拆 $3=3$；
\item 二维表示 $\bij$ 次正则轨道，分拆 $3=2+1$；
\item 符号表示 $\bij$ $\{0\}$，分拆 $3=1+1+1$。
\end{itemize}

\subsubsection{次正则轨道与反射表示}\label{sss:subreg Spr corr} 考虑 $e$ 是次正则幂零元的情形。此时分支群 $A_{e}$ 可以与 \S\ref{sss:subreg} 中引入的 $G'$ 的 Dynkin 图的自同构群等同（见 \cite[\S7.5, Proposition]{Slodowy}）。把 $\sB_{e}$ 的不可约分支与 $G'$ 的 Dynkin 图的顶点等同之后，$A_{e}$ 在 $\cohog{2}{\sB_{e}}$ 上的作用以「重排由不可约分支给出的基」的方式实现，与它在 $G'$ 的 Dynkin 图上的作用相同。例如当 $G=G_{2}$ 时，可写 $\sB_{e}=C_{1}\cup C_{2}\cup C_{3}\cup C_{4}$，其中 $C_{1}$、$C_{2}$、$C_{3}$ 各自与 $C_{4}$ 交于一点，其余两两不交。群 $A_{e}$ 同构于 $S_{3}$，它在 $\cohog{2}{\sB_{e}}$ 上的作用固定 $C_{4}$ 的基本类，并置换 $C_{1}, C_{2}, C_{3}$ 的基本类。

注意 $\cohog{2}{\sB_{e}}^{A_{e}}$ 的维数总是 $r$，即 $G$ 的秩。事实上，作为 $W$ 模，$\cohog{2}{\sB_{e}}^{A_{e}}$ 同构于 $W$ 在 $\frt^{*}$ 上的反射表示。换言之，在 Springer 对应之下，数对 $(e=\text{次正则}, \rho=1)$ 对应 $W$ 的反射表示。

\subsection{评注与推广}

\subsubsection{扩张的对称性} 若改用更精细的上同调理论，$\cohog{*}{\sB_{e}}$ 上的 $W$ 作用可以用多种方式扩张为更大代数的作用。在等变上同调 $\upH^{*}_{\tilG_{e}}(\sB_{e})$ 上，有{\em 分级仿射 Hecke 代数}的作用（见 Lusztig \cite{LuGrHk}）。在 $\sB_{e}$ 的 $\tilG_{e}$ 等变 $K$ 群上，有{\em 仿射 Hecke 代数}的作用（见 Kazhdan--Lusztig \cite{KLAffHk} 与 Chriss--Ginzburg \cite{CG}）。

\subsubsection{群版本} 把 $\sN$ 与 $\frg$ 分别换成幂幺簇 $\sU\subset G$ 与 $G$ 本身，Springer 消解与 Grothendieck 修改就有显然的类比物。当 $\chark$ 很大时，指数映射以 $G$ 等变的方式把 $\sN$ 与 $\sU$ 等同，因此幂零元的 Springer 纤维理论与幂幺元的相应理论完全一样。反常层 $\calS_{\frg}$ 及其不可约直和项的群版本，是{\em 特征层}的先声；这一理论由 Lusztig（\cite{LuChShI}、\cite{LuChShII}、\cite{LuChShIV} 与 \cite{LuChShV}）发展起来，用于研究有限群 $G(\FF_{q})$ 的特征标。

\subsubsection{部分 Springer 消解}\label{sss:partial Spr} 可以在部分旗簇中定义 $\sB_{e}$ 的类比物。设 $\sP$ 是 $G$ 的部分旗簇，分类给定类型的抛物子群 $P\subset G$。可以考虑映射 $\pi:\tN\to\sN$ 的两种类比。

第一种，把 $\tN=T^{*}\sB$ 换成 $T^{*}\sP$，它分类满足 $e\in \Lie\frn_{P}$ 的数对 $(e, P)\in\sN\times\sP$，其中 $\frn_{P}$ 是 $\Lie P$ 的幂零根。设 $\tau_{\sP}: T^{*}\sP\to\sN$ 为第一个投影。一般说来此映射不是满射，其像是某个幂零轨道 $\calO_{\calP}$ 的闭包。轨道 $\calO_{P}$ 的特征性质是：对任意 $P\in\sP$，它与 $\frn_{P}$ 的交在 $\frn_{P}$ 中稠密。它称为与类型 $\sP$ 的抛物子群相联系的{\em Richardson 类}。当 $G$ 为 $A$ 型时，每个幂零类 $\calO$ 都是某个类型 $\sP$ 的抛物子群所联系的 Richardson 类（一般并不唯一）。映射 $T^{*}\sP\to \ov\calO$ 是奇点的消解。对一般的 $G$，并非每个幂零轨道都是 Richardson 类。

第二种，可以考虑分类数对 $(e,P)$ 的子概形 $\tN_{\sP}\subset\sN\times \sP$，其中 $e\in\sN_{P}$，而 $\sN_{P}\subset\Lie P$ 是 $P$ 的幂零锥。此时投影 $\pi_{\sP}: \tN_{\sP}\to \sN$ 是满射，并且是奇点的部分消解。Springer 消解 $\pi$ 可以分解为
\begin{equation*}
\pi: \tN=\tN_{\sB}\xrightarrow{\nu_{\sP}}\tN_{\sP}\xrightarrow{\pi_{\sP}} \sN.
\end{equation*}

有嵌入 $T^{*}\sP\incl \tN_{\sP}$。可以把 $\tau_{\sP}$ 或 $\pi_{\sP}$ 的纤维看作 Springer 纤维的抛物类比物，称之为{\em 部分 Springer 纤维}。$W$ 在 $\sB_{e}$ 上同调上的 Springer 作用，在部分 Springer 纤维上有类比物。更多的信息参见 \cite{BM}。

\subsubsection{Hessenberg 簇} Grothendieck 修改 $\pi_{\frg}:\tg\to \frg$ 有一个推广，其中 $\frg$ 被换成 $G$ 的另一个线性表示。

固定 $G$ 的 Borel 子群 $B$。设 $(V,\rho)$ 是 $G$ 的表示，$V^{+}\subset V$ 是在 $B$ 下稳定的子空间。现在用数对 $(V,V^{+})$ 代替 $(\frg,\frb)$，就得到 Grothendieck 修改的推广。确切地说，设 $\tilV\subset V\times \calB$ 是由满足 $v\in \rho(g)V^{+}$ 的数对 $(v, gB)\in V\times\sB$ 构成的子概形。设 $\pi_{V}:\tilV\to V$ 为第一个投影。$\pi_{V}$ 的纤维称为 {\em Hessenberg 簇}。

我们将在 \S\ref{ss:aff examples} 看到，Hessenberg 簇在研究某些仿射 Springer 纤维时自然出现。关于 Hessenberg 簇的更多信息，见 \cite{GKM} 与 \cite{OY}。

\subsection{习题}

\subsubsection{} 对 $G=\SL_{n}$，确定 $\frg$ 中正则幂零元与次正则幂零元的 Jordan 块大小。

\subsubsection{}\label{exI:nonred Spr} 对 $G=\SL_{2}$ 与 $\SL_{3}$，计算正则幂零元 $e$ 对应的{\em 非约化} Springer 纤维 $\tB_{e}$ 的坐标环。并证明 Springer 纤维 $\sB_{e}$ 确实是约化单点。

{\em 提示}：若把 $e$ 写成上三角矩阵，则 $\tB_{e}$ 落在旗簇 $\sB$ 的大 Bruhat 胞腔中，由此可得计算所需的坐标。

\subsubsection{}\label{exI:211Sp4} 验证 \S\ref{sss:Sp4} 中的论断：设 $G=\Sp(V)$，$V=\jiao{v_{1},v_{2},v_{3},v_{4}}$，辛形式由 $\jiao{v_{i},v_{5-i}}=1$（$i=1,2$）与 $\jiao{v_{i},v_{j}}=0$（$i+j\neq 5$）给出。设 $e: v_{4}\mapsto v_{1}\to 0, v_{2}\mapsto 0,v_{3}\mapsto 0$。则 $\sB_{e}$ 中任一旗 $0\subset V_{1}\subset V_{2}\subset V_{1}^{\bot}\subset V$ 必满足
\begin{equation*}
\jiao{v_{1}}\subset V_{2}\subset\jiao{v_{1},v_{2},v_{3}}.
\end{equation*}
而且这也是旗 $0\subset V_{1}\subset V_{2}\subset V_{1}^{\bot}\subset V$ 属于 $\sB_{e}$ 的唯一条件。

\subsubsection{}\label{exI:inv} 设 $e\in\sN$，$B\subset G$ 为 Borel 子群。
\begin{enumerate}
\item 设 $\alpha$ 为单根。设 $P_{\alpha}\supset B$ 是抛物子群，其 Levi 因子的半单秩为 $1$、根为 $\pm\alpha$。设 $\sP_{\alpha}\cong G/P_{\alpha}$ 是 $G$ 的部分旗簇，分类与 $P_{\alpha}$ 共轭的抛物子群。把投影 $\sB\to \sP_{\alpha}$ 限制到 $\sB_{e}$，得到映射
\begin{equation*}
\pi_{\alpha}: \sB_{e}\to \pi_{\alpha}(\sB_{e}).
\end{equation*}
证明上同调上的拉回 $\pi^{*}_{\alpha}$ 诱导同构
\begin{equation}\label{sa inv}
\cohog{*}{\pi_{\alpha}(\sB_{e})}\cong \cohog{*}{\sB_{e}}^{s_{\alpha}}
\end{equation}
其中 $s_{\alpha}\in W$ 是与 $\alpha$ 相联系的单反射，它经由 Springer 作用作用在 $\cohog{*}{\sB_{e}}$ 上。
\item 能否把 \eqref{sa inv} 推广到其它部分旗簇？
\end{enumerate}

\subsubsection{}\label{exI:SL3} 设 $G=\SL_{3}$，$e\in\sN$ 为次正则元。用两个不可约分支的基本类给出的基（见 \S\ref{sss:SL3 subreg}），计算 $S_{3}$ 在二维空间 $\cohog{2}{\sB_{e}}$ 上的作用。

{\em 提示}：对本题与下一题，习题 \ref{exI:inv} 会很有用。

\subsubsection{}\label{exI:Sp4} 描述 $\Sp_{4}$ 的 Springer 纤维，并具体算出 $G=\Sp_{4}$ 的 Springer 对应。

\subsubsection{} 用 $\sB_{e}$ 的维数公式验证 Springer 消解 $\pi$ 是半小的。

\subsubsection{}\label{exI:dim BX} 设 $X\in\frg$，$X=X_{s}+X_{n}$ 为 $X$ 的 Jordan 分解。设 $H=C_{G}(X_{s})\subset G$，这是 $G$ 的 Levi 子群。设 $\sB^{H}$ 为 $H$ 的旗簇，$\sB^{H}_{X_{n}}$ 为把 $X_{n}$ 看作 $\Lie H$ 中的幂零元时的 Springer 纤维。证明 $\dim\tB_{X}=\dim \sB^{H}_{X_{n}}$。

\subsubsection{} 用习题 \ref{exI:dim BX} 与维数公式 \eqref{Spr dim}，导出一切 $X\in\frg$ 的 $\tB_{X}$ 的维数公式。再用你的公式证明 Grothendieck 修改 $\pi_{\frg}$ 是小映射。

\subsubsection{} 证明 $\sN$ 有理光滑，即其相交上同调复形同构于移位常层 $\Ql[\dim\sN]$。

{\em 提示}：$\bR\pi_{*}\Ql$ 中最大的直和项就是 $\sN$ 的移位 $\IC$ 层，它也是 $\bR\pi_{\frg,*}\Ql$ 的某个直和项的限制。

\subsubsection{}\label{exI:conn} 证明 Springer 纤维 $\sB_{e}$ 是连通的。

{\em 提示}：各 Springer 纤维的 $\upH^{0}$ 被装在某一个层里。

\subsubsection{} 记 $G$ 的单根为 $\{\alpha_{1},\cdots,\alpha_{r}\}$。若抛物子群 $P\subset G$ 的 Levi 商 $L_{P}$ 的根为 $\pm\alpha_{i}$，则称 $P$ 为 $i$ 型。
\begin{enumerate}
\item 设 $1\leq i\leq r$，$P$ 为 $i$ 型抛物子群。设 $\frn_{P}$ 为 $\Lie P$ 的幂零根。证明 $\frn_{P}\cap \calO_{\subreg}$ 在 $\frn_{P}$ 中稠密。
\item 设 $e\in\calO_{\subreg}$。证明对每个 $i$，只有有限多个 $i$ 型抛物子群 $P$ 满足 $e\in\frn_{P}$。对每个这样的 $P$，$\sB_{e}$ 的子簇 $C_{P}:=\{B\in\sB_{e}|B\subset P\}$ 同构于 $\PP^{1}$，称为 $i$ 型曲线。
\item 对 $i$ 型抛物子群 $P$ 与 $j$ 型抛物子群 $Q$（$P\neq Q$），证明 $C_{P}\cap C_{Q}$ 或为空集，或为单点。
\item 设 $P$ 为 $i$ 型抛物子群且 $e\in\frn_{P}$。则对任意 $1\leq j\leq r$、$i\neq j$，$C_{P}$ 恰与 $-\jiao{\alpha^{\vee}_{i}, \alpha_{j}}$ 条 $j$ 型曲线相交。
\item 证明 $\sB_{e}$ 具有 \S\ref{sss:subreg} 中描述的构型。
\item 用习题 \ref{exI:inv} 计算 $W$ 在 $\cohog{2}{\sB_{e}}$ 上的 Springer 作用。
\end{enumerate}

\section{第二讲：仿射 Springer 纤维}\label{s:ASF}

仿射 Springer 纤维是环路群的 Springer 纤维的类比物，由 Kazhdan 与 Lusztig \cite{KL} 引入。粗略地说，对经典群而言，Springer 纤维分类的是 $k$ 向量空间中被 $k$ 线性变换固定的旗；仿射 Springer 纤维分类的则是 $F$ 向量空间中被 $F$ 线性变换固定的格（及其链），其中 $F=k\lr{t}$。仿射 Springer 纤维的上同调群带有仿射 Weyl 群的作用。

本节的设定与 \S\ref{ss:notation} 相同。

\subsection{环路群、parahoric 子群与仿射旗簇}
设 $F=k\lr{t}$ 为单变量 $t$ 的形式 Laurent 级数域。$F$ 有离散赋值 $\val_{F}:F^{\times}\to\ZZ$，满足 $\val_{F}(t)=1$，其赋值环为 $\calO_{F}=k\tl{t}$。

\subsubsection{环路群} {\em 正环路群} $L^{+}G$ 是 $k$ 代数上的群值函子。对任意 $k$ 代数 $R$，定义 $LG(R):=G(R\tl{t})$。可以证明 $LG$ 可由 $k$ 上的一个概形表示，但该概形不是有限型的。\zy{原文此处称 $L^{+}G$ 为「正环路群」，紧接的定义却是 $LG(R):=G(R\tl{t})$，即环路群。下一段才定义正环路群 $L^{+}G$。原文的这句应属笔误，译文照录，并登记为勘误 E1。［来源：原文 \texttt{ZhiweiYunPCMI.tex} 第 1251--1253 行文脉。］}

例如当 $G=\GL_{n}$ 时，$LG(R)=\GL_{n}(R\tl{t})$ 中的元素由 $n^{2}$ 个形式 Laurent 级数 $a_{ij}=\sum_{s\geq 0}a^{(s)}_{ij}t^{s}$（$1\leq i,j\leq n$）给出，其中 $a^{(s)}_{ij}\in R$，并满足一个条件：$\det((a_{ij}))$（它是 $a^{(s)}_{ij}$ 的 $n$ 次多项式）在 $R$ 中可逆。因此此时 $LG$ 是以 $a^{(s)}_{ij}$（$1\leq i,j\leq n$，$s\geq 0$）为坐标的无穷维仿射空间中的开子概形。

类似地可以定义{\em 环路群} $LG$ 为 $k$ 代数 $R$ 上的函子 $LG(R)=G(R\lr{t})$。函子 $LG$ 不再可由概形表示，而是由{\em 归纳概形}表示。归纳概形是范畴意义下的归纳极限 $\varinjlim_{m} X_{m}$：即 $\{X_{m}\}$ 构成 $k$ 上概形的归纳系，而 $\varinjlim_{m} X_{m}$ 是函子 $R\mapsto \varinjlim_{m} X_{m}(R)$。当 $G=\GL_{n}$ 时，可令 $X_{m}$ 为 $LG$ 的子函子，使 $X_{m}(R)$ 由元素落在 $t^{-m}R\tl{t}\subset R\lr{t}$ 中的 $n$ 阶可逆矩阵组成。此时与 $L^{+}G$ 的情形同样的论证表明 $X_{m}$ 可由 $k$ 上概形表示。当 $m<m'$ 时，有自然的闭嵌入 $X_{m}\incl X_{m'}$，而此时的 $LG$ 就是归纳极限 $\varinjlim_{m}X_{m}$。对一般的 $G$，参见 \cite[\S1]{BL} 与 \cite[\S2]{Faltings}。

\subsubsection{仿射 Grassmann 簇}\label{sss:Gr} $G$ 的仿射 Grassmann 簇 $\Gr_{G}$ 定义为预层 $R\mapsto LG(R)/L^{+}G(R)$ 在 $k$ 代数范畴中关于 fpqc 拓扑的层化。特别地，有 $\Gr_{G}(k)=G(F)/G(\calO_{F})$。

当 $G=\GL_{n}$ 时，仿射 Grassmann 簇 $\Gr_{G}$ 可以等同于射影 $R\tl{t}$-子模 $\L\subset R\lr{t}^{n}$ 的模空间，其中对某个 $m\geq 0$ 有
\begin{equation}\label{lattice range}
(t^{m}R\tl{t})^{n}\subset \L\subset (t^{-m}R\tl{t})^{n}.
\end{equation}
这样的 $R\tl{t}$-模 $\L$ 称为 $R\lr{t}^{n}$ 中的{\em 格}。对固定的 $m$，令 $X_{m}$ 为 $\Gr_{G}$ 中分类满足 \eqref{lattice range} 的 $\L$ 的子函子，则 $X_{m}$ 可由 $k$ 上射影概形表示。自然的闭嵌入 $X_{m}\incl X_{m+1}$ 使 $\{X_{m}\}$ 成为射影概形的归纳系，而 $\Gr_{G}$ 可由归纳概形 $\varinjlim_{m}X_{m}$ 表示。

我们来详述 $\Gr_{G}(k)$ 与 $F^{n}=k\lr{t}^{n}$ 中的格之间的双射。设 $\calO_{F}^{n}\subset F^{n}$ 为标准格。设 $\frL_{n}$ 为 $F^{n}$ 中格的集合（当 $R=k$ 时，格就是 $F^{n}$ 的秩 $n$ 的 $\calO_{F}$-子模）。群 $LG(k)=G(F)$ 通过 $LG\ni g:\L\mapsto g\L$ 作用在 $\frL_{n}$ 上。此作用可迁，且标准格 $\calO^{n}_{F}$ 的稳定化子是 $L^{+}G(k)=G(\calO_{F})$。因此该作用诱导出 $G(F)$ 等变双射
\begin{equation}\label{Gr lattices}
\Gr_{G}(k)=G(F)/G(\calO_{F})\isom \frL_{n}.
\end{equation}

对一般的 $G$，$\Gr_{G}$ 总可由形如 $\varinjlim_{m}X_{m}$ 的归纳概形表示，其中 $X_{m}$ 是 $k$ 上射影概形，转移映射 $X_{m}\incl X_{m+1}$ 是闭嵌入。仿射 Schubert 分层给出 $\Gr_{G}$ 被射影概形的典范穷尽，现回忆如下。$L^{+}G$ 通过左平移作用在 $\Gr_{G}$ 上，其轨道以支配余权 $\l\in\xcoch(T)^{+}$ 为指标。用 $\Gr_{G,\l}$ 记过 $\l(t)$ 的 $L^{+}G$ 轨道，用 $\Gr_{G,\leq \l}$ 记 $\Gr_{G,\l}$ 的闭包。则 $\Gr_{G,\leq\l}$ 是射影概形，且 $\Gr_{G}$ 是这些 $\Gr_{G,\leq\l}$ 之并。关于仿射 Grassmann 簇的更多细节，参见 \cite[\S2]{BL}、\cite[\S2]{Faltings} 与 Zhu 的讲义 \cite{Zhu}。

\subsubsection{parahoric 子群} 子群 $L^{+}G\subset LG$ 是一类称为{\em parahoric 子群}的子群的例子。固定 Borel 子群 $B\subset G$，令 $\bI\subset L^{+}G$ 为 $B$ 在「模 $t$ 归约」映射 $L^{+}G\to G$ 下的原像。则 $\bI$ 是 $LG$ 的{\em Iwahori 子群}的一例。一般的 Iwahori 子群是 $\bI$ 在 $LG$ 中的共轭。与 $L^{+}G$ 一样，Iwahori 子群是 $LG$ 的无穷型群子概形。parahoric 子群是 $LG$ 中包含某个 Iwahori 子群、余维有限的连通群子概形。parahoric 子群的精确定义需要相当多的 Bruhat--Tits 理论，参见 Bruhat 与 Tits 的原始论文 \cite{BT1} 以及综述文章 \cite{Tits}。

正如 $G$ 的抛物子群的共轭类与 $G$ 的 Dynkin 图的子集一一对应，$LG$ 的 parahoric 子群在 $LG$ 共轭下的类，与 $G$ 的{\em 扩张 Dynkin 图} $\wt{\Dyn}(G)$（比 $G$ 的 Dynkin 图多一个顶点）的顶点的{\em 真子集}一一对应。扩张 Dynkin 图参见 Kac 的书 \cite[\S4.8]{Kac} 与 Bourbaki \cite[Ch VI]{Bourbaki}；与 parahoric 子群的联系参见 Gross 的综述文章 \cite{Gross}。

每个 $\bP$ 都有典范的群概形正合列
\begin{equation*}
1\to \bP^{+}\to \bP\to L_{\bP}\to1
\end{equation*}
其中 $\bP^{+}$ 是 $\bP$ 的{\em 顺幂幺根}，$L_{\bP}$ 是 $k$ 上的约化群，称为 $\bP$ 的{\em Levi 商}。若 $\bP$ 对应 $\wt{\Dyn}(G)$ 的顶点子集 $J$，则 Levi 商 $L_{\bP}$ 的 Dynkin 图是 $\wt{\Dyn}(G)$ 中由 $J$ 张成的子图。

\subsubsection{仿射旗簇} 对每个 parahoric 子群 $\bP\subset LG$，可以定义相应的{\em 仿射部分旗簇} $\Fl_{\bP}$，它是函子 $R\mapsto LG(R)/\bP(R)$ 在 $k$ 代数范畴中关于 fpqc 拓扑的层化。该函子同样可由归纳概形 $\varinjlim_{m}X_{m}$ 表示，其中每个 $X_{m}$ 是 $k$ 上射影概形，转移映射为闭嵌入。当 $\bP=L^{+}G$ 时，仿射 Grassmann 簇 $\Gr_{G}$ 是 $\Fl_{\bP}$ 的特例。

考虑特例：$\bP=\bI$ 是 $LG$ 的 Iwahori 子群。当 $G$ 单连通时，$\bI$ 是其自身的正规化子，于是可以把 $\Fl_{\bI}$ 等同于 $LG$ 的 Iwahori 子群的模空间，从而得到仿射旗簇的内蕴定义。通常把 $\Fl_{\bI}$ 记为 $\Fl$ 或 $\Fl_{G}$，称为{\em $G$ 的仿射旗簇}；但要提醒的是，只有当 $G$ 单连通时，$\Fl_{\bI}$ 才典范地不依赖于 $\bI$ 的选取。

设 $\bP\subset\bQ$ 是 $LG$ 的两个 parahoric 子群。则有自然投影 $\Fl_{\bP}\to \Fl_{\bQ}$。此投影的纤维同构于 $L_{\bQ}$ 的部分旗簇，对应的抛物子群由 $\bP\to L_{\bQ}$ 的像给出。特别地，有自然投影 $\Fl_{G}\to \Gr_{G}$，其纤维同构于旗簇 $\sB$。

\subsubsection{$\SL_{n}$ 的情形}\label{sss:GLn parahoric} 在 \S\ref{sss:Gr} 中已经看到，$\GL_{n}$ 的仿射 Grassmann 簇可以解释为格的模空间。事实上，$LG$ 的 parahoric 子群及相应的仿射部分旗簇也可以用格来描述。这里考虑 $G=\SL_{n}$ 的情形。

回忆 $F^{n}$ 中格的集合记为 $\frL_{n}$。对任意两个格 $\L_{1},\L_{2}\in\frL_{n}$，定义它们的相对长度为整数
\begin{equation*}
[\L_{1}:\L_{2}]:=\dim_{k}(\L_{1}/\L_{1}\cap\L_{2})-\dim_{k}(\L_{2}/\L_{1}\cap\L_{2}).
\end{equation*}
设 $J\subset \ZZ/n\ZZ$ 为非空子集。设 $\tilJ$ 是 $J$ 在投影 $\ZZ\to\ZZ/n\ZZ$ 下的原像。{\em 周期的 $J$-格链}指函数 $\L:\tilJ\to \frL_{n}$，把每个 $i\in \tilJ$ 送到格 $\L_{i}\in\frL_{n}$，并满足
\begin{itemize}
\item $[\L_{i}:\calO_{F}^{n}]=i$，对一切 $i\in \tilJ$；
\item $\L_{i}\subset \L_{j}$，对 $\tilJ$ 中的 $i<j$；
\item $\L_{i}=t\L_{i+n}$，对一切 $i\in \tilJ$。
\end{itemize}
设 $\frL_{J}$ 为周期的 $J$-格链的集合。对每个 $\{\L_{i}\}_{i\in \tilJ}\in\frL_{J}$，令 $\bP_{\{\L_{i}\}_{i\in J}}\subset LG$ 为所有 $\L_{i}$ 的同时稳定化子。则 $\bP_{\{\L_{i}\}_{i\in J}}$ 是 $LG$ 的 parahoric 子群，称之为 $J$ 型 parahoric 子群。事实上 $LG$ 的所有 parahoric 子群都来自唯一的一条周期 $J$-格链，其中 $J\subset \ZZ/n\ZZ$ 是非空子集且唯一。于是得到 $\sqcup_{J}\frL_{J}$ 与 $LG$ 的 parahoric 子群集合之间的双射。特别地，$L^{+}G$ 对应由 $\L_{i}=t^{i/n}\calO^{n}_{F}$（$i\in \tilJ=n\ZZ$）给出的周期 $\{0\}$-格链。

$G$ 的扩张 Dynkin 图是含 $n$ 个节点的圈，我们用 $\ZZ/n\ZZ$ 循环地为节点编号，使 $0$ 对应比通常 Dynkin 图多出的那个节点。$J\neq\varnothing$ 型的 parahoric 子群对应扩张 Dynkin 图节点的真子集 $\ZZ/n\ZZ-J$。

可以找到周期 $J$-格链的模空间 $\Fl_{J}$，使 $\Fl_{J}(k)=\frL_{J}$。固定任意 $J$ 型 parahoric 子群 $\bP$，$\Fl_{J}$ 可等同于仿射部分旗簇 $\Fl_{\bP}$。特别地，当 $G=\SL_{n}$ 时，仿射旗簇 $\Fl$ 可以等同于{\em 周期完备格链}的模空间，即 $F^{n}$ 中满足 $[\L_{i}:\calO^{n}_{F}]=i$ 与 $\L_{i}=t\L_{i+n}$（对一切 $i\in\ZZ$）的格序列 $\cdots \L_{-1}\subset \L_{0}\subset \L_{1}\cdots$。

\subsubsection{$\Sp_{2n}$ 的情形}\label{sss:Sp parahoric} 现考虑 $G=\Sp_{2n}=\Sp(V)$，其中 $V=k^{2n}$ 配有辛形式。把 $V$ 上的辛形式 $F$ 线性地扩张为 $V\otimes_{k}F$ 上的辛形式 $\jiao{-,-}$。对格 $\L\in \frL_{2n}$，定义其辛对偶为集合 $\L^{\vee}:=\{v\in V\otimes_{k}F|\jiao{v,\L}\subset \calO_{F}\}$。则 $\L^{\vee}$ 仍是 $V\otimes_{k}F$ 中的格。运算 $\L\mapsto \L^{\vee}$ 定义了 $\frL_{2n}$ 上的对合。

设 $J\subset\ZZ/2n\ZZ$ 是非空且在乘 $-1$ 下稳定的子集。设 $\tilJ\subset \ZZ$ 是 $J$ 在自然投影 $\ZZ\to\ZZ/2n\ZZ$ 下的原像。$V\otimes_{k}F$ 中的{\em 周期自对偶 $J$-格链}指满足下述附加条件的周期 $J$-格链（即按 \S\ref{sss:GLn parahoric} 的记号属于 $\frL_{J}$ 的元素）：
\begin{equation*}
\L^{\vee}_{i}=\L_{-i}, \text{ 对一切 } i\in \tilJ.
\end{equation*}
把 $V\otimes_{k}F$ 中周期自对偶 $J$-格链的集合记为 $\frL^{\Sp(V)}_{J}$。这是带 $G(F)=\Sp(V\otimes_{k}F)$ 作用的集合。对任意 $\{\L_{i}\}_{i\in\tilJ}\in \frL^{\Sp(V)}_{J}$，诸 $\L_{i}$ 的同时稳定化子是 $LG$ 的 parahoric 子群，而且 $LG$ 的每个 parahoric 子群都这样得到。对 $J$ 型 parahoric 子群 $\bP$，相应的仿射部分旗簇 $\Fl_{\bP}$ 可以等同于周期自对偶 $J$-格链的模空间，使 $\Fl_{\bP}(k)\cong\frL^{\Sp(V)}_{J}$ 作为 $G(F)$-集。请读者自行写出正交群的类似叙述，见习题 \ref{exII:SO parahoric}。

\subsection{仿射 Springer 纤维}

\subsubsection{仿射 Grassmann 簇中的仿射 Springer 纤维} 对任意 $k$ 代数 $R$，把 $\frg\otimes_{k}R$ 记为 $\frg(R)$。特别地，$\frg(F)=\frg\otimes_{k}F$ 是圈群 $LG$ 的 Lie 代数。对 $g\in LG$，用 $\Ad(g)$ 记它在 $\frg(F)$ 上的伴随作用。

设 $\gamma\in\frg(F):=\frg\otimes_{k}F$ 是正则半单元\footnote{这里涉及的 Lie 代数是域 $F$（它不是代数闭域）上的。称 $\gamma\in\frg(F)$ 为正则半单元，指它作为 $\frg(\overline{F})$ 中的元素正则半单，参见 \S\ref{sss:Gr alt} 的脚注。等价地说，$\gamma$ 正则半单当且仅当它在 $\frc(F)$ 中的像落在 $\frc^{\rs}(F)$ 中。}。考虑 $\Gr_{G}$ 的子函子：它在 $k$ 代数 $R$ 上的取值由
\begin{equation}\label{aff Spr}
\tX_{\gamma}(R)=\{[g]\in \Gr_{G}(R)|\Ad(g^{-1})\gamma\in\frg(R\tl{t})\}
\end{equation}
给出。则 $\tX_{\gamma}$ 是 $\Gr_{G}$ 的闭子归纳概形。令 $\sX_{\gamma}=\tX^{\red}_{\gamma}$ 为 $\tX_{\gamma}$ 底层的约化归纳概形。称 $\sX_{\gamma}$ 为 $\gamma$ 在仿射 Grassmann 簇 $\Gr_{G}$ 中的{\em 仿射 Springer 纤维}。

\subsubsection{用格给出的另一套定义}\label{sss:ASF Gr GLn} 考虑 $G=\GL_{n}$ 的情形。设 $\gamma\in\frg(F)=\gl_{n}(F)$ 是正则半单矩阵。按 \S\ref{sss:Gr} 的做法，把 $\Gr_{G}$ 等同于 $F^{n}$ 中格的模空间；更精确地说，$\Gr_{G}(R)$ 是 $R\lr{t}^{n}$ 中格的集合。此时 $\tX_{\gamma}(R)$ 可以等同于那些满足 $\gamma \L\subset \L$ 的格 $\L\subset R\lr{t}^{n}$，即那些在 $R\lr{t}^{n}$ 的自同态 $\gamma$ 下稳定的格。

当 $G=\SL_{n}$ 时，$\Gr_{G}(R)$ 分类 $R\lr{t}^{n}$ 中满足 $[\L:R\tl{t}^{n}]=0$ 的格 $\L$。此时的仿射 Springer 纤维 $\tX_{\gamma}$ 仍由同样的条件 $\gamma \L\subset \L$ 刻画。

当 $G=\Sp_{2n}$ 时，$\Gr_{G}(R)$ 分类 $R\lr{t}^{2n}$ 中满足 $\L^{\vee}=\L$ 的格 $\L$，参见 \S\ref{sss:Sp parahoric}。此时的仿射 Springer 纤维 $\tX_{\gamma}$ 也由同样的条件 $\gamma \L\subset \L$ 刻画。

下面给出仿射 Springer 纤维最简单的几个例子。

\subsubsection{}\label{sss:reduction rs} 设 $\gamma\in\frt(\calO_{F})$，且其约化 $\ov\gamma\in\frt$ 正则半单。对每个余权 $\l:\Gm\to T$，元素 $t^{\l}:=\l(t)\in T(F)$ 给出点 $[t^{\l}]\in \Gr_{G}(k)$；由于 $\Ad(t^{-\l})\gamma=\gamma\in \frg(\calO_{F})$，该点落在 $\tX_{\gamma}$ 中。约化归纳概形 $\sX_{\gamma}$ 实际上就是离散集 $\{[t^{\l}]\}$，它与 $\xcoch(T)$ 一一对应。更典范地说，圈群 $LT$ 通过它在 $\Gr_{G}$ 上的左平移作用而作用在 $\tX_{\gamma}$ 上。此作用经由商 $\Gr_{T}=LT/L^{+}T$ 分解，并把 $\tX_{\gamma}$ 实现为一个 $\Gr_{T}$-主齐性空间。

\subsubsection{}\label{P1 chain} 考虑 $G=\SL_{2}$ 与 $\gamma=\mat{t}{0}{0}{-t}$ 的情形。此时 $\sX_{\gamma}$ 是 $\PP^{1}$ 的无穷长链。更精确地说，对每个 $n\in\ZZ$，考虑 $\Gr_{G}$ 的子概形 $C_{n}$，它分类满足 $t^{n}\calO_{F}\oplus t^{-n+1}\calO_{F} \subset \L\subset t^{n-1}\calO_{F}\oplus t^{-n}\calO_{F}$ 的格 $\L\subset F^{2}$。则 $C_{n}\cong\PP^{1}$。于是 $\sX_{\gamma}=\cup_{n\in\ZZ}C_{n}$ 是 $\PP^{1}$ 的无穷长链：分支 $C_{n}$ 与 $C_{n+1}$ 交于一点 $t^{-n}\calO_{F}\oplus t^{n}\calO_{F}$，此外诸分支互不相交。

下面给出一种计算 $\sX_{\gamma}$ 的 $k$ 点的方法。使用 $G(F)$ 的{\em Iwasawa 分解}：
\begin{equation*}
G(F)=\bigsqcup_{n\in\ZZ}N(F)\mat{t^{n}}{0}{0}{t^{-n}}G(\calO_{F}).
\end{equation*}
按此分解，$\Gr_{G}$ 中任一点都可由
\begin{equation}\label{Iwasawa elt}
g=\mat{1}{x}{0}{1}\mat{t^{n}}{0}{0}{t^{-n}}
\end{equation}
表示，其中 $x\in F$、$n\in\ZZ$ 唯一，且 $x$ 在 $F/t^{2n}\calO_{F}$ 中有确定的像。由于
\begin{equation*}
\Ad(g^{-1})\gamma=g^{-1}\gamma g=\mat{t}{2t^{1-2n}x}{0}{-t}
\end{equation*}
条件 $\Ad(g^{-1})\gamma\in\frg(\calO_{F})$ 等价于要求 $x\in t^{2n-1}\calO_{F}$。因此 $\sX_{\gamma}(k)=\sqcup_{n\in\ZZ}Y_{n}$，其中 $Y_{n}$ 由形如 \eqref{Iwasawa elt} 且 $x\in t^{2n-1}\calO_{F}/t^{2n}\calO_{F}$ 的元素组成。于是每个 $Y_{n}$ 都可等同于 $k$。容易验证 $Y_{n}\subset C_{n}(k)$。

\subsubsection{}\label{SL2 ram dim 1} 考虑 $G=\SL_{2}$ 与 $\gamma=\mat{0}{t^{2}}{t}{0}$ 的情形。此时 $\sX_{\gamma}$ 恰由那些满足 $t\calO_{F}\oplus \calO_{F}\subset\L\subset\calO_{F}\oplus t^{-1}\calO_{F}$ 的格 $\L\in\Gr_{G}$ 组成。因此 $\sX_{\gamma}\cong\PP^{1}$。这些计算的细节留给读者，见习题 \ref{exII:SL2}。

\subsubsection{共轭不变性} 若 $\gamma,\gamma'\in\frg(F)$ 满足 $\Ad(g)\gamma=\gamma'$（$g\in G(F)$），则 $\Gr_{G}$ 上的左乘 $g$ 限制为同构 $\tX_{\gamma}\cong\tX_{\gamma'}$，从而也有 $\sX_{\gamma}\cong\sX_{\gamma'}$。因此 $\sX_{\gamma}$ 的同构类型在 $\gamma$ 的 $G(F)$-共轭下不变。回忆有映射 $\chi:\frg\to \frc:=\frg\sslash G\cong\frt\sslash W$。对正则半单点 $a\in\frc^{\rs}(F)$，纤维 $\chi^{-1}(a)$ 是单个 $G(F)$-共轭类（这里用到 $F$ 的剩余域代数闭）。因此 $\sX_{\gamma}$ 的同构类型只依赖于 $a=\chi(\gamma)\in\frc^{\rs}(F)$。

与 Springer 纤维不同，$\sX_{\gamma}$ 对某些 $\gamma$ 可以是空集。仿射 Springer 纤维 $\sX_{\gamma}$ 非空，当且仅当 $a=\chi(\gamma)\in\frc(\calO_{F})$。事实上，若 $gG(\calO_{F})\in\sX_{\gamma}(k)$，则 $\Ad(g^{-1})\gamma\in\frg(\calO_{F})$，故 $\chi(\gamma)=\chi(\Ad(g^{-1})\gamma)\in\frc(\calO_{F})$。反之，$\chi$ 有 Kostant 截面 $\ep: \frc\to \frg$，它把 $\frc$ 等同于 $e+\frg_{f}$，其中 $(e,h,f)$ 是 $\frg$ 中的正则 $\sl_{2}$ 三元组。因此对任意 $a\in\frc(\calO_{F})\cap\frc^{\rs}(F)$，有 $\ep(a)\in\frg(\calO_{F})$，于是 $\sX_{\ep(a)}$ 包含 $\Gr_{G}$ 中的单位陪集，从而非空；再由 $\sX_{\gamma}$ 与 $\sX_{\ep(\chi(\gamma))}$ 同构，$\sX_{\gamma}$ 也非空。

对 $a\in\frc(\calO_{F})\cap\frc^{\rs}(F)$，也把 $\sX_{\ep(a)}$ 记为 $\sX_{a}$。上面的讨论表明，所有 Springer 纤维 $\sX_{\gamma}$ 都同构于 $\sX_{a}$，其中 $a=\chi(\gamma)$。

\subsubsection{parahoric 版本} 对每个 parahoric 子群 $\bP\subset LG$，可以用 \eqref{aff Spr} 的条件的类似物（把 $\frg(R\tl{t})$ 换成 $(\Lie\bP)\widehat{\otimes}_{k}R$）定义闭子归纳概形 $\tX_{\bP,\gamma}\subset\Fl_{\bP}$。约化归纳概形 $\sX_{\bP,\gamma}=\tX_{\bP,\gamma}^{\red}\subset\Fl_{\bP}$ 称为 $\gamma$ 的{\em $\bP$ 型仿射 Springer 纤维}。特别地，当 $\bI$ 是 $LG$ 的 Iwahori 子群时，把 $\sX_{\bI,\gamma}$ 记为 $\sY_{\gamma}$。

设 $\bP\subset\bQ$ 是 $LG$ 的两个 parahoric 子群。自然投影 $\Fl_{\bP}\to \Fl_{\bQ}$ 诱导映射 $\sX_{\bP,\gamma}\to \sX_{\bQ,\gamma}$。特别地，有映射 $\sY_{\gamma}\to \sX_{\gamma}$。

\subsection{仿射 Springer 纤维上的对称性}
对 Springer 纤维 $\sB_{e}$，中心化子 $C_{G}(e)$ 作用在它上面。本节研究仿射 Springer 纤维上的类似结构。

\subsubsection{中心化子的作用}\label{sss:centralizer} 设 $G_{\gamma}$ 是 $\gamma$ 在 $G_{F}$（$G$ 经基变换得到的 $F$ 上代数群）中的中心化子。则 $G_{\gamma}$ 是 $F$ 上代数群。由于 $\gamma$ 正则半单，$G_{\gamma}$ 是 $F$ 上的环面。可以定义 $G_{\gamma}$ 的圈群 $LG_{\gamma}$ 为 $k$ 代数上的函子 $R\mapsto G_{\gamma}(R\lr{t})$。

我们断言 $LG_{\gamma}$ 作用在归纳概形 $\tX_{\gamma}$ 上。这可以在 $R$ 点上看出。设 $h\in LG_{\gamma}(R)=G_{\gamma}(R\lr{t})$ 与 $[g]\in \tX_{\gamma}(R)$。则陪集 $[hg]\in \Gr_{G}(R)$ 仍满足
\begin{equation*}
\Ad((hg)^{-1})\gamma=\Ad(g^{-1})\Ad(h^{-1})\gamma=\Ad(g^{-1})\gamma\in\frg(R\tl{t}),
\end{equation*}
这里用到 $h$ 在 $\gamma$ 的中心化子中。因此 $[hg]\in \tX_{\gamma}(R)$。映射 $[g]\mapsto [hg]$（$h\in LG_{\gamma}$，$[g]\in\tX_{\gamma}$）定义了 $LG_{\gamma}$ 在 $\tX_{\gamma}$ 上的作用。它诱导出 $LG_{\gamma}$ 在约化结构 $\sX_{\gamma}$ 上的作用。

\subsubsection{分裂情形} 考虑 $\gamma\in\frt^{\rs}(F)$ 的情形。此时 $G_{\gamma}=T\otimes_{k}F$，且
\begin{equation*}
LG_{\gamma}=LT\cong\xcoch(T)\otimes_{\ZZ}L\Gm,
\end{equation*}
其中 $L\Gm$ 是乘法群 $\Gm$ 的圈群。对任意 $k$ 代数 $R$，$L\Gm(R)=R\lr{t}^{\times}$。容易看出，元素 $a=\sum_{i}a_{i}t^{i}\in R\lr{t}$ 可逆，当且仅当 $a$ 以有限多个幂零系数开头、且第一个非幂零系数在 $R$ 中可逆。当 $R$ 约化时，$a$ 的首项系数必在 $R$ 中可逆，于是 $R\lr{t}^{\times}=t^{\ZZ}\cdot R\tl{t}^{\times}$，而 $R\tl{t}^{\times}=\{(a_{0},a_{1},\cdots)|a_{0}\in R^{\times}, a_{i}\in R, \forall i\geq 1\}$。由此看出约化归纳概形 $(L\Gm)^{\red}\cong \ZZ\times L^{+}\Gm$，且作为概形 $L^{+}\Gm\cong \Gm\times\AA^{\NN}$，其中 $\AA^{\NN}=\Spec k[x_{1},x_{2},\cdots]$。因此当 $\gamma\in\frt^{\rs}(F)$ 时，$(LT)^{\red}\cong\xcoch(T)\times L^{+}T$，而 $L^{+}T$ 是无穷型的仿射概形。

\subsubsection{格 $\L_{\gamma}$}\label{sss:L gamma} 对一般的正则半单 $\gamma\in\frg(F)$，记 $\xcoch(G_{\gamma}):=\Hom_{F}(\Gm,G_{\gamma})$ 为环面 $G_{\gamma}$ 的 $F$-有理余权格。每个 $\l\in\xcoch(G_{\gamma})$ 可视为定义在 $F$ 上的同态 $\Gm\to G_{\gamma}$，于是可以谈论元素 $\l(t)$。对应 $\l\mapsto \l(t)$ 定义了单同态
\begin{equation*}
\xcoch(G_{\gamma})\incl G_{\gamma}(F),
\end{equation*}
其像记为 $\L_{\gamma}$。可以证明商 $\L_{\gamma}\bs (LG_{\gamma})^{\red}$ 是仿射概形，并且是若干 $\Gm^{a}\times \AA^{\NN}$（$a$ 为某整数）的有限不交并。

\subsubsection{$G=\GL_{n}$ 的情形}\label{sss:GLn centralizer} 沿用 \S\ref{sss:ASF Gr GLn} 的设定。假设 $\chark>n$。则 $\gamma$ 的特征多项式 $P(x)=x^{n}+a_{1}x^{n-1}+\cdots+a_{n}\in F[x]$ 可分。于是 $F$ 代数 $F[x]/(P(x))$ 是域之积 $F_{1}\times \cdots \times F_{m}$，且 $\sum_{i=1}^{m}[F_{i}:F]=n$。每个域扩张 $F_{i}/F$ 由添加 $P(x)$ 的不可约因子 $P_{i}(x)$ 的一个根得到，而由 $\chark>n$ 知 $F_{i}$ 必形如 $k\lr{t^{1/e_{i}}}$。此时中心化子 $G_{\gamma}$ 同构于{\em 诱导环面}之积：
\begin{equation*}
G_{\gamma}\cong \prod_{i=1}^{m}\Res^{F_{i}}_{F}\Gm.
\end{equation*}
有 $\xcoch(G_{\gamma})\cong \ZZ^{m}$，而映射 $\xcoch(G_{\gamma})\to G_{\gamma}(F)$ 由
\begin{equation*}
\ZZ^{m}\ni(d_{1},\cdots, d_{m})\mapsto (t^{d_{1}},\cdots, t^{d_{m}})\in F^{\times}_{1}\times\cdots\times F^{\times}_{m}
\end{equation*}
给出。商 $\L_{\gamma}\bs G_{\gamma}(F)$ 同构于 $\prod_{i=1}^{m}F^{\times}_{i}/t^{\ZZ}$。由于每个 $F_{i}$ 都同构于 $k\lr{t^{1/e_{i}}}$，有正合列 $1\to \calO^{\times}_{F_{i}}\to F^{\times}_{i}/t^{\ZZ}\to \ZZ/e_{i}\ZZ\to 0$，从而商 $\L_{\gamma}\bs (LG_{\gamma})^{\red}$ 以有限指数包含群概形 $\prod_{i=1}^{m}L^{+}_{F_{i}}\Gm$。这里 $L^{+}_{F_{i}}\Gm$ 作为概形同构于 $L^{+}\Gm$，只是把素元改名为 $t^{1/e_{i}}$。

另一种做法是：对每个 $i$ 固定素元 $t_{i}\in F_{i}$（例如取 $t_{i}=t^{1/e_{i}}$），并令 $\wt\L_{\gamma}=t_{1}^{\ZZ}\times\cdots\times t^{\ZZ}_{m}\subset \prod_{i}F_{i}^{\times}=G_{\gamma}(F)$。格 $\wt\L_{\gamma}$ 在计算轨道积分时有用，参见 \S\ref{sss:GLn orb coho}。

\subsubsection{$G=\SL_{2}$ 的情形} 设 $G=\SL_{2}$ 与 $\gamma=\mat{0}{t^{n}}{1}{0}$，其中 $n\geq 1$ 为奇数。则 $G_{\gamma}(F)$ 由满足 $a^{2}-t^{n}b^{2}=1$ 的矩阵 $\mat{a}{bt^{n}}{b}{a}$（$a,b\in F$）组成。注意此方程迫使 $a,b\in\calO_{F}$，故 $G_{\gamma}(F)=G_{\gamma}(\calO_{F})$。环面 $G_{\gamma}$ 非分裂，在二次扩张 $E=F(t^{1/2})$ 上分裂。此时格 $\L_{\gamma}=\Hom_{F}(\Gm,G_{\gamma})=0$。把 $a$ 与 $b$ 分别写成 $a=\sum_{i\geq0}a_{i}t^{i}$ 与 $b_{i}=\sum_{i\geq 0}b_{i}t^{i}$\zy{原文此处写作「$b_{i}=\sum_{i\geq0}b_{i}t^{i}$」，下标 $i$ 显系误置，应为「$b=\sum_{i\geq0}b_{i}t^{i}$」。译文照录，并登记为勘误 E2。［来源：原文 \texttt{ZhiweiYunPCMI.tex} 第 1431 行；下文「一旦 $b$ 与 $a_{0}$ 固定」亦要求以 $b$ 为对象。］}，可以看到 $a_{0}=\pm1$，并且一旦 $b$ 与 $a_{0}$ 固定，$a$ 的更高次系数就可用 $(1+t^{n}b^{2})^{1/2}$ 的 Taylor 展开唯一地解出。因此 $(LG_{\gamma})^{\red}\cong L^{+}G_{\gamma}$ 同构于 $\{\pm1\}\times \AA^{\NN}$，同构由 $(a,b)\mapsto (a_{0},b_{0},b_{1},\cdots)$ 给出。

\subsubsection{仿射 Springer 纤维上的对称性}\label{sss:sym ASF} 关于 $LG_{\gamma}$ 在 $\tX_{\gamma}$ 上的作用，Ng\^o 找到了更精确的结论，即该作用经由一个典范的有限维商分解。下面按 \cite[\S3.3]{NgoFL} 概述此事。设 $a=\chi(\gamma)\in\frc(F)^{\rs}$ 是 $\gamma$ 在 $\chi:\frg\to \frc$ 下的像。假定 $a\in\frc(\calO_{F})$，否则 $\sX_{\gamma}$ 为空。

存在 $\frc$ 上的光滑仿射群概形 $J$，称为{\em 正则中心化子群概形}。它的刻画性质是：它经由 $\chi$ 拉回到 $\frg$ 所得的 $\chi^{*}J$ 映入 $\frg$ 上的万有中心化子群概形 $I$，且该映射在正则轨迹 $\frg^{\reg}$ 上是同构。令 $J_{a}$ 为 $J$ 在映射 $a:\Spec\calO_{F}\to \frc$ 下的拉回。则 $J_{a}$ 是 $\calO_{F}$ 上的光滑仿射群概形，其 $F$-纤维是环面 $G_{\gamma}$（即 $J_{a}$ 是 $G_{\gamma}$ 在 $\calO_{F}$ 上的整模型）。可以构造 $J_{a}$ 的正圈群 $L^{+}J_{a}$，以及它的仿射 Grassmann 簇 $P_{a}:=LG_{\gamma}/L^{+}J_{a}$（也称为{\em 局部 Picard 群}）。约化群概形 $P^{\red}_{a}$ 是有限维的，且局部有限型。Ng\^o 证明：$LG_{\gamma}$ 在 $\tX_{\gamma}$（从而在 $\sX_{\gamma}$）上的作用经由局部 Picard 群 $P_{a}$ 分解，且不能经由更小的商分解。相关陈述见 \S\ref{dim formula discussion}。

\subsection{仿射 Springer 纤维的更多例子}\label{ss:aff examples}
本小节再给出几个例子，展示仿射 Springer 纤维的丰富几何。导至这些几何描述的计算一概从略。

下面的所有例子中，仿射 Springer 纤维都是{\em 齐性}的：$\sX_{\gamma}$ 上多了一个 $\Gm$-作用，它来自 $\Gr_{G}$ 上的{\em 圈旋转}，即在素元 $t$ 上作伸缩。关于齐性仿射 Springer 纤维及其在表示论中的应用，参见 \cite{OY}。

\subsubsection{}\label{SL2 ram} 设 $G=\SL_{2}$ 与 $\gamma=\mat{0}{t^{m+1}}{t^{m}}{0}$。当 $m$ 为偶数时，$\sX_{\gamma}$ 恰由那些满足
\begin{equation*}
t^{m/2}\calO_{F}\oplus t^{m/2}\calO_{F}\subset\L\subset t^{-m/2}\calO_{F}\oplus t^{-m/2}\calO_{F}
\end{equation*}
的格 $\L$ 组成。此时 $\sX_{\gamma}$ 与 $\Gr_{G}$ 中对应于余权 $m\alpha^{\vee}/2$ 的 $L^{+}G$-轨道的闭包重合。

当 $m$ 为奇数时，$\sX_{\gamma}$ 恰由那些满足
\begin{equation*}
t^{(m+1)/2}\calO_{F}\oplus t^{(m-1)/2}\calO_{F}\subset\L\subset t^{-(m-1)/2}\calO_{F}\oplus t^{-(m+1)/2}\calO_{F}
\end{equation*}
的格 $\L$ 组成。此时改为考虑 $G'=\PGL_{2}$ 的仿射 Grassmann 簇 $\Gr_{G'}$，它把 $\Gr_{G}$ 作为其一个分支包含在内。$\sX_{\gamma}$ 可以等同于 $\Gr_{G'}$ 中对应于余权 $m\alpha^{\vee}/2$ 的 $L^{+}G'$-轨道的闭包。

无论哪种情形，都有 $\dim\sX_{\gamma}=m$。

\subsubsection{Lusztig--Smelt 例}\label{sss:LS} 设 $G=\GL_{n}$，$\gamma\in\frg(F)$ 的特征多项式为 $P(x)=x^{n}-t^{m}=0$，其中 $(m,n)=1$。若格 $\L\subset F^{n}$ 在 $\gamma$ 下稳定，则它带有环 $R=\calO_{F}[X]/(X^{n}-t^{m})$ 的作用。令 $K=\Frac(R)$，则 $K= k\lr{s}$，其中 $x=s^{m}$、$t=s^{n}$。于是 $R$ 在 $K$ 中的整闭包是 $\tilR:=k\tl{s}$。$\gamma$ 在 $F^{n}$ 上的作用使之成为一维 $K$-向量空间。固定一个 $K$-线性同构 $F^{n}\cong K$，在此同构下，$F^{n}$ 中在 $R$ 下稳定的格就只是{\em 分式 $R$-理想}，即 $K=\Frac(R)$ 的有限生成 $R$-子模。于是 $\sX_{\gamma}$ 可以等同于分式 $R$-理想的模空间。

中心化子 $G_{\gamma}(F)$ 就是 $K^{\times}$，它通过乘法作用在分式 $R$-理想集合上。此作用显然经由 $K^{\times}/R^{\times}$ 分解，而后者正是局部 Picard 群概形 $P_{\gamma}$ 的 $k$ 点群。

存在 $\Gm$ 在 $K$ 上的作用（由域自同构给出），它在 $s$ 上作伸缩（于是 $s^{i}$ 在此作用下有权 $i$）。此作用诱导 $\Gm$ 在 $\sX_{\gamma}$ 上的作用。$\Gm$ 在 $\sX_{\gamma}$ 上的不动点对应由 $s$ 的单项式生成的分式理想。更精确地说，若分式 $R$-理想 $\L\subset K=k\lr{s}$ 被 $\Gm$ 固定，则定义 $M_{\L}=\{i\in\ZZ|s^{i}\in \L\}$；它是 $\ZZ$ 的子集，且在加 $m$ 与加 $n$ 下稳定，因为 $\L$ 是 $R=k\tl{s^{m},s^{n}}$-模。于是 $M_{\L}\subset\ZZ$ 是幺半群 $A_{m,n}:=\ZZ_{\geq 0}m+\ZZ_{\geq0}n\subset\ZZ_{\geq 0}$ 上的有限生成模。对应 $\L\mapsto M_{\L}$ 给出双射
\begin{equation*}
\sX_{\gamma}^{\Gm}(k)\isom\{\mbox{$A_{m,n}$-子模 $M\subset \ZZ$}\}.
\end{equation*}
$\ZZ$ 的任一 $A_{m,n}$-子模都包含所有充分大的整数。因此任意两个这样的 $A_{m,n}$-模 $M$ 与 $M'$ 只差有限多个元素，从而可以定义 $[M:M']=\#(M\bs M')-\#(M'\bs M)$。固定任意 $i\in\ZZ$，满足 $[M:\ZZ_{\geq0}]=i$ 的不动点共有 $\frac{1}{n+m}\binom{n+m}{n}$ 个。对与 $A_{m,n}$-模 $M$ 对应的不动点 $p_{M}$，考虑子簇 $C_{M}=\{p\in\sX_{\gamma}|\lim_{\Gm\ni z\to0}z\cdot p=p_{M}\}$。则 $C_{M}$ 同构于仿射空间，其维数可按 $M$ 组合地表达。诸胞腔 $C_{M}$ 给出 $\sX_{\gamma}$ 的分层。由此可以计算 $\sX_{\gamma}$ 各连通分支的 Poincar\'e 多项式。更多细节以及 $\sY_{\gamma}$ 的类似图景，参见 \cite{LS}。

\subsubsection{} 我们在 \S\ref{sss:LS} 的一个特例中考察 $\sY_{\gamma}$ 的几何。取 $G=\GL_{3}$ 与 $\gamma^{3}-t^{2}=0$。如前引入变量 $s$，使 $t=s^{3}$、$\gamma=s^{2}$，则 $R=k\tl{s^{2},s^{3}}$，其分式域为 $K=k\lr{s}$。仿射 Springer 纤维 $\sY_{\gamma}$ 分类分式 $R$-理想的链 $\L_{0}\subset \L_{1}\subset \L_{2}\subset s^{-3}\L_{0}$。考虑分支 $\sY^{0}_{\gamma}\subset\sY_{\gamma}$，它分类如上并满足 $[\L_{i}:k\tl{s}]=i$（$0\leq i\leq 2$）的链 $\{\L_{i}\}$。

先研究 $\sY^{0}_{\gamma}$ 上的 $\Gm$-不动点。对每个被 $\Gm$ 固定的 $R$-子模 $k\lr{s}$，用一列整数记其相应的幺半群 $A_{3,2}$-模（见 \S\ref{sss:LS}）。例如 $(0,1,2,\cdots)$ 表示标准格 $k\tl{s}$。注意任一被 $\Gm$ 固定的 $R$-分式理想对应的数列，要么是连续的 $(n,n+1,n+2,\cdots)$，要么在第二处至多有一个缺口 $(n,n+2,n+3,\cdots)$。$\sY^{0}_{\gamma}$ 中有如下四个不动点：
\begin{enumerate}
\item $q: (0,1,2,\cdots)\subset (-1,0,1,\cdots)\subset (-2,-1,0,\cdots)$；
\item $p_{0}: (-1,1,2,\cdots)\subset (-1,0,1,\cdots)\subset (-2,-1,0,\cdots)$；
\item $p_{1}: (0,1,2,\cdots)\subset (-2,0,1,\cdots)\subset (-2,-1,0,\cdots)$；
\item $p_{2}:(0,1,2,\cdots)\subset (-1,0,1,\cdots)\subset (-3,-1,0,\cdots)$。
\end{enumerate}
我们的编号方式是：$p_{i}$ 由 $q$ 改变第 $i$ 个格 $\L_{i}$ 得到。

于是 $\sY^{0}_{\gamma}$ 是三个不可约分支之并 $C_{0}\cup C_{1}\cup C_{2}$，每个分支 $C_{i}$ 同构于 $\PP^{1}$。三者都过 $q$，且 $q$ 是任意两个分支唯一的交点。对 $i=0,1,2$ 有 $p_{i}\in C_{i}$。

如 \S\ref{sss:GLn parahoric} 所见，$G$ 的仿射部分旗簇可以自然地用子集 $J\subset \{0,1,2\}$ 编号。令 $\sX_{J,\gamma}$ 为 $\gamma$ 在 $\Fl_{J}$ 中的仿射 Springer 纤维。在投影 $\sY_{\gamma}\to \sX_{J,\gamma}$ 下，$i\notin J$ 的曲线 $C_{i}$ 缩成一点，其余的曲线同构地映到其像上。

\subsubsection{}\label{ex:Sp} 设 $G=\Sp(V)$，其中 $V$ 是 $k$ 上维数为 $2n$ 的辛空间，并假定 $\chark\neq2$。固定分解 $V=U\oplus U^{*}$ 为 $V$ 的 Lagrange 子空间，使辛形式限制在 $U\times U^{*}$ 上为自然配对。考虑 $\gamma=\mat{0}{tX}{Y}{0}\in\frg(F)$，其中 $X\in\Sym^{2}(U)$（视为自伴随映射 $\xi:U^{*}\to U$），$Y\in\Sym^{2}(U^{*})$（视为自伴随映射 $\eta:U\to U^{*}$）。$\gamma$ 正则半单的条件等价于：(1) $\xi$ 与 $\eta$ 都是同构；(2) $\xi\eta\in\GL(U)$ 正则半单，或等价地 $\eta\xi\in\GL(U^{*})$ 正则半单。

仿射 Springer 纤维 $\sX_{\gamma}$ 分类在 $\gamma$ 下稳定的自对偶格 $\L\subset V\otimes_{k}F$（见 \S\ref{sss:ASF Gr GLn}）。考虑 $\Gm$ 在 $V\otimes_{k}F=U\otimes_{k}F\oplus U^{*}\otimes_{k}F$ 上的作用，使 $Ut^{i}$ 的权为 $2i-1$、$U^{*}t^{i}$ 的权为 $2i$。此作用诱导 $\Gm$ 在 $\Gr_{G}$ 与 $\sX_{\gamma}$ 上的作用。先考虑不动点 $\Gr_{G}^{\Gm}$。格 $\L\in\Gr_{G}$ 被 $\Gm$ 固定，当且仅当它是
\begin{equation*}
\Span\{(A_{i}\oplus B_{i})t^{i}; i\in\ZZ \}
\end{equation*}
的 $t$-进完备化，其中 $\cdots \subset A_{-1}\subset A_{0}\subset A_{1}\subset \cdots \subset U$ 是 $U$ 的滤过，满足 $N\ll 0$ 时 $A_{N}=0$、$N\gg 0$ 时 $A_{N}=U$；而 $\cdots \subset B_{-1}\subset B_{0}\subset B_{1}\subset \cdots \subset U^{*}$ 是 $U^{*}$ 的类似滤过，并且在对偶配对下成立
\begin{equation}\label{AB dual}
B_{i-1}=A^{\bot}_{-i}, \forall i\in\ZZ.
\end{equation}
最后一个条件反映了 $\L$ 在辛形式下自对偶这一事实。

格 $\L\in\sX_{\gamma}^{\Gm}$ 于是由 $U$ 的滤过 $A_{\bu}$ 与 $U^{*}$ 的滤过 $B_{\bu}$ 决定，两者在 \eqref{AB dual} 的意义下对偶，并满足附加条件
\begin{equation*}
\eta A_{i}\subset B_{i}; \quad \xi B_{i}\subset A_{i+1}, \quad\forall i\in\ZZ.
\end{equation*}
把上述数据概括为下图
\begin{equation*}
\xymatrix{ A_{-1}\ar@{^{(}->}[rr]\ar[dr]^{\eta} && A_{0}\ar@{^{(}->}[rr]\ar[dr]^{\eta} && A_{1}\ar[dr]^{\eta}& \cdots\\
\cdots & A^{\bot}_{0}\ar@{^{(}->}[rr]\ar[ur]^{\xi} && A^{\bot}_{-1}\ar@{^{(}->}[rr]\ar[ur]^{\xi} && A_{-2}^{\bot}}
\end{equation*}

例如当 $\dim V=4$（从而 $\dim U=2$）时有两种可能。第一种是 $A_{-1}=0\subset A_{0}=U$，它对应标准格 $V\otimes_{k}\calO_{F}$。第二种是 $A_{-1}=0\subset A_{0}\subset A_{1}=U$，其中 $A_{0}\subset U$ 是一条满足 $\xi A^{\bot}_{0}\subset A_{0}$ 的直线，即 $A^{\bot}_{0}\subset U^{*}$ 是二次形式 $X$ 下的迷向直线。这样的直线 $A^{\bot}_{0}$ 有两条，给出另外两个 $\Gm$-不动点。因此 $\sX_{\gamma}^{\Gm}$ 由 3 个点组成。此时 $\sX_{\gamma}=C_{1}\cup C_{2}$，其中 $C_{i}\cong\PP^{1}$，两个分支交于标准格，且每个 $C_{i}$ 含有其余 $\Gm$-不动点之一。

\subsubsection{Bernstein--Kazhdan 例}\label{sss:BK example} 在 \cite[附录]{KL} 中，Bernstein 与 Kazhdan 给出了仿射 Springer 纤维的不可约分支不是有理簇的第一个例子。沿用 \S\ref{ex:Sp} 的记号。设 $\Fl_{\bP}$ 是 $G=\Sp(V)$ 的仿射部分旗簇，它分类满足 $\L^{\vee}=\L$ 与 $\L'^{\vee}=t^{-1}\L$ 的格对 $\L'\subset\L$。

取 \S\ref{ex:Sp} 中的 $\gamma$。同一个 $\Gm$ 作用在 $\sX_{\bP,\gamma}$ 上，而不动点 $\L'\subset \L$ 可以用两对滤过 $(A_{\bu}, B_{\bu})$ 与 $(A'_{\bu}, B'_{\bu})$ 来描述：其中 $(A_{\bu}, B_{\bu})$ 是 \S\ref{ex:Sp} 所述那种 $U$ 与 $U^{*}$ 的滤过，$(A'_{\bu}, B'_{\bu})$ 类似，只是把 \eqref{AB dual} 换成 $B'_{i}=A_{-i}'^{\bot}$。此外，包含 $\L'\subset \L$ 等价于对一切 $i$ 有 $A'_{i}\subset A_{i}$ 与 $B'_{i}\subset B_{i}$。

例如取 $\dim V=6$，并固定滤过的维数：
\begin{equation*}
A_{-1}=0, \quad \dim A_{0}=2, \quad A_{1}=U;\quad  \dim A'_{0}=1, \quad A'_{1}=U.
\end{equation*}
这样的滤过组 $(A_{\bu}, B_{\bu}; A'_{\bu}, B'_{\bu})$ 由满足 $\xi A^{\bot}_{0}\subset A_{0}$ 与 $\eta A'_{0}\subset A'^{\bot}_{0}$ 的完备旗 $0\subset A'_{0}\subset A_{0}\subset U$ 决定。换言之，$A^{\bot}_{0}$ 是 $U^{*}$ 中二次形式 $X$ 下的迷向直线，$A'_{0}$ 是 $U$ 中二次形式 $Y$ 下的迷向直线。数对 $(A^{\bot}_{0}, A'_{0})$ 决定 $Q(X)\times Q(Y)\subset \PP(U)\times\PP(U^{*})$ 中一点，这里 $Q(X)$ 与 $Q(Y)$ 分别是由 $X$ 与 $Y$ 定义的二次曲线。关联关系 $A'_{0}\subset A_{0}$ 在 $Q(X)\times Q(Y)\cong\PP^{1}\times \PP^{1}$ 中定义了一条双次数为 $(2,2)$ 的曲线，从而是一条亏格为 $1$ 的曲线。因此 $\sX^{0}_{\bP,\gamma}$ 的连通分支是一条亏格 $1$ 的曲线。考虑 $\sX_{\bP,\gamma}$ 中收缩到该曲线的点，并取其闭包 $Z$。可以证明 $\dim Z=3=\dim\sX_{\bP,\gamma}$。于是 $\sX_{\bP,\gamma}$ 含有一个非有理的不可约分支 $Z$。更多细节参见 \cite{KL} 的附录。

\subsubsection{「次正则」仿射 Springer 纤维}
当 $\sX_{\gamma}$ 或 $\sY_{\gamma}$ 是一维时，不妨称之为次正则仿射 Springer 纤维，以类比 \S\ref{sss:subreg} 中讨论的次正则 Springer 纤维。若 $\dim\sY_{\gamma}=1$，则它是 $\PP^{1}$ 之并，于是可以定义它的对偶图。在 \cite[命题 7.7]{KL} 中，$\Fl$ 中次正则仿射 Springer 纤维的对偶图得到了分类：它们几乎总是单连型群的扩张 Dynkin 图，只是在 $A$ 型时还可以是无穷长链（见例 \ref{P1 chain}）。

\subsection{仿射 Springer 纤维的几何性质}

\subsubsection{非约化性} 若 $G$ 非平凡、$\gamma\in\frg(F)$ 正则半单，则归纳概形 $\tX_{\gamma}$ {\em 从不}约化。例如在 \S\ref{sss:reduction rs} 考虑的情形中，$\tX_{\gamma}$ 同构于 $\Gr_{T}=LT/L^{+}T$。在 \S\ref{sss:centralizer} 中已经看到，对非约化的环 $R$，$LT(R)=R\lr{t}^{\times}$ 中的元素可以有幂零的首项系数。因此 $\Gr_{T}(R)$ 并不只是 $\xcoch(T)$，后者才是 $\Gr^{\red}_{T}$。这说明 $\Gr_{T}$ 非约化，从而 $\tX_{\gamma}$ 非约化。

下面这个定理是关于 $\sX_{\gamma}$ 的基本有限性陈述。

\begin{theorem}[实质上是 Kazhdan 与 Lusztig {\cite[命题 2.1]{KL}}]\label{th:ft quot} 设 $\gamma\in \frg(F)$ 是正则半单元。则
\begin{enumerate}
\item 存在闭子概形 $Z\subset \sX_{\gamma}$，它在 $k$ 上射影，且 $\sX_{\gamma}=\cup_{\ell\in\L_{\gamma}}\ell\cdot Z$。
\item 归纳概形 $\sX_{\gamma}$ 是 $k$ 上局部有限型的概形。
\item $\L_{\gamma}$ 在 $\sX_{\gamma}$ 上的作用自由，且商 $\L_{\gamma}\bs\sX_{\gamma}$（作为 $k$ 代数上的 fppf 层）可由 $k$ 上真代数空间表示。
\end{enumerate}
\end{theorem}

下面分三步概述该定理的证明。

\subsubsection{第一步归约} 我们说明定理的第 (1) 部分蕴含 (2) 与 (3)。设 $Z$ 是如 (1) 中的射影子概形。为证 (2)，只需说明任一点 $x\in \sX_{\gamma}(k)$ 都有有限型的开邻域。借助 $\L_{\gamma}$-作用，可以假定 $x\in Z(k)$。由于 $Z$ 有限型，集合 $\Sigma:=\{\ell\in \L_{\gamma}|Z\cap\ell\cdot Z\neq\varnothing\}$ 有限。令 $U=\sX_{\gamma}-\cup_{\ell\notin\Sigma}\ell\cdot Z$，则 $U$ 是 $Z$ 的开邻域，从而也是 $x$ 的开邻域。此外 $U$ 包含在有限并 $\cup_{\ell\in\Sigma}\ell\cdot Z$ 中，从而包含在某个 Schubert 簇 $\Gr_{G,\leq\l}$ 中。因此 $U$ 是射影概形 $\sX_{\gamma}\cap\Gr_{G,\leq\l}$ 的开子集，故 $U$ 本身是有限型概形。为证 (3)，注意 fppf 层商 $\L_{\gamma}\bs\sX_{\gamma}$ 是分离代数空间，因为它是 $\sX_{\gamma}$ 对 étale 等价关系 $\L_{\gamma}\times\sX_{\gamma}\subset \sX_{\gamma}\times\sX_{\gamma}$（由作用映射与投影映射给出）的商。由 (1)，有射影概形 $Z$ 到 $\L_{\gamma}\bs\sX_{\gamma}$ 的满射，故 $\L_{\gamma}\bs\sX_{\gamma}$ 真。

\subsubsection{$\gamma$ 落在分裂环面中时 (1) 的证明}\label{sss:split case} 先考虑 $\gamma$ 落在分裂环面中的情形。经 $G(F)$-共轭，可假定 $\gamma\in\frt(F)$。此时 $\L_{\gamma}\cong\xcoch(T)$。固定含 $T$ 的 Borel 子群 $B$，令 $N$ 为 $B$ 的幂幺根。$LG$ 的 Iwasawa 分解给出
\begin{equation*}
\Gr_{G}=LN\cdot \L_{\gamma}L^{+}G/L^{+}G=\sqcup_{\l\in\xcoch(T)}LN\cdot \l(t)L^{+}G/L^{+}G.
\end{equation*}
令 $X:=\sX_{\gamma}\cap (LN\cdot L^{+}G/L^{+}G)\subset \sX_{\gamma}$。只需说明 $X$ 落在某个仿射 Schubert 簇 $\Gr_{\leq \l}$ 中，因为此时它在 $\Gr_{\leq\l}$ 中的闭包 $Z:=\ov{X}$ 便满足 (1) 的条件。为后文备用，注意平移 $\ell\cdot X$（$\ell\in \L_{\gamma}$）互不相交，且覆盖 $\sX_{\gamma}$。

固定 $T$ 关于 $B$ 的正根的一个次序，则可把元素 $u\in N(F)$ 唯一地写成
\begin{equation}\label{root gp}
u=\prod_{\alpha>0}x_{\alpha}(c_{\alpha}),
\end{equation}
其中 $c_{\alpha}\in F$，$x_{\alpha}: \Ga\to N$ 是 $\alpha$ 对应的根群。为说明 $X$ 包含在某个仿射 Schubert 簇中，只需对满足 $[u]\in X(k)$ 的任意 $u\in N(F)$，给出 \eqref{root gp} 中出现的 $c_{\alpha}$ 的赋值下界。可以把 $\Ad(u^{-1})\gamma\in\frb(F)$ 按根分解展开：
\begin{equation*}
\Ad(u^{-1})\gamma=\gamma+\sum_{\alpha>0}P_{\alpha}(\gamma; c)e_{\alpha},
\end{equation*}
其中 $e_{\alpha}\in\frg_{\alpha}$ 是每个根空间固定的基，而 $P_{\alpha}(\gamma;c)$ 是关于 $\{c_{\alpha}\}_{\alpha>0}$ 的 $F$-值多项式函数，并对 $\gamma$ 线性。对 $\alpha$ 的高度作归纳可知 $P_{\alpha}(\gamma;c)$ 具有如下形式
\begin{equation}\label{conj u gamma}
P_{\alpha}(\gamma;c)=\jiao{\alpha,\gamma}c_{\alpha}+\sum_{\beta<\alpha}\jiao{\beta,\gamma}P^{\beta}_{\alpha}(c),
\end{equation}
其中 $P^{\beta}_{\alpha}(c)$ 是只涉及 $\{c_{\alpha'}\}_{\alpha'<\alpha}$ 的多项式，且关于次数 $\alpha$ 齐次（这里约定 $\deg c_{\alpha'}:=\alpha'\in\xch(T)$）。

令 $n=\max_{\alpha>0}\{\val_{F}\jiao{\alpha,\gamma}\}$。此值有限，因为 $\gamma$ 正则半单。若 $[u]\in X(k)$，即 $\Ad(u^{-1})\gamma\in\frg(\calO_{F})$，则对 $\alpha$ 的高度作归纳可知
\begin{equation*}
\val_{F}(c_{\alpha})\geq -(2\textup{ht}(\alpha)-1)n,
\end{equation*}
这就给出了所需的下界，并说明 $X$ 落在某个仿射 Schubert 簇中。

\subsubsection{一般情形下 (1) 的证明} 一般情形下，我们给出一个比 \cite{KL} 中原本的论证更简化的论证，思路相同。作基变换到 $F'=k\lr{t^{1/m}}$，在其上 $\gamma$ 可以共轭进分裂环面。令 $\Gr'_{G}$ 为用域 $F'$ 代替 $F$ 定义的 $G$ 的仿射 Grassmann 簇（于是 $\Gr'_{G}(k)=G(F')/G(\calO'_{F})$），并令 $\sX'_{\gamma}\subset\Gr'_{G}$ 为相应的仿射 Springer 纤维。则 $\Gr'_{G}$ 与 $\sX'_{\gamma}$ 都带有 $\Gamma:=\Gal(F'/F)\cong\mu_{m}$ 的作用，该作用由 $\Gamma$ 在 $F'$ 上的作用诱导，并且有闭嵌入 $\sX_{\gamma}\incl (\sX'_{\gamma})^{\Gamma}$。

令 $\L'_{\gamma}=\xcoch(G_{\gamma}\otimes F')\incl G_{\gamma}(F')$ 为用域 $F'$ 构造的格。$\Gamma$ 作用在 $\xcoch(G_{\gamma}\otimes F')$ 上，不动点为 $\xcoch(G_{\gamma})\cong \L_{\gamma}$。此作用诱导 $\Gamma$ 在 $\L'_{\gamma}$ 上的作用，但它不保持嵌入 $\xcoch(G_{\gamma}\otimes F')\incl G_{\gamma}(F')$。不动点 $(\L'_{\gamma})^{\Gamma}=\xcoch(G_{\gamma}\otimes F')^{\Gamma}$ 未必落在 $G_{\gamma}(F)$ 中，但总以有限指数包含 $\L_{\gamma}$。

由 \S\ref{sss:split case} 中分裂情形的证明，有来自 Iwasawa 分解的有限型局部闭子概形 $X'\subset\sX'_{\gamma}$，使 $\sX'_{\gamma}$ 可以分解为{\em 不交并} $\sqcup_{\ell\in\L'_{\gamma}}\ell\cdot X'$（不是作为概形，而是作为可构造集）。把 $\ell\in \L'_{\gamma}$ 与 $\xcoch(G_{\gamma}\otimes F')$ 中的元素等同，则对 $\sigma\in\Gamma$，$\sigma(\ell)$ 有意义。容易验证：对 $\ell\in\L'_{\gamma}$ 有 $\sigma(\ell\cdot X')=\sigma(\ell)\cdot X'$，尽管 $\sigma$ 不保持嵌入 $\L'_{\gamma}\to G_{\gamma}(F')$。因此 $(\sX'_{\gamma})^{\Gamma}\subset (\L'_{\gamma})^{\Gamma}\cdot X'$。取有限陪集空间 $(\L'_{\gamma})^{\Gamma}/\L_{\gamma}$ 的代表元集 $C$，便知 $(\sX'_{\gamma})^{\Gamma}$ 包含在 $\L_{\gamma}\cdot (C\cdot X')$ 中。由于 $C\cdot X'$ 有限型，$Z'=\ov{C\cdot X'}\cap (\sX'_{\gamma})^{\Gamma}$ 是 $\sX'_{\gamma}$ 的射影子概形，其 $\L_{\gamma}$-平移覆盖 $(\sX'_{\gamma})^{\Gamma}$。最后，$\sX_{\gamma}$ 的射影子概形 $Z=Z'\cap \sX_{\gamma}$ 满足 (1) 的要求。定理 \ref{th:ft quot} 证毕。\qed

\subsubsection{归约到 Levi} \S\ref{sss:split case} 中分裂情形下定理 \ref{th:ft quot} 的证明给出了更多信息。在 Iwasawa 分解中，对 $\l\in\xcoch(T)$，令 $S_{\l}\subset \Gr_{G}$ 为 $\l(t)$ 的 $LN$-轨道。它称为{\em 半无限轨道}，因为它的维数无穷，在 $\Gr_{G}$ 中的余维数也无穷。令 $C_{\l}:=\sX_{\gamma}\cap S_{\l}$，则按 \S\ref{sss:split case} 的记号 $C_{\l}=\l(t)\cdot X$。公式 \eqref{conj u gamma} 表明：$C_{\l}\neq\varnothing$（等价地 $X\neq\varnothing$，或 $\sX_{\gamma}\neq\varnothing$）当且仅当对一切根 $\alpha$ 有 $\jiao{\alpha,\gamma}\in\calO_{F}$；而且此时 $C_{\l}$ 同构于一个维数为
\begin{equation*}
\dim C_{\l}=\sum_{\alpha>0}\val_{F}(\jiao{\alpha,\gamma})=\frac{1}{2}\val_{F}\Delta(\gamma)
\end{equation*}
的几乎仿射空间（即迭代的 $\AA^{1}$-丛）。这里 $\Delta(\gamma)$ 是 $\gamma$ 在 $\frg(F)/\frt(F)$ 上伴随作用之行列式。因此 $\sX_{\gamma}$ 可以分解为同维数的几乎仿射空间之并，以 $\xcoch(T)$ 为指标。但这个分解不是分层：$C_{\l}$ 的闭包 $\ov{C}_{\l}$ 会与其他 $C_{\l'}$ 相交，却肯定不是若干 $C_{\l'}$ 之并。

分裂情形的分解 $\sX_{\gamma}=\sqcup C_{\l}$ 有推广。设 $P$ 是 $G$ 的抛物子群，幂幺根为 $N_{P}$，Levi 子群为 $M_{P}$。令 $\mathfrak{m}_{P}=\Lie M_{P}$，并设 $\gamma\in\mathfrak{m}_{P}(F)$，且作为 $\frg(F)$ 中的元素正则半单。利用广义 Iwasawa 分解 $G(F)=N_{P}(F)M_{P}(F)G(\calO_{F})$，有定义良好的映射 $\Gr_{G}(k)\to \Gr_{M_{P}}(k)$，它把 $nm G(\calO_{F})$（$n\in N_{P}(F)$，$m\in M_{P}(F)$）送到 $m M_{P}(\calO_{F})$。然而此映射不给出归纳概形的映射。尽管如此，此映射的纤维具有无穷维仿射空间的自然结构。把此映射限制到 $\sX_{\gamma}$，得到
\begin{equation*}
\tau: \sX_{\gamma}(k)\to \sX^{M_{P}}_{\gamma}(k),
\end{equation*}
其中 $\sX^{M_{P}}_{\gamma}\subset \Gr_{M_{P}}$ 是 $\gamma$ 与群 $M_{P}$ 的仿射 Springer 纤维。$\tau$ 的纤维若非空，则是维数为 $\frac{1}{2}\val_{F}\Delta^{G}_{M_{P}}(\gamma)$ 的几乎仿射空间，其中
\begin{equation}\label{Delta G L}
\Delta^{G}_{M_{P}}(\gamma):=\det(\ad(\gamma)|\frg(F)/\mathfrak{m}_{P}(F)).
\end{equation}
设 $T\subset M_{P}$，则 $\Gr_{M_{P}}$ 的连通分支以 $\xcoch(T)/R^{\vee}_{M_{P}}$ 为指标，其中 $R^{\vee}_{M_{P}}$ 是 $M_{P}$ 的余根格。若把 $\sX^{M_{P}}_{\gamma}$ 按 $\xcoch(T)/R^{\vee}_{M_{P}}$ 分解为连通分支 $\sX^{M_{P}}_{\gamma}(\l)$ 并取它们在 $\tau$ 下的原像，就得到局部闭子归纳概形之分解 $\sX_{\gamma}=\sqcup\sX_{\gamma,\l}$，指标同样是 $\l\in\xcoch(T)/R^{\vee}_{M_{P}}$。可以证明：$\sX_{\gamma,\l}$ 若非空，则是 $\sX^{M_{P}}_{\gamma}(\l)$ 上的几乎仿射空间丛，纤维维数为 $\frac{1}{2}\val_{F}\Delta^{G}_{M_{P}}(\gamma)$。

\subsubsection{连通性与等维性}\label{sss:equidim} 当 $G$ 单连通时，$\Fl$ 连通，此时仿射 Springer 纤维 $\sY_{\gamma}$ 也连通。参见 \cite[\S4, 引理 2]{KL}。于是当 $G$ 单连通时，对一切 parahoric 子群 $\bP$，$\sX_{\bP,\gamma}$ 都连通，因为自然投影 $\sY_{\gamma}\to \sX_{\bP,\gamma}$ 满射。

\cite{KL} 中还证明了 $\sY_{\gamma}$ 等维。那里的论证类似 Spaltenstein 在 \cite{Spa77} 中对 Springer 纤维的连通性与等维性的证明。

\subsubsection{维数公式}
由定理 \ref{th:ft quot}，$\sX_{\gamma}$ 的维数有定义，且等于代数空间 $L_{\gamma}\bs\sX_{\gamma}$ 的维数。为陈述 $\dim\sX_{\gamma}$ 的公式，还需要一些记号。

考虑伴随作用 $\ad(\gamma):\frg(F)\to \frg(F)$。此映射的核是 $\frg_{\gamma}(F)$，而诱导的自同态 $\ov{\ad(\gamma)}$ 在 $\frg(F)/\frg_{\gamma}(F)$ 上可逆。令 $\Delta(\gamma)\in F^{\times}$ 为 $\ov{\ad(\gamma)}$ 的行列式。这与先前在 $\gamma$ 落在分裂环面 $\frt(F)$ 中时给出的 $\Delta(\gamma)$ 的定义一致\zy{$\Delta(\gamma)$ 与 \S4.2 中判别式 $\Delta(P_{a})$ 用同一字母而含义不同：前者是 $\ad(\gamma)$ 在 $\frg(F)/\frg_{\gamma}(F)$ 上的行列式，后者是多项式 $P_{a}$ 的判别式。来源：原文 \S2.5 该段与 \S4.2 文脉。}。

另一方面，回忆 $\xcoch(G_{\gamma})$ 是 $G_{\gamma}$ 的 $F$-有理余权群，它也是 $G_{\gamma}$ 的最大 $F$-分裂子环面的秩。令
\begin{equation*}
c(\gamma)=r-\rk_{\ZZ}\xcoch(G_{\gamma}).
\end{equation*}
则 $c(\gamma)$ 也是 $G_{\gamma}$ 的最大 $F$-各向异性子环面的秩。

\begin{theorem}[Bezrukavnikov {\cite{Bez}}，Kazhdan--Lusztig 猜想 \cite{KL}]\label{th:aff dim} 设 $\gamma\in \frg(F)$ 是正则半单元。则
\begin{equation}\label{aff Spr dim}
\dim\sX_{\gamma}=\frac{1}{2}(\val_{F}\Delta(\gamma)-c(\gamma)).
\end{equation}
\end{theorem}

\subsubsection{证明梗概}\label{dim formula discussion} 证明中起关键作用的是 $\sX_{\gamma}$ 的{\em 正则点}概念。有赋值映射 $\ev: \sX_{\gamma}\to [\frg/G]$，它把 $[g]\in\sX_{\gamma}\subset \Gr_{G}$ 送到 $\Ad(g^{-1})\gamma$ 模 $t$ 的约化，后者在 $G$ 的伴随作用下定义良好。称 $[g]\in \sX_{\gamma}$ 是{\em 正则点}，若 $\ev([g])$ 落在 $[\frg/G]$ 的开子叠 $[\frg^{\reg}/G]$ 中。令 $\sX^{\reg}_{\gamma}\subset\sX_{\gamma}$ 为 $\sX_{\gamma}$ 中由正则点构成的开子归纳概形。可以证明 $\sX^{\reg}_{\gamma}$ 非空。把 $\sX^{\reg}_{\gamma}$ 在 $\sY_{\gamma}$ 中的原像记为 $\sY^{\reg}_{\gamma}$，则投影 $\sY^{\reg}_{\gamma}\to\sX^{\reg}_{\gamma}$ 是同构。由于如 \S\ref{sss:equidim} 所述 $\sY^{\reg}_{\gamma}$ 等维，有 $\dim\sX^{\reg}_{\gamma}=\dim\sY^{\reg}_{\gamma}=\dim\sY_{\gamma}$。当然 $\dim\sX_{\gamma}\leq\dim \sY_{\gamma}$，因此必有 $\dim \sX^{\reg}_{\gamma}=\dim \sX_{\gamma}=\dim \sY_{\gamma}$。于是只需计算 $\sX^{\reg}_{\gamma}$ 的维数。

回忆在 \S\ref{sss:sym ASF} 中定义了局部 Picard 群 $P_{a}$。$P_{a}$ 在 $\sX_{\gamma}$ 上的作用保持 $\sX^{\reg}_{\gamma}$，事实上 $\sX^{\reg}_{\gamma}$ 是 $P_{a}$-主齐性空间。因此只需计算 $P_{a}=LG_{\gamma}/L^{+}J_{a}$ 的维数。

考虑投影 $\chi_{\frt}: \frt\to\frc=\frt\sslash W$。沿 $a=\chi(\gamma):\Spec\calO_{F}\to\frc$ 拉回，得到有限态射 $\chi_{a}: \chi^{-1}_{\frt}(a)=\Spec A\to \Spec\calO_{F}$，其中 $A$ 是有限平坦的 $\calO_{F}$-代数。群概形 $J_{a}$ 与 $J'_{a}=(\Res^{A}_{\calO_{F}}(T\otimes_{k}A))^{W}$ 之间有密切关系，这里 Weyl 群 $W$ 对角地作用在 $T$ 与 $A$ 上。事实上 $J_{a}$ 与 $J'_{a}$ 在特殊纤维中只差连通分支。特别地，$\dim P_{a}=\dim P'_{a}$，其中 $P'_{a}=LG_{\gamma}/L^{+}J'_{a}$。不难看出
\begin{equation*}
\dim P'_{a}=\dim (\frt\otimes \tilA)^{W}-\dim (\frt\otimes A)^{W}=\dim (\frt\otimes (\tilA/A))^{W},
\end{equation*}
其中 $\tilA$ 是 $A$ 的正规化。由此即得维数公式 \eqref{aff Spr dim}。\qed

\begin{remark} 利用仿射 Springer 纤维与 Hitchin 纤维的联系，Ng\^o \cite[推论 4.16.2]{NgoFL} 证明了 $\sX^{\reg}_{\gamma}$ 实际上在 $\sX_{\gamma}$ 中稠密。特别地，$\sX_{\gamma}$ 的每个不可约分支都是有理簇。然而对更一般的仿射部分旗簇中的仿射 Springer 纤维，这一有理性不再成立，如 \S\ref{sss:BK example} 中 Bernstein--Kazhdan 的例子所示。
\end{remark}

\subsubsection{纯性} Goresky、Kottwitz 与 MacPherson \cite{GKM} 猜想仿射 Springer 纤维的上同调应当是纯的（若 $k=\overline{\FF}_{p}$ 指 Frobenius 权的意义，若 $k=\CC$ 指 Hodge 结构的含义）。仿射 Springer 纤维的纯性将使 \cite{GKM} 的作者们能用等变上同调中的局部化技巧证明非分歧元素的基本引理。这一纯性猜想在一般情形仍悬而未决。在 \cite{GKMPurity} 中，一类称为{\em 等值}的仿射 Springer 纤维被证明是纯的。

\subsubsection{扰动下的不变性} 设 $a,a'\in\frc(\calO_{F})\cap\frc^{\rs}(F)$。若 $a$ 与 $a'$ 在映射 $\frc(\calO_{F})\to\frc(\calO_{F}/t^{N})$ 下的像相同，则称 $a\equiv a'\mod t^{N}$。\cite[命题 3.5.1]{NgoFL} 证明了：对固定的 $a\in\frc(\calO_{F})\cap\frc^{\rs}(F)$，存在 $N$（依赖于 $a$），使只要 $a'\equiv a\mod t^{N}$，就有与作用相容的同构
\begin{equation*}
\tX_{a}\cong\tX_{a'} \textup{ 与 } P_{a}\cong P_{a'}.
\end{equation*}
因此可以说：在 $\frc(\calO_{F})$ 的 $t$-进拓扑下，$\sX_{a}$ 局部常值地依赖于 $a$。

例如取 $G=\GL_{n}$，并设 $a\in\frc(\calO_{F})$ 对应特征多项式 $P(x)=x^{n}+a_{1}x^{n-1}+\cdots+a_{n}$，其根都在 $\calO_{F}$ 中且模 $t$ 互不相同。则对任意 $a'\equiv a\mod t$，$a'$ 的特征多项式也模 $t$ 有互不相同的根。此时 $\tX_{a}$ 与 $\tX_{a'}$ 都是 $P_{a}\cong P_{a'}\cong \Gr_{T}$-主齐性空间。

\subsection{仿射 Springer 表示}
本小节引入仿射情形下 Springer 的 $W$-作用在 $\cohog{*}{\sB_{e}}$ 上的类比物。

\subsubsection{仿射 Weyl 群} 把 $W$ 视为余权格 $\xcoch(T)$ 的自同构群，其中 $T$ 是 $G$ 的固定极大环面。{\em 扩张仿射 Weyl 群} $\tilW$ 是半直积
\begin{equation*}
\tilW=\xcoch(T)\rtimes W.
\end{equation*}
当 $G$ 单连通时，$\xcoch(T)$ 由余根张成，此时 $\tilW$ 是 Coxeter 群，其单反射 $\{s_{0},s_{1}, \cdots,s_{r}\}$ 与 $G$ 的扩张 Dynkin 图的节点一一对应。一般情形下，$\tilW$ 是{\em 仿射 Weyl 群} $\Wa=(\ZZ R^{\vee})\rtimes W$（其中 $\ZZ R^{\vee}$ 是余根格）与交换群 $\Om\cong \xcoch(T)/\ZZ R^{\vee}$ 的半直积；$\Wa$ 是 Coxeter 群。$\tilW$ 自然地通过仿射变换作用在仿射空间 $\xcoch(T)_{\RR}$ 上，其中 $\xcoch(T)$ 通过平移作用。

\begin{theorem}[Lusztig {\cite{L96}}，Sage {\cite{Sage}}]\label{th:aff action} 存在 $\tilW$ 在 $\homog{*}{\sY_{\gamma}}$ 上的典范作用。
\end{theorem}

由于 $\sY_{\gamma}$ 不是有限型的，$\ell$-进同调 $\homog{*}{\sY_{\gamma}}$ 应理解为归纳极限 $\varinjlim_{n}\homog{*}{\sY_{\gamma,n}}$，其中把 $\sY_{\gamma}$ 表为射影子概形之并 $\sY_{\gamma,n}$。

\subsubsection{$\tilW$-作用构造梗概} 只考虑 $G$ 单连通的情形，此时 $\tilW=\Wa$ 由仿射单反射 $s_{0}, \cdots, s_{r}$ 生成。对每个 parahoric 子群 $\bP\subset LG$，都有相应的仿射 Springer 纤维 $\tX_{\bP,\gamma}$。当 $\bP$ 含固定的 Iwahori 子群 $\bI$ 时，有投影 $\pi_{\bP,\gamma}:\tY_{\gamma}\to \tX_{\bP,\gamma}$。

设 $L_{\bP}$ 是 $\bP$ 的 Levi 商，$\frl_{\bP}=\Lie L_{\bP}$。有赋值映射 $\ev_{\bP,\gamma}: \tX_{\bP,\gamma}\to [\frl_{\bP}/L_{\bP}]$，定义如下：对满足 $\Ad(g^{-1})\gamma\in\Lie \bP$ 的 $[g]\in\Fl_{\bP}$，把陪集 $[g]=g\bP$ 送到 $\Ad(g^{-1})\gamma$ 在投影 $\Lie \bP\to\frl_{\bP}$ 下的像。此定义在 $L_{\bP}$ 的伴随作用下不变。有 Cartesian 方块
\begin{equation*}
\xymatrix{\tY_{\gamma}\ar[d]^{\pi_{\bP,\gamma}}\ar[r]^{\ev_{\bI,\gamma}} & [\wt{\frl}_{\bP}/L_{\bP}]\ar[d]^{\pi_{\frl_{\bP}}}\\
\tX_{\bP,\gamma}\ar[r]^{\ev_{\bP,\gamma}} & [\frl_{\bP}/L_{\bP}]}
\end{equation*}
其中 $\pi_{\frl_{\bP}}$ 是约化群 $\frl_{\bP}$ 的 Grothendieck 修改。由 $L_{\bP}$ 的 Springer 理论，$W(L_{\bP})$ 作用在直像复形 $\bR\pi_{\frl_{\bP,*}}\DD$ 上（其中 $\DD$ 表示 $[\wt{\frl}_{\bP}/L_{\bP}]$ 的对偶化复形）。由真基变换，得到 $W(L_{\bP})$ 在 $\bR\pi_{\bP,\gamma,*}\DD_{\tY_{\gamma}}$ 上的作用，从而得到它在 $\homog{*}{\sY_{\gamma}}$ 上的作用。取扩张 Dynkin 图第 $i$ 个节点对应的标准 parahoric 子群 $\bP$，此时 $W(L_{\bP})=\jiao{s_{i}}$，于是得到作用在 $\homog{*}{\sY_{\gamma}}$ 上的对合 $s_{i}$。为验证相邻节点 $i$ 与 $j$ 对应的 $s_{i}$ 与 $s_{j}$ 之间的辫关系，可以取标准 parahoric 子群 $\bP$ 使 $W(L_{\bP})=\jiao{s_{i},s_{j}}$，此时辫关系成立，因为 $W(L_{\bP})$ 作用在 $\homog{*}{\sY_{\gamma}}$ 上。这就说明 $\tilW$ 作用在 $\homog{*}{\sY_{\gamma}}$ 上。\qed

尽管 $\tilW$-作用的构造很简单，这些作用的计算却相当困难。这里的一个新现象是：$\tilW$ 在 $\homog{*}{\sY_{\gamma}}$ 上的作用{\em 未必半单}，下面的例子即可看到这一点。

\subsubsection{$\SL_{2}$ 中的一个例子}\label{sss:tilW SL2 ram} 考虑 $G=\SL_{2}$ 与元素 $\gamma=\mat{0}{t^{2}}{t}{0}$。这是次正则情形。仿射 Springer 纤维 $\sY_{\gamma}$ 有两个不可约分支 $C_{0}$ 与 $C_{1}$，两者都同构于 $\PP^{1}$。这里 $C_{0}$ 参数化格链 $\{\L_{i}\}$，其中 $\L_{0}$ 是标准格 $\calO^{2}_{F}$，而 $\L_{-1}\subset \L_{0}$ 变动。另一分支 $C_{1}$ 参数化格链 $\{\L_{i}\}$，其中 $\L_{-1}=t\calO_{F}\oplus\calO_{F}$，而 $\L_{0}$ 变动。基本类 $[C_{0}], [C_{1}]$ 给出最高同调群 $\homog{2}{\sY_{\gamma}}$ 的自然基。可以证明，仿射 Weyl 群 $\tilW=\jiao{s_{0},s_{1}}$ 在 $\homog{2}{\sY_{\gamma}}$ 上的作用在基 $[C_{0}]$ 与 $[C_{1}]$ 下具有如下形式：
\begin{equation*}
s_{0}=\mat{-1}{2}{0}{1}; \quad s_{1}=\mat{1}{0}{2}{-1}.
\end{equation*}
由此可见，$\homog{2}{\sY_{\gamma}}$ 是 $\tilW$ 的符号表示被平凡表示（由 $[C_{0}]+[C_{1}]$ 张成）的{\em 非平凡扩张}。还可以典范地把 $\tilW$-模 $\homog{2}{\sY_{\gamma},\ZZ}$ 等同于环路群 $LG$ 的仿射余根格。

\subsection{评注与推广}

\subsubsection{与轨道积分的关系} 如 \S\ref{s:orb} 将看到的，仿射 Springer 纤维的上同调与点数，同 $p$-进群 $G(\FF_{q}\lr{t})$ 上的轨道积分密切相关。

\subsubsection{扩张对称性} $\tilW$ 在 $\homog{*}{\sY_{\gamma}}$ 上的作用可以扩张为圈积 $\Sym(\xch(T))\rtimes \tilW$ 的作用。对{\em 齐性}仿射 Springer 纤维 $\sY_{\gamma}$（即带有来自圈旋转的环面作用的那些），等变上同调群 $\upH^{*}_{\Gm}(\sY_{\gamma})$ 带有{\em 分级双仿射 Hecke 代数}的作用，后者是 $\Sym(\xch(T))\rtimes \tilW$ 的形变。细节参见 \cite{OY}。Vasserot 与 Varagnolo \cite{Vass} \cite{VV} 在仿射 Springer 纤维的 $K$-群上构造了{\em 双仿射 Hecke 代数}的作用。

\subsubsection{群版本} 把 $\gamma$ 取在 $G(F)$ 中而非 $\frg(F)$ 中，同样可以定义仿射 Springer 纤维的群版本，仍记为 $\tX_{\gamma}$，其约化结构记为 $\sX_{\gamma}$。对 $k$ 代数 $R$，有
\begin{equation}\label{gp AFS}
\tX_{\gamma}(R)=\{[g]\in \Gr_{G}(R)|g^{-1}\gamma g\in L^{+}G(R)\}.
\end{equation}
不过在群版本中，上述定义有一个有趣的推广。回忆 $LG$ 中的 $L^{+}G$-双陪集以支配余权 $\l\in\xcoch(T)^{+}$ 为指标。对 $\l\in\xcoch(T)^{+}$，把相应的双陪集记为 $(LG)_{\l}$，它是 Schubert 分层 $\Gr_{G,\l}$ 在投影 $LG\to \Gr_{G}$ 下的原像。类似地，可以定义 $(LG)_{\leq \l}$ 为 $\Gr_{G,\l}$ 的闭包 $\Gr_{G,\leq \l}$ 的原像。把 \eqref{gp AFS} 中的条件 $g^{-1}\gamma g\in L^{+}G(R)$ 换成 $g^{-1}\gamma g\in (LG)_{\l}(R)$ 或 $g^{-1}\gamma g\in (LG)_{\leq\l}(R)$，并取约化结构，就得到约化的广义仿射 Springer 纤维 $\sX_{\l,\gamma}$ 与 $\sX_{\leq\l, \gamma}$。有开嵌入 $\sX_{\l,\gamma}\incl\sX_{\leq\l,\gamma}$，其补集是支配余权 $\mu\leq\l$ 对应的 $\sX_{\mu,\gamma}$ 之并。引入 $\sX_{\leq\l,\gamma}$ 的动机是给 $G(F)$ 上球 Hecke 函数的轨道积分以几何解释。

A. Bouthier 建立了 $\sX_{\l,\gamma}$ 的基本几何性质，与定理 \ref{th:ft quot} 和定理 \ref{th:aff dim} 平行。

\begin{theorem}[Bouthier \cite{Bouthier}] 设 $\gamma\in G(F)$ 正则半单，$\l\in\xcoch(T)^{+}$。
\begin{enumerate}
\item 广义仿射 Springer 纤维 $\sX_{\l,\gamma}$ 非空，当且仅当 $[\nu_{\gamma}]\leq\l$，其中 $[\nu_{\gamma}]\in\xcoch(T)^{+}_{\QQ}$ 是 $\gamma$ 的 Newton 点，见 \cite[\S2]{KV}。
\item 归纳概形 $\sX_{\l,\gamma}$ 局部有限型。
\item 有
\begin{equation*}
\dim\sX_{\l,\gamma}=\jiao{\rho,\l}+\frac{1}{2}(\val_{F}\Delta(\gamma)-c(\gamma)),
\end{equation*}
其中 $\rho\in\xch(T)_{\QQ}$ 是正根之和的一半，$\Delta(\gamma)$ 与 $c(\gamma)$ 的定义与 Lie 代数情形类似。
\end{enumerate}
\end{theorem}
此定理的证明用到 Vinberg 半群理论，后者是 $G$ 的一种紧化。

\subsubsection{} 在 \cite{KV} 中，Kottwitz 与 Viehmann 对 Lie 代数 $\frg(F)$ 中的元素 $\gamma$ 定义了仿射 Springer 纤维的两种推广。

\subsubsection{} 作为 Hessenberg 簇的类比，还可以考虑如下情形。设 $(\rho,V)$ 是 $k$ 上约化群 $G$ 的线性表示。设 $\L\subset V\otimes F$ 是在 $G(\calO_{F})$ 下稳定的 $\calO_{F}$-格。对 $v\in V\otimes F$，可以定义 $\Gr_{G}$ 的子归纳概形 $\tX_{v}$：
\begin{equation*}
\tX_{\L, v}(R)=\{[g]\in\Gr_{G}(R)|\rho(g^{-1})v\in\L\widehat{\otimes}_{k}R\}.
\end{equation*}
设 $\sX_{\L,v}$ 为 $\tX_{\L,v}$ 的约化结构。这些归纳概形的上同调与相对迹公式中出现的轨道积分相关。

\subsection{习题}

\subsubsection{}\label{exII:SO parahoric} 设 $G=\SO(V,q)$，其中 $V$ 是 $k$ 上配有二次形式 $q$ 的向量空间。仿照 \S\ref{sss:Sp parahoric} 的风格，用 $V\otimes_{k}F$ 中的自对偶格链来解释 $LG$ 的 parahoric 子群与仿射部分旗簇。

\subsubsection{}\label{exII:SL2} 验证 \S\ref{P1 chain} 与 \S\ref{SL2 ram dim 1} 中对 $G=\SL_{2}$ 给出的仿射 Springer 纤维的描述。

\subsubsection{} 设 $G=\SL_{2}$ 与 $\gamma=\mat{0}{t^{n}}{1}{0}\in\frg(F)$。描述 $\sY_{\gamma}$。

\subsubsection{} 设 $G=\SL_{2}$ 与 $\gamma=\mat{0}{t^{n}}{1}{0}\in\frg(F)$。在 $\sX_{\gamma}$ 上构造一个涉及圈旋转（即该作用在 $t$ 上作伸缩）的非平凡 $\Gm$-作用，并确定其不动点。

\subsubsection{} 设 $G=\SL_{2}$ 与 $\gamma=\mat{t^{n}}{0}{0}{-t^{n}}\in\frg(F)$。设 $T\subset G$ 为对角环面，则 $G_{\gamma}=T\otimes_{k}F$。正则轨迹 $\sX^{\reg}_{\gamma}$（见 \S\ref{dim formula discussion}）是什么？研究 $L^+T$ 在 $\sX_{\gamma}$ 上的轨道。

\subsubsection{} 在 \S\ref{sss:L gamma} 的设定下，证明 $\L_{\gamma}$ 在 $\Gr_{G}$ 上的作用自由，从而它在 $\sX_{\gamma}$ 上的作用也自由。再证明 $\L_{\gamma}$ 在 $\sX_{\gamma}$ 的不可约分支集合上的置换作用也自由。

\subsubsection{} 对 $G=\GL_{n}$，设 $L\subset G$ 是由块对角矩阵构成的 Levi 子群，各块大小为 $n_{1},\cdots, n_{s}$，$\sum_{i}n_{i}=n$。设 $\gamma=(\gamma_{1},\cdots,\gamma_{s})\in\frl(F)$ 作为 $\frg(F)$ 中的元素正则半单。不变量 $\Delta^{G}_{L}(\gamma)$（见 \eqref{Delta G L}）用各 $\gamma_{i}$ 的特征多项式的常见不变量如何表示？

\subsubsection{}\label{exII:SL3} 设 $G=\SL_{3}$ 与 $\gamma=\diag(x_{1}t,x_{2}t,x_{3}t)\in\frg(F)$，其中 $x_{i}\in k$ 两两不同且 $x_{1}+x_{2}+x_{3}=0$。描述仿射 Springer 纤维 $\sX_{\gamma}$ 与 $\sY_{\gamma}$。

注：若你有一整天可以消磨，这是一道好习题。

\subsubsection{} 对 $G=\Sp_{6}$ 与例 \ref{ex:Sp} 中的 $\gamma=\mat{0}{tX}{Y}{0}$，描述 $\sX_{\gamma}$ 与 $\sY_{\gamma}$ 上的 $\Gm$-不动点。

\subsubsection{} 验证 \S\ref{sss:tilW SL2 ram} 中的计算。

\subsubsection{} 设 $G=\SL_{2}$ 与 $\gamma=\mat{t}{0}{0}{-t}$。描述仿射 Springer 纤维 $\sY_{\gamma}$。在以 $\sY_{\gamma}$ 的不可约分支给出的基下，$\tilW=\jiao{s_{0},s_{1}}$ 在 $\homog{2}{\sY_{\gamma}}$ 上的作用是什么？

\section{第三讲：轨道积分}\label{s:orb}
仿射 Springer 纤维在表示论中的意义，体现在它们与轨道积分的密切关系上。轨道积分是 $p$-进群调和分析中出现的某类积分。正如有限群的共轭类对理解其表示是基本的，轨道积分对理解 $p$-进群的表示也是基本的。在某些情形下，轨道积分可以解释为仿射 Springer 纤维上的点数。

\subsection{$p$-进群上的积分}\label{ss:int}

\subsubsection{设定} 设 $F$ 是局部非阿基米德域，即 $F$ 或是 $\QQ_{p}$ 的有限扩张，或是 $\FF_{p}\lr{t}$ 的有限扩张。则 $F$ 有离散赋值 $\val: F^{\times}\to \ZZ$，我们把它归一化为满射。设 $\calO_{F}$ 是 $F$ 的赋值环，$k$ 是剩余域。因此与前面几节不同，这里的 $k$ 是有限域。假定 $\chark$ 相对所讨论的群足够大。

\subsubsection{Haar 测度与积分} 设 $G$ 是 $F$ 上代数群。拓扑群 $G(F)$ 局部紧且完全不连通。它有右不变 Haar 测度 $\mu_{G}$，在相差一个标量的意义下唯一。对可测子集 $S\subset G(F)$，用 $\vol(S,\mu_{G})$ 记它在 $\mu_{G}$ 下的体积。固定紧开子群 $K_{0}\subset G(F)$，可以把 Haar 测度 $\mu_{G}$ 归一化，使 $K_{0}$ 的体积为 $1$。例如，若取 $G$ 在 $\calO_{F}$ 上的整模型 $\calG$，则可取 $K_{0}=\calG(\calO_{F})$。

借助 Haar 测度 $\mu_{G}$，可以对 $G(F)$ 上光滑（即局部常值）、紧支、复值的函数积分。把此函数空间记为 $\sS(G(F))$（其中 $\sS$ 取自 Schwarz）。对 $f\in\sS(G(F))$，积分
\begin{equation*}
\int_{G(F)}f\mu_{G}
\end{equation*}
可以如下计算。可以找到有限指数的子群 $K\subset K_{0}$，使 $f$ 右 $K$-不变，即对一切 $g\in G(F)$、$x\in K$ 有 $f(gx)=f(g)$（见习题 \ref{exIII:K inv}）。于是上述积分化为陪集 $G(F)/K$ 上的加权计数：
\begin{equation*}
\int_{G(F)}f\mu_{G}=\vol(K,\mu_{G})\sum_{[g]\in G(F)/K}f(g)=\frac{1}{[K_{0}:K]}\sum_{[g]\in G(F)/K}f(g).
\end{equation*}
容易验证：只要 $f$ 右 $K$-不变、且 $K$ 在 $K_{0}$ 中有限指数，上式右端就与 $K$ 的选取无关。

\subsubsection{变体} 设 $H\subset G$ 是定义在 $F$ 上的子群，并带有 Haar 测度 $\mu_{H}$。考虑函数 $f\in\sS(H(F)\bs G(F))$，即 $f$ 是 $G(F)$ 上左 $H(F)$-不变、局部常值、且支集模 $H(F)$ 紧的函数，则可以定义积分
\begin{equation}\label{int quot}
\int_{H(F)\bs G(F)}f\frac{\mu_{G}}{\mu_{H}}.
\end{equation}
此积分的算法如下。同样取有限指数的子群 $K\subset G(F)$，使 $f$ 右 $K$-不变。则积分 \eqref{int quot} 可以写成双陪集 $H(F)\bs G(F)/K$ 上的加权和：
\begin{equation*}
\int_{H(F)\bs G(F)}f\frac{\mu_{G}}{\mu_{H}}=\frac{1}{[K_{0}:K]}\sum_{[g]\in H(F)\bs G(F)/K}\frac{f(g)}{\vol(H(F)\cap gKg^{-1}, \mu_{H})}.
\end{equation*}

\subsection{轨道积分}

\subsubsection{轨道积分的定义} 沿用 \S\ref{ss:int} 的设定。把 $G$ 的 Lie 代数记为 $\frg(F)$，以强调它是 $F$ 上的向量空间。设 $\varphi\in\sS(\frg(F))$，$\gamma\in\frg(F)$。考虑映射 $G(F)\to \frg(F)$，$g\mapsto \Ad(g^{-1})\gamma$。则复合 $f: g\mapsto \varphi(\Ad(g^{-1})\gamma)$ 是 $G(F)$ 上的光滑函数。于是 $f$ 局部常值，在 $\gamma$ 于 $G(F)$ 中的中心化子 $G_{\gamma}(F)$ 下左不变，且支集模 $G_{\gamma}(F)$ 紧。

固定 $G(F)$ 上的 Haar 测度 $\mu_{G}$ 与 $G_{\gamma}(F)$ 上的 Haar 测度 $\mu_{G_{\gamma}}$。下面这个积分是 \eqref{int quot} 的特例（只是我们下面把积分变量 $g\in G_{\gamma}(F)\bs G(F)$ 明确写出，而 \eqref{int quot} 中没有）：
\begin{equation*}
O_{\gamma}(\varphi):=\int_{G_{\gamma}(F)\bs G(F)}\varphi(\Ad(g^{-1})\gamma)\frac{\mu_{G}}{\mu_{G_{\gamma}}}.
\end{equation*}
这类积分称为 Lie 代数 $\frg(F)$ 上的{\em 轨道积分}\zy{式中的 $\mu_{G}/\mu_{G_{\gamma}}$ 不是两个测度相除，而是 $G_{\gamma}(F)\bs G(F)$ 上由 $\mu_{G}$ 与 $\mu_{G_{\gamma}}$ 定出的商测度的记号，满足 $\int_{G}f\,\mu_{G}=\int_{G_{\gamma}\bs G}\big(\int_{G_{\gamma}}f(hg)\,\mu_{G_{\gamma}}\big)\,\mathrm{d}(\mu_{G}/\mu_{G_{\gamma}})$。来源：原文 \S3.1 的 \eqref{int quot} 与此处文脉。}。把 $\gamma$ 换成 $G(F)$ 中的元素、$\varphi$ 换成 $\sS(G(F))$ 中的元素，同样可以定义群 $G(F)$ 上的轨道积分。

\subsubsection{具体情形} 本节余下部分限制在如下情形。设 $G$ 是 $F$ 上分裂约化群。把相应的 Chevalley 群概形从 $\ZZ$ 基变换到 $\calO_{F}$，就固定了 $G$ 的一个整模型。以下把 $G$ 视为 $\calO_{F}$ 上的约化群概形。要求 $K_{0}=G(\calO_{F})$ 的体积为 $1$，以此归一化 $G(F)$ 上的 Haar 测度 $\mu_{G}$。

Lie 代数 $\frg(F)$ 含有来自 $\calO_{F}$ 上整模型的典范格 $\frg(\calO_{F})$。我们最感兴趣的是格 $\frg(\calO_{F})$ 的特征函数 $\varphi=\one_{\frg(\calO_{F})}$ 的轨道积分。

\subsubsection{$\gamma$ 的中心化子} 设 $\gamma$ 正则半单，于是其中心化子 $G_{\gamma}$ 是 $F$ 上环面。令 $F^{\ur}=F\wh{\otimes}_{k}\kbar$，它是剩余域代数闭的完备离散赋值域。仍用 $t$ 记 $F$ 的素元，它也是 $F^{\ur}$ 的素元。用 $t$，\S\ref{sss:L gamma} 中的构造给出嵌入 $\xcoch(G_{\gamma}\otimes_{F}F^{\ur})\incl  G_{\gamma}(F^{\ur})$，其像仍记为 $\L_{\gamma}$。此嵌入是 $\Gk$-等变的，故 $\L_{\gamma}$ 带有 $\Gk$ 的作用。把 $\L_{\gamma}$ 看成 $k$ 上的 étale 群概形，则记号 $\L_{\gamma}(k)$ 有意义；若把 $\L_{\gamma}$ 看成普通群，它就恰是 $\L_{\gamma}$ 中的 $\Gk$-不变量。于是 $\L_{\gamma}(k)\subset G_{\gamma}(F)$ 是离散余紧子群。

\subsubsection{$\GL_{n}$ 中的中心化子}\label{sss:GLn orb} 设 $G=\GL_{n}$，$\gamma$ 是 $\frg(F)$ 中的正则半单元，不要求它在 $F$ 上可对角化。假定 $\textup{char}(F)=0$ 或 $\chark>n$。如 \S\ref{sss:GLn centralizer}，$\gamma$ 的特征多项式 $P(x)=x^{n}+a_{1}x^{n-1}+\cdots+a_{n}\in F[x]$ 可分，故 $F$ 代数 $F[x]/(P(x))$ 同构于域之积 $F_{1}\times \cdots \times F_{m}$。有
\begin{equation*}
G_{\gamma}(F)\cong F^{\times}_{1}\times\cdots\times F^{\times}_{m}.
\end{equation*}
此时格 $\L_{\gamma}(k)\subset G_{\gamma}(F)=F^{\times}_{1}\times\cdots\times F^{\times}_{m}$ 由形如 $(t^{d_{1}},\cdots, t^{d_{m}})$（$(d_{1},\cdots, d_{m})\in\ZZ^{m}$）的元素组成。商 $\L_{\gamma}(k)\bs G_{\gamma}(F)$ 同构于 $\prod_{i=1}^{m}F^{\times}_{i}/t^{\ZZ}$。每个 $F^{\times}_{i}/t^{\ZZ}$ 落在正合列
\begin{equation}\label{Fi seq}
0\to \calO^{\times}_{F_{i}}\to F^{\times}_{i}/t^{\ZZ}\to \ZZ/e_{i}\ZZ\to 0
\end{equation}
中，其中 $e_{i}$ 是扩张 $F_{i}/F$ 的分歧指数，因此商 $\L_{\gamma}(k)\bs G_{\gamma}(F)$ 紧。

\subsubsection{用计数表示轨道积分}\label{sss:X gamma} 考虑 $G(F)/G(\calO_{F})$ 的如下子集
\begin{equation*}
X_{\gamma}:=\{[g]\in G(F)/G(\calO_{F})|\Ad(g^{-1})\gamma\in\frg(\calO_{F})\}.
\end{equation*}
这是仿射 Springer 纤维的集合论版本。

群 $G_{\gamma}(F)$ 按规则 $G_{\gamma}(F)\ni h:[g]\mapsto[hg]$ 作用在 $X_{\gamma}$ 上。对任意自由交换子群 $L\subset G_{\gamma}(F)$，它在 $G(F)/G(\calO_{F})$ 上的左平移作用自由（因为稳定子必有限），从而它在 $X_{\gamma}$ 上的作用也自由。

更一般地，对任意离散余紧子群 $L\subset G_{\gamma}(F)$，商群胚 $L\bs X_{\gamma}$ 是{\em 有限型}的，即有有限多个同构类，且每个对象的自同构群有限。对有限型群胚 $Y$，定义 $Y$ 的基数为
\begin{equation}\label{card groupoid}
\#Y:=\sum_{y\in \textup{Ob}(Y)/\cong}\frac{1}{\#\Aut(y)}.
\end{equation}

下一个引理直接来自定义，其证明留作习题 \ref{exIII:Ogamma}。

\begin{lemma}\label{l:Ogamma} 设 $\gamma$ 是 $\frg(F)$ 中的正则半单元，$L\subset G_{\gamma}(F)$ 是任意离散余紧子群。则
\begin{equation*}
O_{\gamma}(\one_{\frg(\calO_{F})})=\frac{1}{\vol(G_{\gamma}(F)/L,\mu_{G_{\gamma}})}\#\left(L\bs X_{\gamma}\right),
\end{equation*}
其中右端的基数按 \eqref{card groupoid} 解释。
\end{lemma}

\subsubsection{$G=\GL_{n}$ 的情形与分式理想}\label{sss:GLn orb formula} 沿用 \S\ref{sss:GLn orb} 的情形。在 $G(F)/G(\calO_{F})$ 与 $F^{n}$ 中 $\calO_{F}$-格的集合的等同（见 \eqref{Gr lattices}）下，有
\begin{equation*}
X_{\gamma}\cong\{\textup{格 }\L\subset F^{n}|\gamma \L\subset \L\}.
\end{equation*}
此双射把 $[g]\in X_{\gamma}$ 送到格 $\L=g\calO^{n}_{F}$。

再给出 $X_{\gamma}$ 的另一种解释。设 $P(x)=x^{n}+a_{1}x^{n-1}+\cdots+a_{n}$ 是 $\gamma$ 的特征多项式。令
\begin{equation*}
A=\calO_{F}[x]/(P(x))
\end{equation*}
为 $\frg(F)$ 中由 $\gamma$ 生成的交换 $\calO_{F}$-子代数。全分式环 $\Frac(A)$ 是次数为 $n$ 的有限 étale $F$-代数，而 $A$ 是其中的一个序。$A$ 在 $F^{n}$ 上的典范作用把 $F^{n}$ 实现为秩 $1$ 的自由 $\Frac(A)$-模。回忆{\em 分式 $A$-理想}指有限生成的 $A$-子模 $M\subset \Frac(A)$。若取元素 $e\in F^{n}$ 作为该 $\Frac(A)$-模结构的基，就得到双射
\begin{equation}\label{frac ideals}
\{\mbox{分式 $A$-理想}\}\bij X_{\gamma},
\end{equation}
它把 $M\subset \Frac(A)$ 送到 $\calO_{F}$-格 $M\cdot e\subset F^{n}$。

借助代数 $A$，有典范同构
\begin{equation*}
G_{\gamma}(F)\cong \Frac(A)^{\times}.
\end{equation*}
此同构使 $G_{\gamma}(F)$ 通过左平移在 $X_{\gamma}$ 上的作用，与 $\Frac(A)^{\times}$ 通过乘法在分式 $A$-理想集合上的作用相互缠绕。

当 $A$ 恰是 Dedekind 整环之积（即 $A$ 是 $\Frac(A)$ 中的极大序）时，所有分式 $A$-理想都是主理想；这等价于说 $G_{\gamma}(F)\cong \Frac(A)^{\times}$ 在 $X_{\gamma}$ 上的作用可迁。一般情形下，主分式理想构成 $G_{\gamma}(F)$ 下的齐性空间 $\Frac(A)^{\times}/A^{\times}$；一般情形下计数 $X_{\gamma}$ 的困难来自环 $A$ 的奇异性。

把 $G_{\gamma}(F)\cong\prod_{i}F^{\times}_{i}$ 上的 Haar 测度 $\mu_{G_{\gamma}}$ 归一化，使 $\prod_{i}\calO^{\times}_{F_{i}}$ 的体积为 $1$。取素元 $t_{i}\in F_{i}$，可作出格 $L_{0}=t^{\ZZ}_{1}\times \cdots\times t^{\ZZ}_{m}\subset G_{\gamma}(F)$。此时 $G_{\gamma}(F)/L_{0}\cong \prod_{i}\calO_{F_{i}}^{\times}$ 的体积为 $1$。因此引理 \ref{l:Ogamma} 给出
\begin{equation}\label{GLn L0}
O_{\gamma}(\one_{\frg(\calO_{F})})=\#(L_{0}\bs X_{\gamma}).
\end{equation}
借助 \eqref{frac ideals}，可以把 \eqref{GLn L0} 解释为：$O_{\gamma}(\one_{\frg(\calO_{F})})$ 等于分式 $A$-理想在乘以 $t_{i}$ 的各次幂之下的等价类的个数。

\subsection{与仿射 Springer 纤维的关系}\label{ss:orb ASF}
从本小节起，限制在 $F$ 为局部函数域的情形，即 $F=k\lr{t}$，其中 $k=\FF_{q}$ 是有限域。设 $\gamma\in\frg(F)$ 正则半单。仿射 Grassmann 簇与仿射 Springer 纤维 $\sX_{\gamma}$ 的定义（\S\ref{s:ASF} 中给出）在基域 $k$ 为有限域时仍有意义，于是得到 $\Gr_{G}$ 的子归纳概形 $\sX_{\gamma}$，两者都定义在 $k$ 上。

下述引理由定义即可看出。

\begin{lemma}\label{l:X gamma k} $k$-有理点集 $\sX_{\gamma}(k)$ 与 \S\ref{sss:X gamma} 中定义的集合 $X_{\gamma}$ 相同，两者都作为 $\Gr_{G}(k)=G(F)/G(\calO_{F})$ 的子集。
\end{lemma}

\subsubsection{} 把基域从 $k$ 变到 $\kbar$，由定理 \ref{th:ft quot} 知 $\L_{\gamma}\bs \sX_{\gamma,\kbar}$ 是 $\kbar$ 上的真代数空间。那里的证明实际上说明此代数空间定义在 $k$ 上，仍记为 $\L_{\gamma}\bs \sX_{\gamma}$。这里要强调：$\L_{\gamma}$ 被看作 $k$ 上的 étale 群概形，其 $\kbar$-点就是此前称为 $\L_{\gamma}$ 的普通群。

由引理 \ref{l:Ogamma} 与引理 \ref{l:X gamma k}，自然会期望轨道积分 $O_{\gamma}(\one_{\frg(\calO_{F})})$ 能表为商 $\L_{\gamma}\bs \sX_{\gamma}$ 上 $k$-点的个数。这一关系在 $G=\GL_{n}$ 时形式最简洁。

\subsubsection{$\GL_{n}$ 的情形}\label{sss:GLn orb coho}
在 \S\ref{sss:GLn orb formula} 的情形中，选取素元 $t_{i}\in F_{i}$ 后定义了格 $L_{0}=t^{\ZZ}_{1}\times \cdots\times t^{\ZZ}_{m}$。把基域从 $k$ 变到 $\kbar$，可以用同样的素元 $t_{i}\in F_{i}$，按 \S\ref{sss:GLn centralizer} 类似地定义格 $\wt\L_{\gamma}$（注意 $F_{i}\wh\otimes_{k}\kbar$ 可能分裂为域之积，但 $t_{i}\otimes 1$ 会投影为每个因子中的素元）。$\Gk$ 在 $\wt\L_{\gamma}$ 上的作用给它 $k$ 上 étale 群概形的结构，正如 $\L_{\gamma}$。有 $\wt\L_{\gamma}(k)=L_{0}$。把定理 \ref{th:ft quot} 中的 $\L_{\gamma}$ 换成 $\wt\L_{\gamma}$，有类似的陈述。特别地，$\wt\L_{\gamma}\bs\sX_{\gamma}$ 是 $k$ 上的真代数空间，且有射影概形到它的满射。

\begin{theorem}\label{th:GLn orb} 设 $G=\GL_{n}$，$\gamma$ 是 $\frg(F)$ 中的正则半单元。把 $G_{\gamma}(F)$ 上的 Haar 测度归一化，使其极大紧子群的体积为 $1$。则
\begin{equation}\label{GLn orb}
O_{\gamma}(\one_{\frg(\calO_{F})})=\#(\wt\L_{\gamma}\bs\sX_{\gamma})(k)=\sum_{i}(-1)^{i}\Tr\left(\Frob_{k}, \cohog{i}{\wt\L_{\gamma}\bs\sX_{\gamma,\kbar},\Ql}\right).
\end{equation}
\end{theorem}

\eqref{GLn orb} 中的第二个等号来自 Grothendieck--Lefschetz 迹公式。与 \eqref{GLn L0} 比较，只需说明 $L_{0}\bs X_{\gamma}=\wt\L_{\gamma}(k)\bs \sX_{\gamma}(k)$ 与 $(\wt\L_{\gamma}\bs\sX_{\gamma})(k)$ 相同。先对商叠的 $k$-点作一些一般性说明。

\subsubsection{商的 $k$-点}\label{sss:k pts quot} 考虑商叠 $\sY=[A\bs X]$，其中 $X$ 是 $k$ 上概形，$A$ 是 $k$ 上作用在 $X$ 上的代数群。按定义，$\sY(k)$ 是由数对 $(S, f)$ 构成的群胚，其中 $S\to \Spec k$ 是左 $A$-主齐性空间，$f:S\to X$ 是 $A$-等变态射。$A$-主齐性空间 $S$ 的同构类由 Galois 上同调 $\cohog{1}{k, A}:=\cohog{1}{\Gk,A(\kbar)}$ 分类。对每个类 $\xi\in \cohog{1}{k,A}$，设 $S_{\xi}$ 是类为 $\xi$ 的 $k$ 上 $A$-主齐性空间。可以用 $X_{\xi}:=A\bs (S_{\xi}\times X)$ 定义 $X$ 在 $k$ 上的扭形式。还有作用在 $k$-概形 $X_{\xi}$ 上的内形式 $A_{\xi}:=\Aut_{A}(S_{\xi})$。容易看出 $A$-等变态射 $f:S_{\xi}\to X$ 与 $X_{\xi}(k)$ 一一对应。因此得到群胚分解
\begin{equation}\label{k pts quot}
\sY(k)=[A\bs X](k)\cong\bigsqcup_{\xi\in\cohog{1}{k,A}}A_{\xi}(k)\bs X_{\xi}(k).
\end{equation}

\subsubsection{定理 \ref{th:GLn orb} 的证明} 已化归为证明 $(\wt\L_{\gamma}\bs\sX_{\gamma})(k)= \wt\L_{\gamma}(k)\bs \sX_{\gamma}(k)$。由 \eqref{k pts quot}，只需说明 $\cohog{1}{k,\wt\L_{\gamma}}=0$。沿用 \S\ref{sss:GLn orb coho} 的记号。由 $\wt\L_{\gamma}$ 的定义可见，作为 $\Gk$-模它有形式
\begin{equation*}
\wt\L_{\gamma}\cong \bigoplus_{i=1}^{m}\Ind^{\Gk}_{\Gal(\kbar/k_{i})}\ZZ,
\end{equation*}
其中 $k_{i}$ 是 $F_{i}$ 的剩余域。因此
\begin{equation*}
\cohog{1}{k,\wt\L_{\gamma}}\cong\bigoplus_{i=1}^{m}\cohog{1}{k_{i},\ZZ}=0.
\end{equation*}
\qed

\subsection{稳定轨道积分} 设定与 \S\ref{ss:orb ASF} 相同。对一般的 $G$，公式 \eqref{GLn orb} 的推广并不直接。具体说来，轨道积分 $O_{\gamma}(\one_{\frg(\calO_{F})})$ 本身并没有上同调解释。问题在于未必能找到 $\wt\L_{\gamma}$ 的类比物、使其一阶 Galois 上同调消没，从而 $(\wt\L_{\gamma}\bs\sX_{\gamma})(k)$ 就简单地等于 $\wt\L_{\gamma}(k)\bs \sX_{\gamma}(k)$。由公式 \eqref{k pts quot}，一个自然的补救办法是把 $\sX_{\gamma}$ 的各扭形式放在一起考虑。这提示我们不要只看轨道积分 $O_{\gamma}(\one_{\frg(\calO_{F})})$，而要看若干个轨道积分 $O_{\gamma'}(\one_{\frg(\calO_{F})})$ 之和。

\subsubsection{稳定共轭} 固定正则半单元 $\gamma\in\frg(F)$。称 $\gamma'\in\frg(F)$ {\em 稳定共轭}于 $\gamma$，若它与 $\gamma$ 在同一 $G(\ov{F})$-轨道中。等价地，$\gamma'$ 稳定共轭于 $\gamma$，当且仅当 $\chi(\gamma')=\chi(\gamma)\in\frc(F)$。对稳定共轭于 $\gamma$ 的 $\gamma'$，可以附上一个 Galois 上同调类 $\inv(\gamma,\gamma')\in\cohog{1}{F,G_{\gamma}}$，它在 $\cohog{1}{F,G}$ 中的像平凡。对应 $\gamma'\mapsto\inv(\gamma,\gamma')$ 给出带点集的双射
\begin{equation}\label{stable class}
\{\gamma'\in\frg(F) \mbox{ 稳定共轭于 } \gamma\}/G(F) \cong \ker(\cohog{1}{F,G_{\gamma}}\to \cohog{1}{F,G}).
\end{equation}

\subsubsection{$\GL_{n}$ 的情形} 当 $G=\GL_{n}$ 时，沿用 \S\ref{sss:GLn centralizer} 的记号。有
\begin{equation*}
\cohog{1}{F,G_{\gamma}}\cong \prod_{i=1}^{m}\cohog{1}{F_{i},\Gm}=\jiao{1}.
\end{equation*}
因此由 \eqref{stable class}，所有稳定共轭于 $\gamma$ 的元素实际上都在 $G(F)$ 中共轭于 $\gamma$。当然，此陈述也可以用友矩阵直接证明。

\subsubsection{$\SL_{2}$ 的情形}\label{sss:st conj SL2} 考虑 $G=\SL_{2}$。设 $a\in k^{\times}-(k^{\times})^{2}$。令
\begin{equation*}
\gamma=\mat{0}{at}{t}{0}, \quad \gamma'=\mat{0}{at^{2}}{1}{0}
\end{equation*}
为 $\sl_{2}(F)$ 中两个正则半单元。由于它们的行列式都是 $-at^{2}$，两者稳定共轭。然而它们在 $\SL_{2}(F)$ 下并不共轭。可以证明 $\gamma$ 的稳定共轭类恰由两个 $\SL_{2}(F)$-轨道组成，分别以 $\gamma$ 与 $\gamma'$ 为代表，见习题 \ref{exIII:stable conj}。

\subsubsection{稳定轨道积分的定义} 设 $\varphi\in\sS(\frg(F))$。定义 $\varphi$ 关于 $\gamma$ 的{\em 稳定轨道积分}为
\begin{equation}\label{SO}
SO_{\gamma}(\varphi)=\sum_{\gamma'}O_{\gamma'}(\varphi),
\end{equation}
其中 $\gamma'$ 遍历 $\frg(F)$ 中稳定共轭于 $\gamma$ 的元素的 $G(F)$-轨道。对稳定共轭于 $\gamma$ 的 $\gamma'$，有 $F$-群之间的{\em 典范同构} $G_{\gamma}\cong G_{\gamma'}$。因此一旦固定 $G_{\gamma}(F)$ 上的 Haar 测度，就得到其他中心化子 $G_{\gamma'}(F)$ 上的典范 Haar 测度。\eqref{SO} 中的 $O_{\gamma'}(\varphi)$ 正是按此选取来定义的。

\subsubsection{上同调的稳定部分}\label{sss:st H} 商群概形 $LG_{\gamma}/\L_{\gamma}$ 作用在 $\L_{\gamma}\bs\sX_{\gamma}$ 上。分支群 $\pi_{0}(LG_{\gamma}/\L_{\gamma})$ 是 $k$ 上 étale 群概形，其 $\kbar$-点作用在 $\cohog{*}{L_{\gamma}\bs\sX_{\gamma,\kbar},\Ql}$ 上。定义此上同调群的{\em 稳定部分}为此作用下的不变量
\begin{equation*}
\cohog{*}{\L_{\gamma}\bs\sX_{\gamma,\kbar}}_{\st}=\cohog{*}{\L_{\gamma}\bs\sX_{\gamma,\kbar}}^{\pi_{0}(LG_{\gamma}/\L_{\gamma})(\kbar)}.
\end{equation*}
可以证明：若把 $\L_{\gamma}$ 换成 $G_{\gamma}(F^{\ur})$ 中任一与 $\L_{\gamma}$ 可公度的 $\Gk$-稳定自由交换子群 $\L$，类似定义的稳定部分上同调与上述典范同构。

\begin{theorem}[Goresky--Kottwitz--MacPherson {\cite[定理 15.8]{GKM}} 与 Ng\^{o} {\cite[推论 8.2.10]{NgoFL}} 的特例]\label{th:orb coho} 有
\begin{equation}\label{d}
SO_{\gamma}(\one_{\frg(\calO_{F})})=\frac{1}{\vol(K_{\gamma},\mu_{G_{\gamma}})}\sum_{i}(-1)^{i}\Tr\left(\Frob_{k}, \cohog{i}{\L_{\gamma}\bs\sX_{\gamma,\kbar},\Ql}_{\st}\right).
\end{equation}
这里 $K_{\gamma}\subset G_{\gamma}(F)$ 是环面 $G_{\gamma}$ 的 parahoric 子群。
\end{theorem}

稳定轨道积分的这一上同调解释，是基本引理之证明的出发点（见 \cite{GKM} 与 \cite{NgoFL}）。

\subsubsection{}\label{sss:K gamma} 简略说明 parahoric 子群 $K_{\gamma}$ 的定义。令 $\tilK_{\gamma}\subset G_{\gamma}(F)$ 为极大紧子群。则存在 $\calO_{F}$ 上典范光滑群概形 $\calG_{\gamma}$，其 $F$-纤维为 $G_{\gamma}$，其 $\calO_{F}$-点为 $\tilK_{\gamma}$。此群概形 $\calG_{\gamma}$ 就是环面 $G_{\gamma}$ 的有限型 Néron 模型。令 $\calG^{\circ}_{\gamma}\subset\calG_{\gamma}$ 为去掉 $\calG_{\gamma}$ 特殊纤维的非中性分支后得到的开子群概形。则按定义 $K_{\gamma}=\calG^{\circ}_{\gamma}(\calO_{F})$。正环路群 $L^{+}\calG^{\circ}_{\gamma}$ 在相差幂零元的意义下是 $LG_{\gamma}$ 的中性分支。

\subsubsection{定理 \ref{th:orb coho} 的证明梗概}\label{sss:pf SO} 下面概述的论证更接近 \cite{GKM} 而非 \cite{NgoFL}。可以证明存在 étale $k$-子群 $\wt\L\subset LG_{\gamma}$，它与 $\L_{\gamma}$ 可公度，且{\em 满射}到 $k$ 上 étale 群概形 $\pi_{0}(LG_{\gamma})$（其核必有限）。群 $\wt\L$ 是 \S\ref{sss:GLn orb coho} 中定义的 $\wt\L_{\gamma}$ 的类比物。作商叠 $[\wt\L\bs \sX_{\gamma}]$，则有同构
\begin{equation}\label{st coho quot}
\cohog{*}{\L_{\gamma}\bs\sX_{\gamma,\kbar},\Ql}_{\st}\cong \cohog{*}{[\wt\L\bs\sX_{\gamma,\kbar}],\Ql}.
\end{equation}
事实上，由 \S\ref{sss:st H} 末尾的讨论，把 $\L_{\gamma}$ 换成可公度的格时上式左端不变，故可把 $\L_{\gamma}$ 缩小，假定 $\L_{\gamma}\subset \wt\L$。另一方面，上式右端可用与映射 $\L_{\gamma}\bs\sX_{\gamma,\kbar}\to[\wt\L\bs\sX_{\gamma,\kbar}]$ 相联系的 Leray 谱序列计算；此映射是有限离散群 $(\wt\L/\L_{\gamma})(\kbar)$ 下的主齐性空间，因此 $\cohog{*}{[\wt\L\bs\sX_{\gamma,\kbar}],\Ql}=\cohog{*}{\L_{\gamma}\bs\sX_{\gamma,\kbar},\Ql}^{(\wt\L/\L_{\gamma})(\kbar)}$。由于按 $\wt\L$ 的选取，$(\wt\L/\L_{\gamma})(\kbar)$ 满射到 $\pi_{0}(LG_{\gamma}/\L_{\gamma})(\kbar)$，便有
\begin{equation*}
\cohog{*}{\L_{\gamma}\bs\sX_{\gamma,\kbar},\Ql}^{\pi_{0}(LG_{\gamma}/\L_{\gamma})(\kbar)}=\cohog{*}{\L_{\gamma}\bs\sX_{\gamma,\kbar},\Ql}^{(\wt\L/\L_{\gamma})(\kbar)},
\end{equation*}
由此即得 \eqref{st coho quot}。

由 \S\ref{sss:k pts quot} 与 \eqref{k pts quot} 的讨论，有
\begin{equation}\label{L bs X decomp}
[\wt\L\bs\sX_{\gamma,\kbar}](k)\cong\bigsqcup_{\xi\in\cohog{1}{k,\wt\L}} \wt\L(k)\bs\sX_{\gamma,\xi}(k).
\end{equation}
包含 $\wt\L\subset LG_{\gamma}$ 给出
\begin{equation*}
\theta: \cohog{1}{k,\wt\L}\surj\cohog{1}{k,LG_{\gamma}}\cong\cohog{1}{F, G_{\gamma}}.
\end{equation*}
第一个满射来自 Lang 定理，因为商 $LG_{\gamma}/\wt\L$ 连通；第二个来自 Lang 的另一定理，它断言 $\cohog{1}{F^{\ur}, G_{\gamma}}$ 消没\footnote{参见 \cite[Ch.X, \S7, p.170, Application and Example (b)]{Serre}。设 $K$ 是剩余域完全的完备离散赋值域，$K^{\ur}$ 为其极大非分歧扩张。Lang 定理断言 $K^{\ur}$ 是 $C_{1}$-域。因此对 $K$ 上任一环面 $A$ 有 $\cohog{1}{K^{\ur}, A}=0$。}。对每个满足 $\sX_{\gamma,\xi}(k)\neq\varnothing$ 的 $\xi\in\cohog{1}{k,\wt\L}$，可以证明 $\theta(\xi)$ 在 $\cohog{1}{F,G}$ 中的像平凡。因此由 \eqref{stable class}，可以对每个 $\xi\in\cohog{1}{k,\wt\L}$ 附上一个稳定共轭于 $\gamma$ 的元素 $\gamma'$，它在 $G(F)$-共轭下唯一，且满足 $\inv(\gamma,\gamma')=\theta(\xi)$。可以证明 $\sX_{\gamma,\xi}(k)$ 与集合 $X_{\gamma'}$ 一一对应。因此 \eqref{L bs X decomp} 蕴含
\begin{equation}\label{pre SO}
\#[\wt\L\bs\sX_{\gamma}](k)=\#\ker(\theta)\sum_{\gamma'}\#(\wt\L(k)\bs X_{\gamma'}),
\end{equation}
其中求和遍历稳定共轭于 $\gamma$ 的那些 $\gamma'$ 的 $G(F)$-轨道。把引理 \ref{l:Ogamma} 用于离散余紧子群 $\wt\L(k)$，得
\begin{equation*}
O_{\gamma'}(\one_{\frg(\calO_{F})})=\frac{1}{\vol(G_{\gamma}(F)/\wt\L(k),\mu_{G_{\gamma}})}\#(\wt\L(k)\bs X_{\gamma'}).
\end{equation*}
代入 \eqref{pre SO} 的右端，得到
\begin{equation}\label{next SO}
\SO_{\gamma}(\one_{\frg(\calO_{F})})=\frac{1}{\#\ker(\theta)\vol(G_{\gamma}(F)/\wt\L(k),\mu_{G_{\gamma}})}\#[\wt\L\bs\sX_{\gamma}](k).
\end{equation}
由 Grothendieck--Lefschetz 迹公式，$\#[\wt\L\bs\sX_{\gamma,\kbar}](k)$ 等于 $\cohog{*}{[\wt\L\bs\sX_{\gamma,\kbar}],\Ql}$ 上的交错 Frobenius 迹；由 \eqref{st coho quot}，后者可以等同于 $\cohog{*}{\L_{\gamma}\bs\sX_{\gamma,\kbar},\Ql}_{\st}$。因此定理归结为等式 \eqref{next SO} 与如下体积等式
\begin{equation}\label{vol id}
\#\ker(\theta)\vol(G_{\gamma}(F)/\wt\L(k),\mu_{G_{\gamma}})=\vol(K_{\gamma},\mu_{G_{\gamma}}).
\end{equation}
为证此式，令 $C=L^{+}\calG^{\circ}_{\gamma}\cap \wt\L$（其中 $\calG^{\circ}_{\gamma}$ 是 $G_{\gamma}$ 的连通 Néron 模型，其 $\calO_{F}$-点为 $K_{\gamma}$，见 \S\ref{sss:K gamma}）。这是 $k$ 上有限 étale 群。有 $k$ 上群归纳概形的短正合列
\begin{equation*}
1\to C\to L^{+}\calG^{\circ}_{\gamma}\times\wt\L\to (LG_{\gamma})^{\red}\to 1.
\end{equation*}
与之相伴的 $\Gk$-上同调六项正合列给出
\begin{equation*}
1\to C(k)\to K_{\gamma}\times \wt\L(k)\to G_{\gamma}(F)\to \cohog{1}{k,C}\to \ker(\theta)\to1,
\end{equation*}
由此并利用 $\#C(k)=\#\cohog{1}{k,C}$，即得 \eqref{vol id}。\qed

\subsection{$\SL_{2}$ 中的例子}
由定理 \ref{th:orb coho}，为计算轨道积分，不仅要知道仿射 Springer 纤维 $\sX_{\gamma}$ 的几何，还要知道 Frobenius 在其上同调上的作用。\S\ref{ss:aff examples} 中已见过代数闭域上仿射 Springer 纤维的许多例子，这里重点放在 Frobenius 作用上。

本小节取 $G=\SL_{2}$，并假设 $\chark>2$。我们在此情形计算几个轨道积分，并用显式计算验证定理 \ref{th:orb coho}。

\subsubsection{非分歧情形：$\gamma$}\label{unram sl2} 设 $\gamma=\mat{0}{at}{t}{0}\in\sl_{2}(F)$ 是正则半单元，其中 $a\in k^{\times}-(k^{\times})^{2}$。

设 $E\subset\mathfrak{gl}_{2}(F)$ 是 $\gamma$ 在 $\mathfrak{gl}_{2}(F)$ 中的中心化子。则 $E=\{u+v\gamma|u,v\in F\}$ 是 $F$ 的非分歧二次扩张，由添加 $\sqrt{a}$ 得到。因此 $E=k_{E}\lr{t}$，其中 $k_{E}= k(\sqrt{a})$。有 $G_{\gamma}(F)=(E^{\times})^{\Nm=1}=\ker(\Nm: E^{\times}\to F^{\times})$，它是紧群。固定 $G_{\gamma}(F)$ 上总体积为 $1$ 的 Haar 测度。

设 $\sX_{\gamma}$ 是 $\gamma$ 的仿射 Springer 纤维，它是 $k$ 上概形。由引理 \ref{l:Ogamma}，
\begin{equation*}
O_{\gamma}(\one_{\frg(\calO_{F})})=\#X_{\gamma}=\#\sX_{\gamma}(k).
\end{equation*}
在 \S\ref{P1 chain} 中已证明 $\sX_{\gamma,\kbar}$ 是可数多条有理曲线 $C_{n}\cong\PP^{1}$ 之并，指标为整数 $n\in\ZZ$。由于 $\gamma$ 在 $E$ 上可对角化，每个分支 $C_{n}$ 实际上定义在 $k_{E}=k(\sqrt{a})$ 上。格 $\L_{\gamma}\subset G_{\gamma}(F^{\ur})$ 包含在 $G_{\gamma}(E)\cong E^{\times}$ 中，并由素元 $t\in E$ 生成。把分支 $C_{n}$ 编号，使 $t\in \L_{\gamma}$ 把 $C_{n}$ 送到 $C_{n+1}$。令 $x_{n+1/2}:=C_{n}\cap C_{n+1}$，它是 $\sX_{\gamma}$ 的 $k_{E}$-点。

非平凡对合 $\sigma\in\Gal(k_{E}/k)$ 在 $G_{\gamma}(E)$ 上的作用是取逆，因此它在 $\L_{\gamma}$ 上的作用也是取逆。标准格 $\calO^{2}_{F}$ 同时落在 $C_{0}$ 与 $C_{1}$ 中，故它就是点 $x_{1/2}$。由于 $x_{1/2}$ 定义在 $k$ 上，它被 $\sigma$ 固定。$\sigma$ 在 $\sX_{\gamma,k_{E}}$ 上的作用与其在 $\L_{\gamma}$ 上的作用（取逆）相容，因此唯一可能是
\begin{equation*}
\sigma(C_{n})=C_{1-n}, \quad \sigma(x_{n+1/2})=x_{-n+1/2}, \quad\forall n\in\ZZ.
\end{equation*}
$\sigma$ 的作用可用下图表示
\begin{equation*}
\xymatrix{ x_{-3/2} \ar@/^{1pc}/[rrrr]_{\sigma}
\ar@{=} '[dr]_{C_{-1}} '[rr]_{C_{0}}  '[drrr]_{C_{1}}  [rrrr]_{C_{2}}
&  & x_{1/2}  & & x_{5/2} \\
& x_{-1/2}\ar@/_1pc/[rr]^{\sigma} &  & x_{3/2}}
\end{equation*}
图中每条双线代表一条 $\PP^{1}$。由此看出 $X_{\gamma}=\sX_{\gamma}(k)$ 只含一个点 $x_{1/2}$，即标准格 $\calO^{2}_{F}$。这意味着
\begin{equation}\label{Oga}
O_{\gamma}(\one_{\frg(\calO_{F})})=1.
\end{equation}

\subsubsection{非分歧情形：$\gamma'$} 现考虑元素 $\gamma'=\mat{0}{at^{2}}{1}{0}\in\sl_{2}(F)$。习题 \ref{exIII:stable conj} 说明 $\gamma'$ 稳定共轭于 $\gamma$，但在 $\SL_{2}(F)$ 下不与 $\gamma$ 共轭。不过 $\gamma'$ 在 $\SL_{2}(E)$ 下共轭于 $\gamma$。因此仿射 Springer 纤维 $\sX_{\gamma'}$ 在 $k_{E}$ 上仍与 $\sX_{\gamma}$ 形貌相同，只是 $\sigma\in\Gal(k_{E}/k)$ 的作用不同。

考虑 $\sX_{\gamma',k_{E}}$ 的一个分支，其 $k_{E}$-点由满足 $t\calO_{E}\oplus\calO_{E}\subset \L\subset \calO_{E}\oplus t^{-1}\calO_{E}$ 的 $\gamma'$-稳定格 $\L\subset E^{2}$ 组成。此分支由定义在 $k$ 上的条件截出，故在 $\sigma$ 下稳定，记之为 $C'_{0}$。把 $\sX_{\gamma',k_{E}}$ 的其余分支编号为 $C'_{n}$（$n\in\ZZ$），使生成元 $t\in\L_{\gamma'}$ 把 $C'_{n}$ 送到 $C'_{n+1}$。令 $x'_{n+1/2}=C'_{n}\cap C'_{n+1}\in \sX_{\gamma'}(k_{E})$。由于 $\sigma$ 在 $\sX_{\gamma,k_{E}}$ 上的作用与其在 $\L_{\gamma'}$ 上的作用（取逆）相容，唯一可能是
\begin{equation*}
\sigma(C'_{n})=C'_{-n}, \quad \sigma(x'_{n+1/2})=x'_{-n-1/2},\quad\forall n\in\ZZ.
\end{equation*}
$\sigma$ 的作用可用下图表示
\begin{equation*}
\xymatrix{ & & x'_{-1/2} \ar@/_{1pc}/@<-1ex>[r]_{\sigma}& x'_{1/2}\\
x'_{-5/2}  \ar@/_2pc/[rrrrr]^{\sigma} 
\ar@{=} '[r]^{C'_{-2}} '[urr]_{C'_{-1}}  '[urrr]_{C'_{0}}  '[rrrr]_{C'_{1}} [rrrrr]^{C'_{2}}
&  x'_{-3/2} \ar@/_{1pc}/[rrr]^{\sigma}& & & x'_{3/2}  & x'_{5/2}}
\end{equation*}
因此没有点 $x'_{n+1/2}$ 定义在 $k$ 上。分支 $C'_{0}$ 是唯一定义在 $k$ 上的分支，它在 $k$ 上必同构于 $\PP^{1}$，因为它在 $k_{E}$ 上如此（有限域上没有非平凡 Brauer--Severi 簇）。于是 $\sX_{\gamma'}(k)=C'_{0}(k)\cong\PP^{1}(k)$ 有 $q+1$ 个元素。因此
\begin{equation}\label{Ogb}
O_{\gamma'}(\one_{\frg(\calO_{F})})=q+1.
\end{equation}
把 \eqref{Oga} 与 \eqref{Ogb} 相加，得到
\begin{equation*}
SO_{\gamma}(\one_{\frg(\calO_{F})})=q+2.
\end{equation*}

\subsubsection{非分歧情形：上同调} 商 $\L_{\gamma}\bs\sX_{\gamma}$ 是结点有理曲线，由 $\PP^{1}$ 把两个 $k$-点粘成一个结点得到。

现在考虑商 $\L_{\gamma'}\bs\sX_{\gamma'}$。在 $k_{E}$ 上它同样是结点有理曲线，只有一个结点 $y$，它是所有 $x_{j}$ 的像。虽然 $y$ 是商 $\L_{\gamma'}\bs\sX_{\gamma'}$ 的 $k$-点，但它的任一原像 $x_{j}$ 都不定义在 $k$ 上。因此 $C'_{0}(k)$ 中的 $q+1$ 个点仍单射地映入商，此外还有点 $y$。于是 $(\L_{\gamma'}\bs\sX_{\gamma'})(k)$ 由 $q+2$ 个点组成。由于 $LG_{\gamma'}/\L_{\gamma'}$ 连通，$\L_{\gamma'}\bs\sX_{\gamma'}$ 的上同调的稳定部分就是整个 $\cohog{*}{\L_{\gamma'}\bs\sX_{\gamma'}}$，其上的 Frobenius 迹的交错和等于 $(\L_{\gamma'}\bs\sX_{\gamma'})(k)$ 的基数。此特例验证了公式 \eqref{d}。还要指出：$\sigma$ 在一维空间 $\cohog{1}{\L_{\gamma'}\bs\sX_{\gamma',\kbar},\Ql}$ 上的作用为 $-1$，因此 Frobenius 的 Grothendieck--Lefschetz 迹公式给出 $\#(\L_{\gamma'}\bs\sX_{\gamma'})(k)= 1+1+q$，而非 $1-1+q=q$；后者是把 $\PP^{1}$ 上两个 $k$-点等同所得的结点有理曲线上 $k$-点的个数。

\subsubsection{分歧情形：轨道积分}\label{sss:ram orb SL2} 考虑元素 $\gamma=\mat{0}{t^{2}}{t}{0}$ 与 $\gamma'=\mat{0}{at^{2}}{a^{-1}t}{0}$，其中 $a\in k^{\times}-(k^{\times})^{2}$。同样 $\gamma$ 与 $\gamma'$ 稳定共轭，但在 $\SL_{2}(F)$ 中不共轭。

设 $E$ 是 $\gamma$ 在 $\gl_{2}(F)$ 中的中心化子。则 $E$ 是 $F$ 的分歧二次扩张，$G_{\gamma}(F)=(E^{\times})^{\Nm=1}$。类似地，设 $E'$ 是 $\gamma'$ 在 $\gl_{2}(F)$ 中的中心化子。则 $E'$ 是 $F$ 的另一分歧二次扩张，$G_{\gamma'}(F)=(E'^{\times})^{\Nm=1}$。在紧群 $G_{\gamma}(F)$ 与 $G_{\gamma'}(F)$ 上选取总体积为 $1$ 的 Haar 测度。

两种情形下 $\sX_{\gamma}$ 与 $\sX_{\gamma'}$ 都同构于 $k$ 上的 $\PP^{1}$。事实上 \S\ref{SL2 ram dim 1} 已证明它们在 $\kbar$ 上同构于 $\PP^{1}$，故在 $k$ 上也必同构于 $\PP^{1}$。$\sX_{\gamma}(k)$ 与 $\sX_{\gamma'}(k)$ 都由满足 $t\calO_{F}\oplus\calO_{F}\subset \L\subset \calO_{F}\oplus t^{-1}\calO_{F}$ 的格组成，因此作为 $\Gr_{G}$ 的子簇有 $\sX_{\gamma}=\sX_{\gamma'}$。由引理 \ref{l:Ogamma}，
\begin{equation*}
O_{\gamma}(\one_{\frg(\calO_{F})})=O_{\gamma'}(\one_{\frg(\calO_{F})})=\#\sX_{\gamma}(k)=q+1.
\end{equation*}
因此
\begin{equation}\label{ram so}
SO_{\gamma}(\one_{\frg(\calO_{F})})=2(q+1).
\end{equation}

\subsubsection{分歧情形：上同调} 在 \S\ref{sss:ram orb SL2} 的设定下 $\L_{\gamma}=0$。$LG_{\gamma}$ 的分支群是 $\ZZ/2\ZZ$，但它对 $\sX_{\gamma}$ 的作用平凡。因此上同调的稳定部分就是整个 $\cohog{*}{\sX_{\gamma},\Ql}$，其上 Frobenius 的交错和给出 $\sX_{\gamma}(k)\cong\PP^{1}(k)$ 的基数。
然而 $G_{\gamma}(F)=(E^{\times})^{\Nm=1}$ 的 parahoric 子群在其中指数为 $2$（$K_{\gamma}$ 由那些在 $k$ 中约化为 $1$ 的 $e\in (E^{\times})^{\Nm=1}$ 组成）。因此公式 \eqref{d} 的右端在 $\#\sX_{\gamma}(k)$ 前多出因子 $2=\vol(K_{\gamma},d_{\gamma}g)^{-1}$。这与 \eqref{ram so} 相符，可见特例中公式 \eqref{d} 成立。

\subsection{关于基本引理的评注} 回到 \S\ref{unram sl2} 的情形。若取 $O_{\gamma'}(\one_{\frg(\calO_{F})})$ 与 $O_{\gamma}(\one_{\frg(\calO_{F})})$ 之差而非其和，会发生什么？这一差是否有类似于定理 \ref{th:orb coho} 的几何解释？

\subsubsection{$\kappa$-轨道积分} 线性组合 $O_{\gamma}(\one_{\frg(\calO_{F})})-O_{\gamma'}(\one_{\frg(\calO_{F})})$ 是 $\kappa$-轨道积分的一例。更一般地，设 $\kappa$ 是 $\cohog{1}{F,G_{\gamma}}$ 的特征标，定义 $\varphi\in \sS(\frg(F))$ 的 $\kappa$-轨道积分为
\begin{equation*}
O^{\kappa}_{\gamma}(\varphi)=\sum_{\gamma'}\kappa(\inv(\gamma,\gamma'))O_{\gamma'}(\varphi),
\end{equation*}
求和遍历 $\gamma$ 的稳定共轭类中的 $G(F)$-轨道。

\subsubsection{基本引理的陈述} Langlands--Shelstad 猜想，又称{\em 基本引理}，断言：$\gamma\in\frg(F)$ 的 $\one_{\frg(\calO_{F})}$ 的 $\kappa$-轨道积分，等于较小群 $H$ 的某个元素 $\gamma_{H}\in\frh(F)$ 的稳定轨道积分，相差一个简单因子。用公式写出，基本引理即恒等式
\begin{equation*}
O^{\kappa}_{\gamma}(\one_{\frg(\calO_{F})})=\Delta(\gamma,\gamma_{H})\cdot SO_{\gamma_{H}}(\one_{\frh(\calO_{F})}).
\end{equation*}
较小群 $H$ 同时依赖于 $G$ 与 $\kappa$，称为 $(G,\kappa)$ 的{\em 内窥群}。数 $\Delta(\gamma,\gamma_{H})$ 称为{\em 转移因子}；当 $\gamma$ 从它的稳定共轭类中适当选取、且 $G_{\gamma}$ 与 $H_{\gamma_{H}}$ 上的测度选取适当时，它是 $q$ 的整数次幂（依赖于 $\gamma$ 与 $\gamma_{H}$）。

\subsubsection{一个简单情形} 在 \S\ref{unram sl2} 的情形中，取 $\cohog{1}{F,G_{\gamma}}\cong\ZZ/2\ZZ$ 上的非平凡特征标 $\kappa$，相应的内窥群 $H$ 同构于环面 $G_{\gamma}$；但一般情形下它未必同构于 $G$ 的子群。此时基本引理是恒等式
\begin{equation*}
O_{\gamma'}(\one_{\frg(\calO_{F})})-O_{\gamma}(\one_{\frg(\calO_{F})})=q=q\cdot SO_{\gamma_{H}}(\one_{\frh(\calO_{F})}),
\end{equation*}
其中把 $H$ 与 $G_{\gamma}$ 等同后取 $\gamma_{H}=\gamma$。

另一方面，在 \S\ref{sss:ram orb SL2} 的分歧情形中，$\gamma$ 关于非平凡 $\kappa$ 的 $\kappa$-轨道积分消没。这也许无需显式计算轨道积分即可解释：一般说来，为使 $O^{\kappa}_{\gamma}(\one_{\frg(\calO_{F})})$ 可能非零，$\kappa$ 必须经过 $\cohog{1}{F,G_{\gamma}}$ 的进一步商群分解。确切陈述见 \cite[命题 8.2.7]{NgoFL}。

\subsubsection{关于证明的评注} 一般 $G$ 与函数域 $F$ 的基本引理由 B-C. Ng\^o 证明 \cite{NgoFL}。定理 \ref{th:orb coho} 可以推广到 $\kappa$-轨道积分，只需把 $\L_{\gamma}\bs\sX_{\gamma}$ 的上同调的稳定部分换成 $\kappa$-同型部分。Goresky、Kottwitz 与 MacPherson \cite{GKM} 用这一推广把基本引理重述为仿射 Springer 纤维 $\sX_{\gamma}$ 与其内窥表亲 $\sX^{H}_{\gamma_{H}}$ 的上同调群之间的关系。他们还证明了一些特殊但极不平凡情形下的基本引理。Ng\^o 的证明建立在这一上同调重述之上，但另用了一个新要素，即 Hitchin 纤维——它可以看作仿射 Springer 纤维的「整体」类比。这将是下一讲的主题。

\subsection{习题}

本组习题中，$k=\FF_{q}$ 是 $\chark\neq2$ 的有限域，$F=k\lr{t}$。

\subsubsection{}\label{exIII:K inv} 设 $H$ 是 $F$ 上代数群，$f\in\sS(H(F))$。证明存在紧开子群 $K\subset H(F)$，使 $f$ 在 $K$ 下左右不变。

\subsubsection{}\label{exIII:Ogamma} 证明引理 \ref{l:Ogamma}。

\subsubsection{}\label{exIII:stable conj} 设 $a\in k^{\times}-(k^{\times})^{2}$，$G=\SL_{2}$。考虑 \S\ref{sss:st conj SL2} 中的矩阵 $\gamma,\gamma'$。证明 $\gamma$ 与 $\gamma'$ 稳定共轭，但在 $\SL_{2}(F)$ 下不共轭。

\subsubsection{} 设 $F$ 上 $G=\GL_{n}$，$\gamma=\diag(\gamma_{1},\cdots,\gamma_{m})$ 是 $\frg(F)^{\rs}$ 中的块对角矩阵，其中 $\gamma_{i}\in\mathfrak{gl}_{n_{i}}(F)^{\rs}$。考虑轨道积分 $O_{\gamma}(\one_{\frg(\calO_{F})})$ 当 $q=\#k$ 趋于无穷时的渐近行为。求最小的整数 $d$，使
\begin{equation*}
O_{\gamma}(\one_{\frg(\calO_{F})})=O(q^{d})
\end{equation*}
（$q\to\infty$）。注意右端的 $O(\cdot)$ 是分析学家的 $O$，而左端的 $O$ 指轨道积分\zy{原文此处玩双关：$O_{\gamma}$ 的 $O$ 是 orbital integral 的首字母，$O(q^{d})$ 的 $O$ 是大 $O$ 记号。译文保留两者同形。来源：原文 \texttt{ZhiweiYunPCMI.tex} 第 2141 行原文的自注。}。能否用 $\gamma_{i}$ 的特征多项式解释 $d$？

关于 $O_{\gamma}(\one_{\frg(\calO_{F})})$ 的显式估计，见 \cite{Y-Orb}。

\subsubsection{} 设 $G=\SL_{3}$，$\gamma=\left(\begin{array}{ccc} 0 & 0 & t^{4}\\ 1 & 0 & 0 \\ 0 & 1& 0 \end{array}\right)\in\frg(F)$。计算 $O_{\gamma}(\one_{\frg(\calO_{F})})$。

{\em 提示}：使用 \S\ref{sss:LS} 中引入的胞腔分解。

\subsubsection{} 设 $G=\SL_{2}$。设 $f$ 是满足以下条件的元素 $X\in\frg(\calO_{F})$ 的特征函数：其约化 $\ov{X}\in\frg(k)$ 是正则幂零元。设 $\gamma\in\frg(F)$ 是正则半单元。
\begin{enumerate}
\item 证明除非 $\det(\gamma)\in t\calO_{F}$，有 $O_{\gamma}(f)=0$。
\item 当 $\det(\gamma)\in t\calO_{F}$ 时，证明
\begin{equation*}
O_{\gamma}(f)=O_{\gamma}(\one_{\frg(\calO_{F})})-O_{t^{-1}\gamma}(\one_{\frg(\calO_{F})}).
\end{equation*}
\end{enumerate}

\subsubsection{} 设 $G=\GL_{2}$，$\gamma=\mat{0}{t^{n}}{1}{0}$，其中 $n\geq 1$ 为奇数。设 $G(F)_{d}$ 是满足 $\val_{F}(\det g)=d$ 的 $g\in G(F)$ 之集。固定 $G(F)$ 上的 Haar 测度，使 $G(\calO_{F})$ 的体积为 $1$。证明对任意整数 $d\geq0$，
\begin{equation*}
\int_{G(F)_{d}}\one_{\frg(\calO_{F})}(g^{-1}\gamma g)\one_{\calO^{2}_{F}}((0,1)g)dg
\end{equation*}
等于平面曲线 $y^{2}-t^{n}=0$ 的闭子概形 $Z$ 的个数，其中 $Z$ 满足：(1) $Z$ 的底拓扑空间是点 $(y,t)=(0,0)$；(2) $\dim_{k}\calO_{Z}=d$。

注：此习题把轨道积分与曲线上的点 Hilbert 概形联系起来。这一联系已被 \cite{Y-Orb} 用来给出 $\GL_{n}$ 的轨道积分的估计。整体类比见 \S\ref{sss:MHilb} 与 \S\ref{sss:Hilb}。

\section{第四讲：Hitchin 纤维化}

20 世纪 80 年代后半期，Hitchin 在研究规范场论时引入了著名的可积系统，即 Higgs 丛的模空间。差不多同时，Kazhdan 与 Lusztig 引入仿射 Springer 纤维，作为 Springer 纤维的自然类比。此后十余年里这两个对象一直彼此无关，直到 B-C. Ng\^o 看出两者之间的联系。Ng\^o 的基本洞见可以概括为：Hitchin 纤维是仿射 Springer 纤维的整体类比，而仿射 Springer 纤维是 Hitchin 纤维的局部版本。这里的「整体」指涉及整体函数域或代数曲线的对象，而不只是局部函数域或形式圆盘。这一观察连同精巧的技术工作，使 Ng\^o 证明了 Langlands 与 Shelstad 所猜想的轨道积分的基本引理。下面在略为一般的设定中回顾 Hitchin 的可积系统，并精确说明它与仿射 Springer 纤维的联系。

\subsection{Hitchin 模叠}

\subsubsection{设定} 回到 \S\ref{ss:notation} 的设定。此外，固定 $k$ 上一条代数曲线 $X$（假定 $k$ 代数闭），它光滑、射影且连通。

\subsubsection{向量丛的模叠}
存在模叠 $\Bun_{n}$，分类 $X$ 上秩 $n$ 的向量丛。对任意 $k$-代数 $R$，$\Bun_{n}(R)$ 是 $X_{R}:=X\times_{\Spec k}\Spec R$ 上秩 $n$ 向量丛（局部自由凝聚层）的群胚。叠 $\Bun_{n}$ 是代数叠，见 \cite[定理 4.6.2.1]{LMB}。此外，它在 $k$ 上光滑且局部有限型。

\subsubsection{$G$-主齐性空间}
回忆 $X$ 上的{\em（右）$G$-主齐性空间}是指 $X$ 上概形 $\calE$，带有逐纤维的 $G$-作用，使得在 $X$ 的 étale 拓扑下局部地，$\calE$ 变成 $G\times X$，而 $G$-作用变成 $G$ 在第一个因子上的右平移作用。

对一般约化群 $G$，有 $X$ 上 $G$-主齐性空间的模叠 $\Bun_{G}$。对 $k$-代数 $R$，$\Bun_{G}$ 的 $R$-点是 $X_{R}$ 上 $G$-主齐性空间的群胚。于是 $\Bun_{G}$ 也是 $k$ 上光滑且局部有限型的代数叠。

\subsubsection{相伴丛}\label{sss:ass bun} 设 $(V,\rho)$ 是 $G$ 的 $k$-表示，$\calE$ 是 $X$ 上 $G$-主齐性空间。则存在 $X$ 上向量丛 $\rho(\calE)$，其全空间为
\begin{equation*}
\textup{Tot}(\rho(\calE))=G\bs (\calE\times V),
\end{equation*}
其中 $G$ 通过 $g\cdot(e,v)=(eg^{-1}, \rho(g)v)$ 作用在 $(e,v)\in\calE\times V$ 上。称向量丛 $\rho(\calE)$ 相伴于 $\calE$ 与 $\rho$。

当 $G=\GL_{n}$ 时，有群胚的等价
\begin{equation}\label{GLn vec}
\{ \mbox{$X$ 上秩 $n$ 的向量丛 $\calV$}\}\cong\{\mbox{$X$ 上的 $\GL_{n}$-主齐性空间 $\calE$}\}.
\end{equation}
方向 $\calV\mapsto\calE$ 把向量丛 $\calV$ 送到 $\calV$ 的标架构成的 $\GL_{n}$-主齐性空间，即取 $\calE=\un{\Isom}_{X}(\calO^{n}_{X},\calV)$，附带 $\GL_{n}$ 在平凡丛 $\calO^{n}_{X}$ 上的自然作用。另一方向 $\calE\mapsto\calV$ 把 $X$ 上 $\GL_{n}$-主齐性空间 $\calE$ 送到相伴于 $\calE$ 与 $\GL_{n}$ 的标准表示 $\St$ 的向量丛 $\St(\calE)$。等价 \eqref{GLn vec} 给出叠的典范同构 $\Bun_{n}\cong\Bun_{\GL_{n}}$。

\subsubsection{}\label{sss:Bun Sp} 对其他典型群 $G$，$G$-主齐性空间有更显明的向量丛描述。例如当 $G=\SL_{n}$ 时，$X$ 上 $G$-主齐性空间等同于数对 $(\calV,\iota)$，其中 $\calV$ 是 $X$ 上秩 $n$ 的向量丛，$\iota$ 是线丛同构 $\wedge^{n}\calV\cong\calO_{X}$。

当 $G=\Sp_{2n}$ 时，$X$ 上 $G$-主齐性空间 $\calE$ 的群胚等价于数对 $(\calV,\omega)$ 的群胚，其中 $\calV$ 是 $X$ 上秩 $2n$ 的向量丛，$\omega:\wedge^{2}\calV\to\calO_{X}$ 是凝聚层的 $\calO_{X}$-线性映射，且在每个几何纤维上给出辛形式。一个方向把 $\Sp_{2n}$-主齐性空间 $\calE$ 送到 $(\calV,\omega)$：其中 $\calV=\St(\calE)$ 是相伴于 $\calE$ 与 $\Sp_{2n}$ 的标准表示 $\St$ 的向量丛，而 $\calV$ 上的辛形式 $\omega$ 来自 $\Sp_{2n}$-表示之间的典范映射 $\wedge^{2}(\St)\to \one$（其中 $\one$ 是平凡表示）。

\subsubsection{Higgs 丛} 固定 $X$ 上线丛 $\calL$。$X$ 上{\em 秩 $n$ 的 $\calL$-扭 Higgs 丛}是指数对 $(\calV,\varphi)$，其中 $\calV$ 是 $X$ 上秩 $n$ 的向量丛，$\varphi:\calV\to\calV\otimes\calL$ 是 $\calO_{X}$-线性映射。自同态 $\varphi$ 称为 $\calV$ 上的{\em Higgs 场}。

存在模叠 $\calM_{n,\calL}$，分类 $X$ 上秩 $n$ 的 $\calL$-扭 Higgs 丛。忘掉 Higgs 场的态射 $\calM_{n,\calL}\to \Bun_{n}$ 可表。因此 $\calM_{n,\calL}$ 也是 $k$ 上代数叠。

\subsubsection{$G$-Higgs 丛} {\em $\calL$-扭 $G$-Higgs 丛}是指数对 $(\calE,\varphi)$，其中 $\calE$ 是 $X$ 上 $G$-主齐性空间，$\varphi$ 是 $X$ 上向量丛 $\Ad(\calE)\otimes\calL$ 的整体截面。这里 $\Ad(\calE)$ 是相伴于 $\calE$ 与 $G$ 的伴随表示 $(\frg,\Ad)$ 的向量丛，见 \S\ref{sss:ass bun}。称 $\varphi$ 是 $\calE$ 上的{\em $\calL$-扭 Higgs 场}。

当 $G=\GL_{n}$ 时，$\calL$-扭 $G$-Higgs 丛的概念等价于秩 $n$ 的 $\calL$-扭 Higgs 丛。事实上，对每个 $\calL$-扭 $G$-Higgs 丛 $(\calE,\varphi)$，得到 Higgs 丛 $(\calV=\St(\calE),\phi)$；其中把 $\phi:\calV\to\calV\otimes\calL$ 看作 $\underline{\End}(\calV)\otimes\calL$ 的整体截面，它在典范同构 $\underline{\End}(\calV)\cong\Ad(\calE)$ 下对应 $\varphi$。

还有 $X$ 上 $\calL$-扭 Higgs $G$-主齐性空间的模叠 $\calM_{G,\calL}$。$\calM_{G,\calL}$ 的 $R$-点是 $X_{R}$ 上 $\calL_{R}$-扭 $G$-Higgs 丛的群胚，其中 $\calL_{R}$ 表示 $\calL$ 到 $X_{R}$ 的拉回。遗忘态射 $\calM_{G,\calL}\to \Bun_{G}$ 可表，故 $\calM_{G,\calL}$ 是 $k$ 上代数叠。当 $G=\GL_{n}$ 时，有叠的典范同构 $\calM_{n,\calL}\cong\calM_{\GL_{n},\calL}$。

\subsubsection{例子} 当 $G=\SL_{n}$ 时，$X$ 上 $\calL$-扭 $G$-Higgs 丛等同于三元组 $(\calV,\iota,\varphi)$，其中 $\calV$ 是 $X$ 上秩 $n$ 的向量丛，$\iota:\wedge^{n}\calV\isom\calO_{X}$，且 $\varphi: \calV\to\calV\otimes\calL$ 满足 $\Tr(\varphi)=0$。

当 $G=\Sp_{2n}$ 时，$X$ 上 $\calL$-扭 $G$-Higgs 丛等同于三元组 $(\calV,\omega,\varphi)$，其中 $\calV$ 是 $X$ 上秩 $2n$ 的向量丛，$\omega:\wedge^{2}\calV\to\calO_{X}$ 逐纤维非退化，且 $\varphi: \calV\to\calV\otimes\calL$ 使得对 $\calV$ 的一切局部截面 $u$ 与 $v$，$\omega(\varphi(u),v)+\omega(u,\varphi(v))=0$ 作为 $\calL$ 的局部截面成立。

\subsubsection{Hitchin 模叠作为余切丛}
Hitchin 的原始论文 \cite{Hi} 考虑了 $\calL=\omega_{X}$ 为 $X$ 上 $1$-形式层的特例。此特例尤其重要，因为 $\calM_{G,\omega_{X}}$ 与 $\Bun_{G}$ 的余切丛密切相关。对点 $\calE\in\Bun_{G}(R)$，即 $X_{R}$ 上 $G$-主齐性空间，$\Bun_{G}$ 在 $\calE$ 处的余切复形由复形 $\Ad(\calE)^{\vee}\otimes\omega_{X}$ 在 $X_{R}$ 上的导出整体截面给出。用 $\frg$ 上的 Killing 形式可以把 $\Ad(\calE)^{\vee}$ 与 $\Ad(\calE)$ 等同，因此 $\Bun_{G}$ 在 $\calE$ 处的余切复形是 $\bR\Gamma(X_{R}, \Ad(\calE)\otimes\omega_{X})$，它集中在 $0$ 与 $1$ 次。特别地，当 $\calE\in \Bun_{G}(k)$ 的自同构群有限（例如 $\calE$ 稳定）时，$\calE$ 处的{\em Zariski 余切空间}是 $\cohog{0}{X,\Ad(\calE)\otimes\omega_{X}}$；也就是说，$\Bun_{G}$ 在 $\calE$ 处的余切向量等同于 $\calE$ 上 $\omega_{X}$-扭 Higgs 场。因此 $T^{*}\Bun_{G}$（适当定义）与 $\calM_{G,\omega_{X}}$ 共享开子叠 $T^{*}\Bun^{s}_{G}$，其中 $\Bun^{s}_{G}$ 是稳定 $G$-丛的开子叠。

\subsection{Hitchin 纤维化}

\subsubsection{$\GL_{n}$ 的 Hitchin 纤维化}
对 $X$ 上 $\calL$-扭 Higgs 丛 $(\calV,\varphi)$，在 $X$ 上局部地可把 $\varphi$ 看作元素为 $\calL$ 的局部截面的矩阵，从而可取此矩阵的特征多项式。此多项式各系数是 $\calL^{\otimes i}$ 的良定义整体截面，$1\leq i\leq n$。更内蕴地说，$\varphi$ 诱导映射 $\wedge^{i}\varphi:\wedge^{i}\calV\to\wedge^{i}\calV\otimes\calL^{\otimes i}$，可取
\begin{equation*}
a_{i}(\varphi):=\Tr(\wedge^{i}\varphi)\in\cohog{0}{X,\calL^{\otimes i}}.
\end{equation*}
这样就定义了态射
\begin{equation*}
f: \calM_{n,\calL}\to\calA_{n,\calL}:=\prod_{i=1}^{n}\cohog{0}{X,\calL^{i}},
\end{equation*}
它把 $(\calV,\varphi)$ 送到 $(a_{1}(\varphi),\cdots, a_{n}(\varphi))$。把 $\calA_{n,\calL}$ 看作 $k$ 上仿射空间。态射 $f$ 称为 $G=\GL_{n}$ 情形下的{\em Hitchin 纤维化}。

\subsubsection{一般的 Hitchin 纤维化}\label{sss:Hit fib general} 对 \S\ref{ss:notation} 中一般连通约化群 $G$，$\GL_{n}$ 情形中特征多项式的系数由 $\frg$ 上 $G$-不变的基本多项式代替。回忆 $\frc=\frg\sslash G=\Spec\Sym(\frg^{*})^{G}$。Chevalley 定理说 $\Sym(\frg^{*})^{G}\cong\Sym(\frt^{*})^{W}$，而后者是 $r$ 元多项式环。固定 $\Sym(\frg^{*})^{G}$ 作为 $k$-代数的齐次生成元 $f_{1},\cdots, f_{r}$，其次数 $d_{1}\leq \cdots \leq d_{r}$ 是典范确定的，尽管 $f_{i}$ 本身不典范。当 $G$ 几乎单时，数 $e_{i}=d_{i}-1$ 是 $G$ 的{\em 指数}。把 $f_{i}$ 看作 $G$-不变对称多重线性函数 $\frg^{\otimes d_{i}}\to k$，则对 $X$ 上任一 $G$-主齐性空间 $\calE$，$f_{i}$ 诱导相伴丛之间的映射
\begin{equation*}
f_{i}: \Ad(\calE)^{\otimes d_{i}}\to\calO_{X}.
\end{equation*}
这又诱导
\begin{equation*}
f^{\calL}_{i}:(\Ad(\calE)\otimes\calL)^{\otimes d_{i}}\to\calL^{\otimes d_{i}}.
\end{equation*}
若 $\varphi$ 是 $\calE$ 上 $\calL$-扭 Higgs 场，可把 $f^{\calL}_{i}$ 作用在 $(\Ad(\calE)\otimes\calL)^{\otimes d_{i}}$ 的截面 $\varphi^{\otimes d_{i}}$ 上，得到
\begin{equation*}
a_{i}(\varphi):=f^{\calL}_{i}(\varphi^{\otimes d_{i}})\in\cohog{0}{X,\calL^{\otimes d_{i}}},  \quad i=1,2,\cdots, r.
\end{equation*}
对应 $(\calE,\varphi)\mapsto(a_{i}(\varphi))_{1\leq i\leq r}$ 定义了 $\calM_{G,\calL}$ 的{\em Hitchin 纤维化}
\begin{equation*}
f=f_{G,\calL}: \calM_{G,\calL}\to \calA_{G,\calL}:=\prod_{i=1}^{r}\cohog{0}{X,\calL^{\otimes d_{i}}}.
\end{equation*}
目标 $\calA_{G,\calL}$ 同样看作 $k$ 上仿射空间，称为{\em Hitchin 基}。

定义 Hitchin 基 $\calA_{G,\calL}$ 更内蕴的方式如下。仿射概形 $\frc=\Spec \Sym(\frg^{*})^{G}$ 带有 $\Gm$-作用，它诱导坐标环上的分次。设 $\Tot^{\times}(\calL)\to X$ 是 $\calL$ 全空间中零截面的补集。考虑 $\frc$ 在 $X$ 上的 $\calL$-扭版本：
\begin{equation*}
\frc_{X,\calL}:=(\frc\times \Tot^{\times}(\calL))/\Gm,
\end{equation*}
其中 $\l\in\Gm$ 按 $\l: (c,\wt{x})\mapsto (\l c, \l^{-1}\wt{x})$ 作用在两个坐标上。这是 $X$ 上的仿射空间丛，纤维同构于 $\frc$。于是 $\calA_{G,\calL}$ 可典范地等同于映射 $\frc_{X,\calL}\to X$ 的截面构成的模空间。特别地，每个点 $a\in\calA_{G,\calL}$ 给出映射 $a:X\to [\frc/\Gm]$。

\subsubsection{一般正则半单的轨迹} 把 $\calL$ 在 $X$ 的一般点 $\eta$ 处平凡化并限制 $a_{i}$ 到 $\eta$，得多项式 $P_{a}(y)=y^{n}-a_{1}y^{n-1}+a_{2}y^{n-2}+\cdots +(-1)^{n}a_{n}\in F[y]$，其中 $F=k(X)$ 是 $X$ 的函数域。当 $P_{a}(y)$ 是 $F[y]$ 中的可分多项式时，称这样的 $a$ {\em 一般正则半单}。$a$ 的一般正则半单性等价于判别式 $\Delta(P_{a})\in\cohog{0}{X,\calL^{n(n-1)}}$ 不消没，因此它定义了开子概形 $\Ah_{n,\calL}\subset \calA_{n,\calL}$。

对一般 $G$，把 $a\in\calA_{G,\calL}$ 看作映射 $a:X\to [\frc/\Gm]$（见 \S\ref{sss:Hit fib general} 末）。称 $a$ {\em 一般正则半单}，若 $a$ 把 $X$ 的一般点映入开子叠 $[\frc^{\rs}/\Gm]$。这定义了开子概形 $\Ah_{G,\calL}\subset\calA_{G,\calL}$，推广了上面对 $\Ah_{n,\calL}$ 的构造。

\subsubsection{几何性质}\label{sss:geom Hit} 当 $\deg \calL>2g-2$ 时，叠 $\calM|_{\Ah_{G,\calL}}$ 光滑，见 \cite[定理 4.14.1]{NgoFL}。此时态射 $f_{G,\calL}$ 在 $\Ah_{G,\calL}$ 上平坦，见 \cite[推论 4.16.4]{NgoFL}。当 $G$ 半单时，还有更开稠的子集 $\calA^{\textup{ani}}_{G,\calL}\subset \Ah_{G,\calL}$，在其上 $\calM$ 是 Deligne--Mumford 叠且映射 $f_{G,\calL}$ 真，见 \cite[命题 6.1.3]{NgoFL}。与仿射 Springer 纤维的几何所涉及的无穷维相比，Hitchin 纤维化的几何性质好得多，然而它与仿射 Springer 纤维密切相关，这一点将在 \S\ref{ss:rel Hit ASF} 看到。

\subsubsection{推广}\label{sss:general Hitchin} 设 $H$ 是 $k$ 上约化群，处于约化群的短正合列
\begin{equation*}
1\to H_{1}\to H\to A\to 1.
\end{equation*}
设 $(V,\rho)$ 是 $H$ 的表示。则可考虑数对 $(\calE,\psi)$，其中 $\calE$ 是 $X$ 上 $H$-主齐性空间，$\psi$ 是相伴丛 $\rho(\calE)$ 的截面。等价地，这样的数对等同于态射 $X\to [V/H]$。可以证明存在代数叠 $\calM_{H, \rho}$ 分类这样的数对。每个 $H$-主齐性空间诱导 $A$-主齐性空间，故有态射 $\alpha: \calM_{H, \rho}\to \Bun_{A}$。固定 $X$ 上 $A$-主齐性空间 $\calE_{A}$，把原像 $\alpha^{-1}(\calE_{A})$ 记为 $\calM_{H,\rho,\calE_{A}}$。

为回到 Hitchin 模叠，考虑 $H=G\times\Gm$、$H_{1}=G$、$A=\Gm$ 的情形。取 $V=\frg$，$H$ 的作用 $\rho$ 定义如下：$G$ 按伴随表示作用，$\Gm$ 按伸缩作用在 $V$ 上。$H$-主齐性空间是由 $G$-主齐性空间与 $X$ 上线丛 $\calL$ 构成的数对。固定线丛 $\calL$（这等价于固定 $A=\Gm$-主齐性空间），得到同构
\begin{equation*}
\calM_{H,\rho, \calL}\cong\calM_{G, \calL}.
\end{equation*}

对如上一般的 $(1\to H_{1}\to H\to A\to 1,V,\rho,\calE_{A})$，可如下定义 Hitchin 基的类比。令 $\frc_{H,\rho}=V\sslash H_{1}$。这是 $\frc$ 的类比，它带有 $A$ 的作用。再作 $\frc_{H,\rho}$ 在 $X$ 上的扭版本
\begin{equation*}
\frc_{H,\rho,\calE_{A}}:=\calE_{A}\twtimes{A}\frc_{H,\rho}.
\end{equation*}
然后把 $\calA_{H,\rho,\calE_{A}}$ 定义为映射 $\frc_{H,\rho,\calE_{A}}\to X$ 的截面构成的模空间。态射 $[V/H]\to [\frc_{H,\rho}/A]$ 进而诱导 Hitchin 纤维化的类比
\begin{equation*}
f_{H,\rho,\calE_{A}}: \calM_{H,\rho,\calE_{A}}\to \calA_{H,\rho,\calE_{A}}.
\end{equation*}

\subsubsection{例子}\label{sss:MHilb} 考虑 $H=\GL(U)\times\Gm\times\Gm$ 与 $V=\End(U)\oplus U^{*}$。作用 $\rho(g,s_{1}, s_{2})$ 在 $V$ 上由 $(A,u^{*})\mapsto (s_{1}gAg^{-1}, s_{2}gu^{*})$ 给出。模叠 $\calM_{H,\rho}$ 通过只记住两个 $\Gm$-主齐性空间而映入 $\Pic(X)\times\Pic(X)$。固定 $(\calL_{1},\calL_{2})\in \Pic(X)\times\Pic(X)$，它在 $\calM_{H,\rho}$ 中的原像分类三元组 $(\calU,\varphi,\beta)$，其中 $(\calU,\varphi)$ 是 $X$ 上秩 $n=\dim V$ 的 Higgs 丛，$\beta$ 是 $\calO_{X}$-线性映射 $\calU\to \calL_{2}$。此情形下的 Hitchin 基与经典 Hitchin 基 $\calA_{n,\calL_{1}}$ 相同。\S\ref{sss:Hilb} 将把此模空间与曲线的 Hilbert 概形联系起来。

\subsection{Hitchin 纤维} Hitchin 的一个重要观察是：Hitchin 纤维化 $f$ 的纤维可以用「交换」的语言描述，即用 $X$ 的某些有限覆叠上的线丛来描述。下面对 $G=\GL_{n}$ 与 $G=\Sp_{2n}$ 详述这一观察。

\subsubsection{$\GL_{n}$ 的情形与谱曲线}\label{sss:spectral curve} 对 $a=(a_{1},\cdots, a_{n})\in\calA_{n,\calL}$，可以定义一条曲线 $Y_{a}$，带有 $n$ 次态射 $\pi_{a}: Y_{a}\to X$。构造如下。$\calL$ 的全空间可以写成 $X$ 上的相对谱
\begin{equation*}
\Sigma:=\Tot(\calL)=\Spec(\calO\oplus\calL^{\otimes -1}\oplus\calL^{\otimes -2}\oplus\cdots)=\Spec\Sym(\calO\oplus\calL^{\otimes -1}).
\end{equation*}
设 $\pi: \Sigma\to X$ 是投影。考虑 $X$ 上凝聚层之间的映射
\begin{equation*}
\iota_{a}: \calL^{\otimes -n}\to \pi_{*}\calO_{\Sigma}=\calO\oplus\calL^{\otimes -1}\oplus\calL^{\otimes -2}\oplus\cdots,
\end{equation*}
它在坐标下由 $((-1)^{n}a_{n}, (-1)^{n-1}a_{n-1},\cdots, -a_{1},1,0,\cdots)$ 给出。由伴随性，$\iota_{a}$ 诱导映射 $\iota'_{a}: \pi^{*}\calL^{\otimes -n}\to \calO_{\Sigma}$，其像记为 $\calI_{a}$。则 $\calI_{a}$ 是 $\Sigma$ 上的理想层。定义{\em 谱曲线} $Y_{a}$ 为 $\Sigma$ 中由 $\calI_{a}$ 定义的闭子概形
\begin{equation*}
Y_{a}=\Spec(\calO\oplus\calL^{\otimes-1}\oplus\calL^{\otimes-2}\oplus\cdots)/\calI_{a}.
\end{equation*}
若在某开子集 $U\subset X$ 上平凡化 $\calL$，并把 $a_{i}$ 看作 $U$ 上函数，则 $Y_{a}|_{U}$ 是 $U\times\AA^{1}$ 中由方程 $y^{n}-a_{1}y^{n-1}+a_{2}y^{n-2}+\cdots +(-1)^{n}a_{n}=0$ 定义的子概形（其中 $y$ 是 $\AA^{1}$ 的坐标）。投影 $\pi_{a}: Y_{a}\to X$ 是 $n$ 次有限平坦的。曲线 $Y_{a}$ 称为 $a$ 的{\em 谱曲线}，因为 $\pi_{a}$ 的纤维正是特征多项式 $y^{n}-a_{1}y^{n-1}+a_{2}y^{n-2}+\cdots +(-1)^{n}a_{n}$ 的根。

当 $a\in\Ah_{n,\calL}$ 时，曲线 $Y_{a}$ 约化，因而在某 Zariski 稠密开子集上光滑；此时存在模叠 $\cPic(Y_{a})$，分类 $Y_{a}$ 上无挠且一般秩为 $1$ 的凝聚 $\calO_{Y_{a}}$-模，见 \cite{AK}。分类 $Y_{a}$ 上线丛的通常 Picard 叠 $\Pic(Y_{a})$ 是 $\cPic(Y_{a})$ 的开子叠，它通过张量积作用在 $\cPic(Y_{a})$ 上。

\begin{prop}\label{p:cPic} 设 $a\in\Ah_{n,\calL}(k)$。设 $\calM_{a}$ 是 $f:\calM_{n,\calL}\to \calA_{n,\calL}$ 在 $a$ 上的纤维。则有叠的典范同构
\begin{equation*}
\cPic(Y_{a})\cong\calM_{a}.
\end{equation*}
\end{prop}

\subsubsection{证明梗概}\label{sss:pf cPic}
给出态射 $\cPic(Y_{a})\to \calM_{a}$。对 $Y_{a}$ 上任意凝聚层 $\calF$，正向像 $\pi_{a,*}\calF$ 是 $X$ 上凝聚层，带有映射 $\varphi_{\calF}: \pi_{a,*}\calF\otimes\calL^{-1}\to\pi_{a,*}\calF$，因为 $\calF$ 是 $\calO_{Y_{a}}$-模，而 $\calO_{Y_{a}}$ 把 $\calL^{-1}$ 作为第二个直和项包含在内。当 $\calF$ 无挠且一般秩为 $1$ 时，$\calV:=\pi_{a,*}\calF$ 在 $X$ 上无挠（因而是向量丛），秩为 $n$，且映射 $\varphi_{\calF}$ 诱导 Higgs 场 $\varphi:\calV\to\calV\otimes\calL$。对应 $\calF\mapsto (\calV,\varphi)$ 定义了态射 $\cPic(Y_{a})\to \calM_{a}$，可以证明它是同构。\qed

\subsubsection{$G=\Sp_{2n}$ 的情形} 此时 $a\in\calA_{G,\calL}$ 是多元组 $(a_{1},\cdots, a_{n})$，其中 $a_{i}\in\cohog{0}{X,\calL^{\otimes2i}}$。对 $a\in\calA_{G,\calL}$，同样可以定义谱曲线 $Y_{a}$，它是 $\calL$ 全空间中由局部生成元
\begin{equation*}
P_{a}(y)=y^{2n}+a_{1}y^{2n-2}+\cdots+a_{n}
\end{equation*}
所生成理想截出的闭子概形。注意 $Y_{a}$ 带有对合 $\sigma(y)=-y$，投影 $\pi_{a}: Y_{a}\to X$ 在此对合下不变。现设 $P_{a}(y)$ 限制到 $X$ 的一般点时可分，从而 $Y_{a}$ 约化。对合 $\sigma$ 诱导紧化 Picard 叠 $\cPic(Y_{a})$ 上的对合 $\sigma$。设 $(-)^{\vee}$ 是 $Y_{a}$ 上凝聚层的相对 Serre 对偶函子，即 $\calF^{\vee}=\un{\Hom}_{\calO_{X}}(\calF,\calO_{X})$，按自然方式看作 $\calO_{Y_{a}}$-模。设 $\cPrym(Y_{a};\sigma)$ 是数对 $(\calF,\iota)$ 的模叠，其中 $\calF\in\cPic(Y_{a})$，$\iota$ 是同构 $\iota: \sigma^{*}\calF\cong\calF^{\vee}$，且满足 $(\sigma^{*}\iota)^{\vee}=-\iota$。称 $\cPrym(Y_{a};\sigma)$ 为 $Y_{a}$ 关于对合 $\sigma$ 的{\em 紧化 Prym 叠}。与 $\GL_{n}$ 情形类似，$\calM_{a}$ 有如下描述。

\begin{prop} 设 $a\in\calA_{G,\calL}$ 使得 $Y_{a}$ 约化。则有叠的典范同构
\begin{equation*}
\cPrym(Y_{a})\cong\calM_{a}.
\end{equation*}
\end{prop}

\subsubsection{证明梗概}
对 $a\in\calA_{G,\calL}$，$\calM_{a}$ 中的点是三元组 $(\calV,\omega, \varphi)$，其中 $\calV$ 是 $X$ 上秩 $2n$ 的向量丛，$\omega: \wedge^{2}\calV\to\calO_{X}$ 是 \S\ref{sss:Bun Sp} 中 $\calV$ 上的辛形式，$\varphi:\calV\to\calV\otimes\calL$ 满足对 $\calV$ 的局部截面 $u,v$ 有 $\om(\varphi u, v)+\om(u, \varphi v)=0$，且 $\varphi$ 的特征多项式为 $P_{a}(y)$。由命题 \ref{p:cPic}，Higgs 丛 $(\calV,\varphi)$ 给出点 $\calF\in\cPic(Y_{a})$。辛形式可以看作同构 $\jmath: (\calV,\varphi)\isom (\calV^{\vee}, -\varphi^{\vee})$，满足 $\jmath^{\vee}=-\jmath$。注意 Higgs 丛 $(\calV^{\vee}, -\varphi^{\vee})$ 在命题 \ref{p:cPic} 的同构下对应 $\sigma^{*}\calF^{\vee}$。于是同构 $\jmath$ 变成 $\iota:  \sigma^{*}\calF\isom\calF^{\vee}$，满足 $(\sigma^{*}\iota)^{\vee}=-\iota$。\qed

\subsubsection{例 \ref{sss:MHilb} 续}\label{sss:Hilb}
固定 $X$ 上两个线丛 $\calL_{1}$ 与 $\calL_{2}$，设 $\calH_{\calL_{1},\calL_{2}}$ 是 \S\ref{sss:MHilb} 中模叠 $\calM_{H,\rho}$ 在 $(\calL_{1},\calL_{2})$ 上的纤维。则 $\calH_{\calL_{1},\calL_{2}}$ 分类 $(\calV,\varphi,\beta)$，其中 $\calV$ 是 $X$ 上秩 $n$ 的向量丛，$\varphi:\calU\to\calU\otimes\calL_{1}$ 是 Higgs 场，$\beta: \calU\to \calL_{2}$。有 Hitchin 型映射 $\calH_{\calL_{1},\calL_{2}}\to\calA_{n,\calL_{1}}$，把 $(\calV,\varphi,\beta)$ 送到 $(a_{i}(\varphi))_{1\leq i\leq n}$。设 $a\in\Ah_{n,\calL_{1}}$，$\calH_{\calL_{1},\calL_{2},a}$ 是 $\calH_{\calL_{1},\calL_{2}}$ 在 $a$ 上的纤维。考虑与 \S\ref{sss:spectral curve} 中相同的谱曲线 $Y_{a}$。用命题 \ref{p:cPic} 可把 $(\calU,\varphi)$ 等同于点 $\calF\in\cPic(Y_{a})$。映射 $\beta:\pi_{a,*}\calF=\calU\to\calL_{2}$ 由伴随性给出映射 $b:\calF\to \pi^{!}_{a}\calL_{2}$。这里 $\pi^{!}_{a}\calL_{2}\cong\pi^{*}_{a}\calL_{2}\otimes\omega_{Y_{a}/X}$，而相对对偶复形 $\omega_{Y_{a}/X}=\omega_{Y_{a}}\otimes\pi_{a}^{*}\omega_{X}^{-1}$ 是 $Y_{a}$ 上线丛，因为 $Y_{a}$ 是平面曲线，从而 Gorenstein。由于 $\calF$ 无挠且一般是线丛，$b$ 是凝聚层的单射。因此数据 $(\calU,\varphi,\beta)$ 变成线丛 $\pi^{!}_{a}\calL_{2}$ 的凝聚子层 $\calF$ 的数据。由于 $\pi^{!}_{a}\calL_{2}$ 是线丛，子层 $\calF$ 由商 $(\pi^{!}_{a}\calL_{2})/\calF$ 的支撑决定，而后者是 $Y_{a}$ 的零维子概形。于是得到典范同构
\begin{equation*}
\calH_{\calL_{1},\calL_{2},a}\cong\Hilb(Y_{a}),
\end{equation*}
其中 $\Hilb(Y_{a})$ 是 $Y_{a}$ 的各种长度零维子概形的 Hilbert 概形的不交并。

总之，一族平面曲线的点 Hilbert 概形可以实现在 \S\ref{sss:general Hitchin} 的框架中，作为 Hitchin 型模空间。

\subsubsection{Hitchin 纤维上的对称性}\label{sss:sym Hit}
在 $G=\GL_{n}$、$a\in\Ah_{n,\calL}$ 的情形，已把 Hitchin 纤维 $\calM_{a}$ 与谱曲线 $Y_{a}$ 的紧化 Picard 叠 $\cPic(Y_{a})$ 等同。通常 Picard 叠 $\Pic(Y_{a})$ 通过张量积作用在 $\cPic(Y_{a})$ 上。因此有 $\Pic(Y_{a})$ 在 $\calM_{a}$ 上的作用。当 $Y_{a}$ 光滑时，此作用单可迁。

在 $G=\Sp_{2n}$、$a\in\Ah_{G,\calL}$ 的情形，Prym 叠 $\Prym(Y_{a};\sigma)$ 定义为数对 $(\calL,\iota)$ 的模叠，其中 $\calL\in\Pic(Y_{a})$，$\iota:\sigma^*\calL\cong\calL^{-1}$ 满足 $(\sigma^{*}\iota)^{\vee}=\iota$。则 $\Prym(Y_{a};\sigma)$ 通过张量积作用在 $\calM_{a}\cong\cPrym(Y_{a};\sigma)$ 上。

对一般的 $G$ 与 $a\in\Ah_{G,\calL}$，同样可以定义交换群叠 $\calP_{a}$，它作用在相应的 Hitchin 纤维 $\calM_{a}$ 上。事实上，经由映射 $a:X\to [\frc/\Gm]$，\S\ref{sss:sym ASF} 中下降到 $[\frc/\Gm]$ 的正则中心化子群概形 $J$ 拉回为 $X$ 上光滑群概形 $J_{a}$，它一般地是环面。叠 $\calP_{a}$ 于是分类 $X$ 上 $J_{a}$-主齐性空间。

\subsection{与仿射 Springer 纤维的关系}\label{ss:rel Hit ASF}
本小节陈述 Ng\^o 所观察到的 Hitchin 纤维与仿射 Springer 纤维之间的精确关系。

\subsubsection{} 固定点 $a\in\Ah_{G,\calL}(k)$。把 $a$ 看作映射 $X\to [\frc/\Gm]$，设 $U_{a}\subset X$ 是 $[\frc^{\rs}/\Gm]$ 在 $a$ 下的原像。由于 $a\in\Ah_{G,\calL}$，$U_{a}$ 非空，故补集 $X-U_{a}$ 由 $X$ 的有限多个 $k$-点组成。对每个点 $x\in X-U_{a}$，设 $\calO_{x}$ 是 $X$ 在 $x$ 处的完备局部环，其分式域为 $F_{x}$，剩余域 $k(x)=k$。固定 $\calL$ 在 $x$ 附近的平凡化，可把 $a_{x}=a|_{\Spec\calO_{x}}$ 等同于 $\frc(\calO_{x})\cap\frc^{\rs}(F_{x})$ 中的元素。设 $\ep(a_{x})\in\frg^{\rs}(F_{x})$ 是它在 Kostant 截面中的对应点，则得到仿射 Grassmann 簇 $\Gr_{G, x}=L_{x}G/L^{+}_{x}G$ 中的仿射 Springer 纤维 $\sX_{a_{x}}:=\sX_{\ep(a_{x})}$（把 $x$ 放在下标中，是为了强调环路群的定义用的是域 $F_{x}$，它同构于 $k\lr{t}$）。中心化子的环路群 $LG_{\ep(a_{x})}$ 作用在 $\sX_{a_{x}}$ 上，且此作用如 \S\ref{sss:sym ASF} 那样经过局部 Picard 群 $P_{a_{x}}$ 分解。另一方面，\S\ref{sss:sym Hit} 提到整体 Picard 叠 $\calP_{a}$ 作用在 Hitchin 纤维 $\calM_{a}$ 上。Ng\^o 的乘积公式大意是说：模去局部与整体 Picard 叠的作用后，$\calM_{a}$ 是所有点 $x\in X-U_{a}$ 处仿射 Springer 纤维 $\sX_{a_{x}}$ 的乘积。

\begin{theorem}[乘积公式，Ng\^o {\cite[定理 4.6]{NgoHit} 与 \cite[命题 4.15.1]{NgoFL}}]\label{th:prod} 对 $a\in\Ah(k)$，存在典范态射
\begin{equation*}
\calP_{a}\twtimes{\prod_{x\in X-U_{a}}P_{a_{x}}}\left(\prod_{x\in X-U_{a}}\sX_{a_{x}}\right)\to \calM_{a},
\end{equation*}
它是叠的同胚。这里记号 $P\twtimes{H}Y$（其中 $H$ 右作用在 $P$ 上、左作用在 $Y$ 上）表示 $P\times Y$ 对作用 $h\cdot (p,y)=(ph^{-1}, hy)$ 的商，其中 $p\in P, y\in Y, h\in H$\zy{原文此处用 homeomorphism（同胚）一词；对象是模叠，读者可按叠意义下的同构理解，译文照原文作「同胚」。同一记号 $\twtimes{}$ 在 \S4.2 已用于 $\frc_{H,\rho,\calE_{A}}=\calE_{A}\twtimes{A}\frc_{H,\rho}$。来源：原文 \S4.4 该定理文脉与 \S4.2。}。
\end{theorem}

当 $G=\GL_{n}$ 时，乘积公式可以用谱曲线的紧化 Picard 叠以更熟悉的方式重述；后者实际上对一切约化曲线都有意义。设 $C$ 是 $k$ 上约化射影曲线。对每个点 $x\in C(k)$，可以定义局部 Picard 群 $P_{x}(C)$，其 $k$-点是 $F^{\times}_{x}/\calO^{\times}_{x}$，其中 $\calO_{x}$ 是 $C$ 在 $x$ 处的完备局部环，$F_{x}$ 是其全分式环（一般是域之积）。还有 $\cPic(C)$ 的局部类比 $\ov{P}_{x}(C)$，其 $k$-点是 $F_{x}$ 的分式 $\calO_{x}$-理想，参看 \S\ref{sss:GLn orb formula}。于是有乘积公式的如下变体，其证明与定理 \ref{th:prod} 类似。

\begin{prop}\label{p:Pic prod} 设 $C$ 是 $k$ 上约化射影曲线，$Z=C-C^{sm}$ 是 $C$ 的奇点轨迹。则有典范态射
\begin{equation*}
\Pic(C)\twtimes{\prod_{x\in Z}P_{x}(C)}\left(\prod_{x\in Z}\ov{P}_{x}(C)\right)\to\cPic(C),
\end{equation*}
它是叠的同胚。
\end{prop}

乘积公式提供了 Hitchin 纤维的几何与仿射 Springer 纤维的几何之间的联系，而后者如前所见与轨道积分密切相关。这一联系使得可以通过研究 Hitchin 纤维的上同调来接近基本引理。用 Hitchin 纤维代替仿射 Springer 纤维的好处是：Hitchin 纤维化具有更好的几何性质，这一点已在 \S\ref{sss:geom Hit} 看到。

\subsection{Springer 作用的整体版本}

定理 \ref{th:prod} 中的乘积公式提示：应当存在仿射 Springer 理论的整体类比，其中仿射 Springer 纤维被 Hitchin 纤维取代。这样的理论在作者本人自 \cite{GS} 开始的一系列论文中得到发展。

\subsubsection{Iwahori 水平结构} 现定义仿射 Springer 纤维 $\sY_{\gamma}$ 的 Hitchin 型类比。仍固定曲线 $X$ 及其上线丛 $\calL$。设 $\Mpar_{G,\calL}$ 是分类 $(x,\calE,\varphi, \calE^{B}_{x})$ 的模叠，其中 $x\in X$，$(\calE,\varphi)$ 是 $\calL$-扭 Higgs $G$-丛，$\calE^{B}_{x}$ 是把 $x$ 处的纤维 $\calE_{x}$ 约化为 $B$-主齐性空间 $\calE^{B}_{x}$ 的约化（这里 $B\subset G$ 是 $G$ 的 Borel 子群），且与 $\varphi$ 相容。

当 $G=\GL_{n}$ 时，$\Mpar_{G,\calL}$ 除 $\calL$-扭 Higgs 丛 $(\calV,\varphi)$ 之外，还分类点 $x\in X$ 与纤维 $\calV_{x}$ 的完整旗。这样的完整旗等同于凝聚子层链
\begin{equation*}
\calV(-x)=\calV_{0}\subset \calV_{1}\subset\cdots\subset\calV_{n-1}\subset\calV_{n}=\calV,
\end{equation*}
使得 $\calV_{i}/\calV_{i-1}$ 是支撑在 $x$ 处、长度为 $1$ 的层。$\varphi$ 与完整旗的相容性条件是：对每个 $0\leq i\leq n$，$\varphi$ 限制为映射 $\calV_{i}\to\calV_{i}\otimes\calL$。

有 Hitchin 纤维化的类比
\begin{equation*}
\fpar_{G,\calL}:\Mpar_{G,\calL}\to \calA_{G,\calL}\times X,
\end{equation*}
它除 $a=f_{G,\calL}(\calE,\varphi)\in\calA_{G,\calL}$ 之外，还记下数据中的点 $x\in X$。

\begin{theorem}[见 {\cite[定理 3.3.3]{GS}}]\label{th:GS} 设 $\deg\calL>2g-1$。则复形 $\bR\fpar_{!}\Ql$ 在 $\Ah_{G,\calL}\times X$ 的某开子集 $(\Ah_{G,\calL}\times X)'$ 上的限制，带有 $\tilW=\xcoch(T)\rtimes W$ 的自然作用。\footnote{当 $\textup{char}(k)=0$ 时可取 $(\Ah_{G,\calL}\times X)'=\Ah_{G,\calL}\times X$；一般情形下，限制到开子集 $(\Ah_{G,\calL}\times X)'$ 并不妨碍应用于仿射 Springer 纤维的问题。}
\end{theorem}

\subsubsection{扩张对称性} 正如仿射 Springer 纤维的情形，定理 \ref{th:GS} 中的 $\tilW$-作用可扩张为双仿射 Hecke 代数的分级版本的作用，见 \cite[定理 6.1.6]{GS}。此外还有乘积公式，把 $(a,x)\in\Ah_{G,\calL}\times X$ 上的纤维 $\Mpar_{a,x}$ 与仿射 Springer 纤维 $\sY_{a_{x}}$ 及 $y\neq x$ 处仿射 Springer 纤维 $\sX_{a_{y}}$ 之积联系起来。由此诱导的 $\tilW$-作用在茎 $(\bR\fpar_{!}\Ql)_{a,x}\cong\cohog{*}{\Mpar_{a,x}}$ 上，与定理 \ref{th:aff action} 中 $\cohog{*}{\sY_{a_{x}}}$ 上的仿射 Springer 作用相容。Hitchin 纤维上同调与仿射 Springer 纤维上同调之间的这一 $\tilW$-作用联系，可用于证明关于仿射 Springer 作用的结论。这样的应用见 \cite{YSph}。

\subsection{习题}

\subsubsection{} 设 $G=\SL_{n}$。用谱曲线描述 $\Ah_{G,\calL}$ 上的 Hitchin 纤维。

\subsubsection{} 设 $G=\SL_{n}$。计算 $\calA_{G,\calL}$ 与 Hitchin 纤维 $\calM_{a}$ 的维数。当 $\calL=\omega_{X}$ 时，验证两者维数相同。这是此情形下 Hitchin 纤维化为 Lagrange 纤维化的数值证据。

\subsubsection{} 描述 $G=\SO_{n}$ 的 Hitchin 基。

\subsubsection{} 对 $G=\SO_{n}$，用谱曲线描述 Hitchin 纤维。

\subsubsection{} 设 $C$ 是 $k$ 上有理曲线，只有唯一奇点 $x_{0}$，且它是单分支的（即 $x_{0}$ 在正规化 $\tilC\cong\PP^{1}$ 中的原像是一个点）。设 $\overline{P}_{x_{0}}$ 是完备局部环 $\wh\calO_{C,x_{0}}$ 的分式理想的模空间。证明存在典范同胚
\begin{equation*}
\overline{P}_{x_{0}}\to \cPic(C).
\end{equation*}
并解释为何这是命题 \ref{p:Pic prod} 的特例。


\clearpage
\phantomsection
\addcontentsline{toc}{section}{参考文献}
\begin{thebibliography}{99}

\bibitem{AK} Altman, Allen B., Kleiman, Steven L., \emph{Compactifying the Picard scheme}, Adv. in Math. \textbf{35} (1980), no. 1, 50--112.

\bibitem{BL} Beauville, Arnaud, Laszlo, Yves, \emph{Conformal blocks and generalized theta functions}, Comm. Math. Phys. \textbf{164} (1994), no. 2, 385--419.

\bibitem{BBD} Be{\u\i}linson, A. A., Bernstein, J., Deligne, P., \emph{Faisceaux pervers}, Analysis and topology on singular spaces, I, Ast\'erisque \textbf{100}, Soc. Math. France, Paris, 1982, 5--171.

\bibitem{Bez} Bezrukavnikov, Roman, \emph{The dimension of the fixed point set on affine flag manifolds}, Math. Res. Lett. \textbf{3} (1996), no. 2, 185--189.

\bibitem{BM} Borho, Walter, MacPherson, Robert, \emph{Partial resolutions of nilpotent varieties}, Analysis and topology on singular spaces, II, III, Ast\'erisque \textbf{101}, Soc. Math. France, Paris, 1983, 23--74.

\bibitem{Bourbaki} Bourbaki, N., \emph{\'El\'ements de math\'ematique. Fasc. XXXIV. Groupes et alg\`ebres de Lie. Chapitre IV: Groupes de Coxeter et syst\`emes de Tits. Chapitre V: Groupes engendr\'es par des r\'eflexions. Chapitre VI: syst\`emes de racines}, Actualit\'es Scientifiques et Industrielles, No. 1337, Hermann, Paris, 1968, 288 pp. (loose errata).

\bibitem{Bouthier} Bouthier, Alexis, \emph{Dimension des fibres de Springer affines pour les groupes}, Transform. Groups \textbf{20} (2015), no. 3, 615--663.

\bibitem{Brieskorn} Brieskorn, E., \emph{Singular elements of semi-simple algebraic groups}, Actes du Congr\`es International des Math\'ematiciens, Gauthier-Villars, Paris, 1971, 279--284.

\bibitem{BT1} Bruhat, F., Tits, J., \emph{Groupes r\'eductifs sur un corps local}, Inst. Hautes \'Etudes Sci. Publ. Math. (1972), no. 41, 5--251.

\bibitem{CG} Chriss, Neil, Ginzburg, Victor, \emph{Representation theory and complex geometry}, Modern Birkh\"auser Classics, Birkh\"auser Boston, Inc., Boston, MA, 2010, x+495.

\bibitem{dC} de Cataldo, Mark, \emph{Perverse sheaves and the topology of algebraic varieties}, These Proceedings.

\bibitem{DLP} De Concini, C., Lusztig, G., Procesi, C., \emph{Homology of the zero-set of a nilpotent vector field on a flag manifold}, J. Amer. Math. Soc. \textbf{1} (1988), no. 1, 15--34.

\bibitem{Faltings} Faltings, Gerd, \emph{Algebraic loop groups and moduli spaces of bundles}, J. Eur. Math. Soc. (JEMS) \textbf{5} (2003), no. 1, 41--68.

\bibitem{GKM} Goresky, Mark, Kottwitz, Robert, Macpherson, Robert, \emph{Homology of affine Springer fibers in the unramified case}, Duke Math. J. \textbf{121} (2004), no. 3, 509--561.

\bibitem{GKMPurity} Goresky, Mark, Kottwitz, Robert, MacPherson, Robert, \emph{Purity of equivalued affine Springer fibers}, Represent. Theory \textbf{10} (2006), 130--146 (electronic).

\bibitem{Gross} Gross, B. H., \emph{Parahorics}.

\bibitem{Hi} Hitchin, Nigel, \emph{Stable bundles and integrable systems}, Duke Math. J. \textbf{54} (1987), no. 1, 91--114.

\bibitem{HottaSpringer} Hotta, R., Springer, T. A., \emph{A specialization theorem for certain Weyl group representations and an application to the Green polynomials of unitary groups}, Invent. Math. \textbf{41} (1977), no. 2, 113--127.

\bibitem{Kac} Kac, Victor G., \emph{Infinite-dimensional Lie algebras}, Cambridge University Press, Cambridge, 1990, xxii+400.

\bibitem{KL80} Kazhdan, David, Lusztig, George, \emph{A topological approach to Springer's representations}, Adv. in Math. \textbf{38} (1980), no. 2, 222--228.

\bibitem{KL} Kazhdan, D., Lusztig, G., \emph{Fixed point varieties on affine flag manifolds}, Israel J. Math. \textbf{62} (1988), no. 2, 129--168.

\bibitem{KV} Kottwitz, Robert, Viehmann, Eva, \emph{Generalized affine Springer fibres}, J. Inst. Math. Jussieu \textbf{11} (2012), no. 3, 569--609.

\bibitem{Laumon} Laumon, G{\'e}rard, \emph{Transformation de Fourier homog\`ene}, Bull. Soc. Math. France \textbf{131} (2003), no. 4, 527--551.

\bibitem{LuGreen} Lusztig, G., \emph{Green polynomials and singularities of unipotent classes}, Adv. in Math. \textbf{42} (1981), no. 2, 169--178.

\bibitem{LuIC} Lusztig, G., \emph{Intersection cohomology complexes on a reductive group}, Invent. Math. \textbf{75} (1984), no. 2, 205--272.

\bibitem{LuChShI} Lusztig, George, \emph{Character sheaves. I}, Adv. in Math. \textbf{56} (1985), no. 3, 193--237.

\bibitem{LuChShII} Lusztig, George, \emph{Character sheaves. II, III}, Adv. in Math. \textbf{57} (1985), no. 3, 226--265, 266--315.

\bibitem{LuChShIV} Lusztig, George, \emph{Character sheaves. IV}, Adv. in Math. \textbf{59} (1986), no. 1, 1--63.

\bibitem{LuChShV} Lusztig, George, \emph{Character sheaves. V}, Adv. in Math. \textbf{61} (1986), no. 2, 103--155.

\bibitem{KLAffHk} Kazhdan, David, Lusztig, George, \emph{Proof of the Deligne-Langlands conjecture for Hecke algebras}, Invent. Math. \textbf{87} (1987), no. 1, 153--215.

\bibitem{LMB} Laumon, G{\'e}rard, Moret-Bailly, Laurent, \emph{Champs alg\'ebriques}, Ergebnisse der Mathematik und ihrer Grenzgebiete. 3. Folge. A Series of Modern Surveys in Mathematics [Results in Mathematics and Related Areas. 3rd Series. A Series of Modern Surveys in Mathematics] \textbf{39}, Springer-Verlag, Berlin, 2000, xii+208.

\bibitem{LuGrHk} Lusztig, George, \emph{Cuspidal local systems and graded Hecke algebras. I}, Inst. Hautes \'Etudes Sci. Publ. Math. (1988), no. 67, 145--202.

\bibitem{LS} Lusztig, G., Smelt, J. M., \emph{Fixed point varieties on the space of lattices}, Bull. London Math. Soc. \textbf{23} (1991), no. 3, 213--218.

\bibitem{L96} Lusztig, George, \emph{Affine Weyl groups and conjugacy classes in Weyl groups}, Transform. Groups \textbf{1} (1996), no. 1-2, 83--97.

\bibitem{NgoHit} Ng{\^o}, Bao Ch{\^a}u, \emph{Fibration de Hitchin et endoscopie}, Invent. Math. \textbf{164} (2006), no. 2, 399--453.

\bibitem{NgoFL} Ng{\^o}, Bao Ch{\^a}u, \emph{Le lemme fondamental pour les alg\`ebres de Lie}, Publ. Math. Inst. Hautes \'Etudes Sci. (2010), no. 111, 1--169.

\bibitem{OY} Oblomkov, Alexei, Yun, Zhiwei, \emph{Geometric representations of graded and rational Cherednik algebras}.

\bibitem{Rossmann} Rossmann, W., \emph{Invariant eigendistributions on a semisimple Lie algebra and homology classes on the conormal variety. II. Representations of Weyl groups}, J. Funct. Anal. \textbf{96} (1991), no. 1, 155--193.

\bibitem{Sage} Sage, D. S., \emph{A construction of representations of affine Weyl groups}, Compositio Math. \textbf{108} (1997), no. 3, 241--245.

\bibitem{Serre} Serre, Jean-Pierre, \emph{Corps locaux}, Hermann, Paris, 1968, 245.

\bibitem{Slo4} Slodowy, Peter, \emph{Four lectures on simple groups and singularities}, Communications of the Mathematical Institute, Rijksuniversiteit Utrecht \textbf{11}, Rijksuniversiteit Utrecht, Mathematical Institute, Utrecht, 1980, ii+64.

\bibitem{Slodowy} Slodowy, Peter, \emph{Simple singularities and simple algebraic groups}, Lecture Notes in Mathematics \textbf{815}, Springer, Berlin, 1980, x+175.

\bibitem{Spa76} Spaltenstein, N., \emph{The fixed point set of a unipotent transformation on the flag manifold}, Nederl. Akad. Wetensch. Proc. Ser. A {\bf 79}=Indag. Math. \textbf{38} (1976), no. 5, 452--456.

\bibitem{Spa77} Spaltenstein, Nicolas, \emph{On the fixed point set of a unipotent element on the variety of Borel subgroups}, Topology \textbf{16} (1977), no. 2, 203--204.

\bibitem{Spaltenstein} Spaltenstein, Nicolas, \emph{Classes unipotentes et sous-groupes de Borel}, Lecture Notes in Mathematics \textbf{946}, Springer-Verlag, Berlin-New York, 1982, ix+259.

\bibitem{Springer} Springer, T. A., \emph{Trigonometric sums, Green functions of finite groups and representations of Weyl groups}, Invent. Math. \textbf{36} (1976), 173--207.

\bibitem{SpringerPure} Springer, T. A., \emph{A purity result for fixed point varieties in flag manifolds}, J. Fac. Sci. Univ. Tokyo Sect. IA Math. \textbf{31} (1984), no. 2, 271--282.

\bibitem{St} Steinberg, Robert, \emph{On the desingularization of the unipotent variety}, Invent. Math. \textbf{36} (1976), 209--224.

\bibitem{StConj} Steinberg, Robert, \emph{Conjugacy classes in algebraic groups}, Lecture Notes in Mathematics, Vol. 366, Springer-Verlag, Berlin-New York, 1974, vi+159.

\bibitem{Tits} Tits, J., \emph{Reductive groups over local fields}, Automorphic forms, representations and $L$-functions, Proc. Sympos. Pure Math., XXXIII, Amer. Math. Soc., Providence, R.I., 1979, 29--69.

\bibitem{Vass} Vasserot, Eric, \emph{Induced and simple modules of double affine Hecke algebras}, Duke Math. J. \textbf{126} (2005), no. 2, 251--323.

\bibitem{VV} Varagnolo, M., Vasserot, E., \emph{Finite-dimensional representations of DAHA and affine Springer fibers: the spherical case}, Duke Math. J. \textbf{147} (2009), no. 3, 439--540.

\bibitem{GS} Yun, Zhiwei, \emph{Global Springer theory}, Adv. Math. \textbf{228} (2011), no. 1, 266--328.

\bibitem{Y-Orb} Yun, Zhiwei, \emph{Orbital integrals and Dedekind zeta functions}, The legacy of Srinivasa Ramanujan, Ramanujan Math. Soc. Lect. Notes Ser. \textbf{20}, Ramanujan Math. Soc., Mysore, 2013, 399--420.

\bibitem{YSph} Yun, Zhiwei, \emph{The spherical part of the local and global Springer actions}, Math. Ann. \textbf{359} (2014), no. 3-4, 557--594.

\bibitem{Zhu} Zhu, Xinwen, \emph{Introduction to the affine Grassmannian and the geometric Satake isomorphism}, These Proceedings.

\end{thebibliography}


\end{document}
